\UseRawInputEncoding
\documentclass[a4paper,12pt]{article}
\usepackage[utf8]{inputenc}
\usepackage[T2A]{fontenc}
\usepackage[english,russian]{babel}
\usepackage{amsmath,amsfonts,amssymb}
\usepackage{hyperref}
\usepackage{listings}
\usepackage{xcolor}
\usepackage{geometry}
\geometry{top=2cm,bottom=2cm,left=2.5cm,right=2cm}

% Настройка листингов для поддержки кириллицы
\lstset{
    language=C++,
    basicstyle=\ttfamily\small,
    keywordstyle=\color{blue},
    commentstyle=\color{green!60!black},
    stringstyle=\color{red},
    numbers=left,
    numberstyle=\tiny,
    breaklines=true,
    frame=single,
    tabsize=2,
    extendedchars=true,
    literate= % правила замены русских букв
        {а}{{\"a}}1
        {б}{{\"b}}1
        {в}{{\"v}}1
        {г}{{\"g}}1
        {д}{{\"d}}1
        {е}{{\"e}}1
        {ё}{{\"e}}1
        {ж}{{\"z}}1
        {з}{{\"z}}1
        {и}{{\"i}}1
        {й}{{\"i}}1
        {к}{{\"k}}1
        {л}{{\"l}}1
        {м}{{\"m}}1
        {н}{{\"n}}1
        {о}{{\"o}}1
        {п}{{\"p}}1
        {р}{{\"r}}1
        {с}{{\"s}}1
        {т}{{\"t}}1
        {у}{{\"u}}1
        {ф}{{\"f}}1
        {х}{{\"h}}1
        {ц}{{\"c}}1
        {ч}{{\"c}}1
        {ш}{{\"s}}1
        {щ}{{\"s}}1
        {ъ}{{\"\'}}1
        {ы}{{\"y}}1
        {ь}{{\"\'}}1
        {э}{{\"e}}1
        {ю}{{\"u}}1
        {я}{{\"a}}1
        {А}{{\"A}}1
        {Б}{{\"B}}1
        {В}{{\"V}}1
        {Г}{{\"G}}1
        {Д}{{\"D}}1
        {Е}{{\"E}}1
        {Ё}{{\"E}}1
        {Ж}{{\"Z}}1
        {З}{{\"Z}}1
        {И}{{\"I}}1
        {Й}{{\"I}}1
        {К}{{\"K}}1
        {Л}{{\"L}}1
        {М}{{\"M}}1
        {Н}{{\"N}}1
        {О}{{\"O}}1
        {П}{{\"P}}1
        {Р}{{\"R}}1
        {С}{{\"S}}1
        {Т}{{\"T}}1
        {У}{{\"U}}1
        {Ф}{{\"F}}1
        {Х}{{\"H}}1
        {Ц}{{\"C}}1
        {Ч}{{\"C}}1
        {Ш}{{\"S}}1
        {Щ}{{\"S}}1
        {Ъ}{{\"\'}}1
        {Ы}{{\"Y}}1
        {Ь}{{\"\'}}1
        {Э}{{\"E}}1
        {Ю}{{\"U}}1
        {Я}{{\"A}}1
        {_}{{\_}}1           
        {<}{{\textless}}1    
        {>}{{\textgreater}}1}

\title{Фреймворк \texttt{robosd} \\ (роботы -- космическая мечта)}
\author{Юсупов Андрей Николаевич}
\date{\today}

\begin{document}

\maketitle
\tableofcontents
\newpage

\section{Введение}

Фреймворк \texttt{robosd} предназначен для разработки программного обеспечения робототехнических систем. 
Основные цели проекта:
\begin{itemize}
    \item максимальная переносимость между различными аппаратными платформами и операционными системами;
    \item чёткое разделение на \textbf{backend} (жёсткое реальное время, минимальные задержки) и \textbf{frontend} (управление, логирование, взаимодействие с пользователем);
    \item модульная архитектура, позволяющая легко добавлять новые устройства и алгоритмы;
    \item встроенные средства диагностики, логирования и работы с конфигурационными файлами.
\end{itemize}

Фреймворк активно использует препроцессорную настройку через макросы, определённые в файле \texttt{robosd\_app\_tuning.hpp} (не показан в данной документации, но подразумевается). 
Все компоненты расположены в каталогах \texttt{core}, \texttt{burst}, \texttt{net}, \texttt{servo} и других, однако в этом документе основное внимание уделено ядру (\texttt{core}).

\section{Общая архитектура}

Ключевая идея -- разделение на два контекста выполнения:
\begin{description}
    \item[Backend] -- поток (или прерывание) с наивысшим приоритетом, выполняющийся строго периодически. В нём запрещены блокировки, вызовы операционной системы, динамическое выделение памяти (кроме специально подготовленного пула). Здесь выполняются регуляторы, обработка датчиков, ШИМ.
    \item[Frontend] -- обычные потоки ОС, в которых работают интерфейсы, настройка, ведение логов, обработка команд. Взаимодействие с backend'ом осуществляется через безаалочные очереди и разделяемые данные, защищённые специальными примитивами синхронизации.
\end{description}

Для синхронизации между контекстами предусмотрены классы \texttt{system::guard} и \texttt{system::critical}, которые корректно обрабатывают ситуации вызова из разных контекстов (например, из backend'а нельзя входить в критическую секцию, блокирующую frontend).

\section{Ядро (\texttt{core})}

Каталог \texttt{core} содержит базовые механизмы, общие для всех остальных подсистем: определение типов, макросы, работу со строками, системное окружение, управление памятью, логирование, конфигурацию, списки и т.д.

\subsection{Общие определения (\texttt{robosd\_common.hpp})}

Файл подключает \texttt{robosd\_app\_tuning.hpp} и \texttt{robosd\_target.hpp}, после чего задаёт основные типы и вспомогательные макросы.

\subsubsection{Типы и макросы}

\begin{lstlisting}[caption={Определения типов}]
typedef ROBO_TYPE_RANDOM random_t;
typedef ROBO_TYPE_TIME_US time_us_t;
typedef ROBO_TYPE_TIME_MS time_ms_t;
typedef ROBO_CHAR char_t;
typedef ROBO_CONST_STRING cstr;
\end{lstlisting}

Макросы \texttt{ROBO\_UNICODE\_ENABLED} управляют использованием широких символов (wchar\_t). В зависимости от этого макросы \texttt{RT} (строковый литерал) раскрываются либо в обычную строку, либо в L"string".

Также определяются макросы отладки и ассертов:
\begin{lstlisting}
#define ROBO_APP_CRASH() robo::crash(__func__, __FILE__, __LINE__)
#define ROBO_APP_ASSERT(x) if(!(x)) ROBO_APP_CRASH()
\end{lstlisting}

Функция \texttt{robo::crash} останавливает выполнение с сохранением места ошибки.

\subsubsection{Вспомогательные функции}

\begin{itemize}
    \item \texttt{int32\_t hash(cstr src, int32\_t begin = 0)} -- хеш-функция (простая аддитивная с xor).
    \item \texttt{template<typename T> constexpr T pi, deg2rad, rad2deg, g} -- математические константы.
    \item \texttt{template<typename T> constexpr T csqrt(T x)} -- рекурсивное вычисление квадратного корня (constexpr).
    \item Классы \texttt{ostram} / \texttt{istram} -- простые потоки для сериализации в память.
    \item \texttt{template<typename T> constexpr T abs, min, max, saturate} -- утилиты.
    \item Пространство \texttt{digit} -- функции для округления при сдвигах (\texttt{rsh}, \texttt{lsh}, \texttt{round}, \texttt{pack}).
    \item Реализация \texttt{catan2} -- аппроксимация арктангенса.
    \item \texttt{generator\_t} -- генератор таблиц для синуса (например, для ШИМ).
    \item \texttt{debouncer\_t} -- антидребезг контактов.
\end{itemize}

\subsection{Системное окружение (\texttt{robosd\_system.hpp}, \texttt{robosd\_system.cpp})}

Это центральная часть ядра, обеспечивающая абстракцию над аппаратурой и ОС.

\subsubsection{Класс \texttt{system}}

Класс \texttt{robo::system} реализован как синглтон (через \texttt{instance\_()}) и предоставляет статические методы для доступа ко всем сервисам.

Основные группы методов:
\begin{itemize}
    \item Управление временем: \texttt{clock\_period\_us()}, \texttt{time\_us()}, \texttt{time\_ms()}.
    \item Жизненный цикл: \texttt{begin()}, \texttt{start(period\_us)}, \texttt{stop()}, \texttt{finish()}.
    \item Обработчики циклов: \texttt{backend\_loop()}, \texttt{frontend\_loop()}.
    \item Печать и форматирование: \texttt{printf()}, \texttt{sprintf()}.
    \item Подсистемы: \texttt{ini}, \texttt{lib}, \texttt{shared}, \texttt{consol}.
    \item Служебные классы синхронизации: \texttt{guard}, \texttt{critical}, \texttt{fall}, \texttt{lazzyboy}.
\end{itemize}

\subsubsection{Синхронизация}

Класс \texttt{system::guard} автоматически определяет контекст вызова и выполняет соответствующие действия:
\begin{itemize}
    \item Если вызван из frontend -- входит в критическую секцию, останавливая backend и другие frontend-потоки.
    \item Если вызван из backend -- только блокирует доступ к разделяемым данным (lock/unlock), не приостанавливая себя.
\end{itemize}

Класс \texttt{system::critical} используется только во frontend и защищает данные от одновременного доступа нескольких frontend-потоков.

Класс \texttt{system::fall} предназначен для обозначения текущего потока как backend. В конструкторе он ожидает начала следующего такта (если включена квази-синхронизация) и вызывает \texttt{env::fall()}.

\subsubsection{Окружение (\texttt{system::env})}

Класс \texttt{env} содержит чистые виртуальные (или переопределяемые) методы, которые должны быть реализованы для каждой целевой платформы. В каталоге \texttt{core/platform} находятся готовые реализации:
\begin{itemize}
    \item \texttt{win+vs19} -- Windows, Visual Studio 2019.
    \item \texttt{ubuntu} -- Linux (Ubuntu).
    \item \texttt{raspberry+debain+visualdbg} -- Raspberry Pi.
    \item \texttt{keil+cpp6} -- Keil для ARM.
    \item \texttt{robosd\_system\_std.cpp} -- реализация на стандартной библиотеке C++ (для тестирования).
    \item \texttt{robosd\_system\_weak.cpp} -- слабые символы (weak) для переопределения пользователем.
\end{itemize}

Каждая реализация предоставляет:
\begin{itemize}
    \item \texttt{begin()}, \texttt{finish()} -- инициализация и завершение.
    \item \texttt{start()}, \texttt{stop()} -- запуск/остановка периодического таймера.
    \item \texttt{is\_frontend()}, \texttt{is\_backend()} -- определение контекста.
    \item \texttt{critical\_enter()}, \texttt{critical\_leave()} -- для frontend.
    \item \texttt{enter()}, \texttt{leave()} -- блокировка backend'а.
    \item \texttt{lock()}, \texttt{unlock()} -- быстрая блокировка для backend'а.
    \item \texttt{fall()}, \texttt{comeback()} -- переключение в режим backend.
    \item \texttt{realtime\_us()} -- высокоточное время.
    \item \texttt{rand()}, \texttt{sleep()}, \texttt{wakeup()} -- утилиты.
    \item \texttt{print()} -- низкоуровневый вывод.
    \item \texttt{sprintf()} -- форматирование (если не используется стандартное).
    \item \texttt{mem\_alloc()}, \texttt{mem\_free()}, \texttt{mem\_size()} -- аллокация (переопределяется для backend'а).
\end{itemize}

\subsubsection{Подсистемы}

\textbf{INI-файлы} (\texttt{system::ini}) -- чтение конфигурации из .ini-файлов. Реализации под Windows (через \texttt{GetPrivateProfileString}), Linux (через \texttt{minIni}) и встроенный парсер (\texttt{robosd\_ini\_parser.*}) для native-сборок.

\textbf{Динамические библиотеки} (\texttt{system::lib}) -- загрузка разделяемых библиотек (DLL/so) и получение адресов функций. Для Windows используются \texttt{LoadLibrary}/\texttt{GetProcAddress}, для Linux -- \texttt{dlopen}/\texttt{dlsym}, а также есть "нативная" реализация для статической компоновки (\texttt{robosd\_system\_native.cpp}).

\textbf{Разделяемая память} (\texttt{system::shared}) -- создание и отображение разделяемой памяти с синхронизацией (мьютекс). Реализации для Windows (File Mapping + Mutex) и Linux (shm\_open + mmap + sem\_wait?).

\textbf{Консоль} (\texttt{system::consol}) -- обработка сигналов завершения (Ctrl+C, закрытие окна). Устанавливает обработчики, вызывает пользовательскую функцию обратного вызова.

\subsection{Строки (\texttt{robosd\_string.hpp}, \texttt{robosd\_string.cpp})}

Класс \texttt{robo::string} предоставляет удобную работу со строками с поддержкой Unicode (если включено). Внутри используется \texttt{std::basic\_string<wchar\_t>} или \texttt{std::string}.

Основные возможности:
\begin{itemize}
    \item Конструкторы от C-строки, от форматной строки (с \texttt{va\_list} или переменным числом аргументов).
    \item Загрузка из INI-файла: \texttt{load(section, key)}.
    \item Преобразование в числа: \texttt{to\_number<T>()}, \texttt{to\_number\_array<T>()}, \texttt{scan\_numbers()}.
    \item Конвертация в ASCII (если строка в Unicode): \texttt{ascii()}, \texttt{ascii(char*)}.
    \item Форматирование с использованием внутренних буферов (frontend/backend), что исключает конфликты при параллельном вызове.
    \item Статические методы \texttt{sprintf\_backend\_} и \texttt{sprintf\_frontend\_} для безопасного форматирования в соответствующих контекстах.
\end{itemize}

Для ускорения преобразования чисел в строки предусмотрен метод \texttt{utoa\_n} (если включён \texttt{ROBO\_STRING\_UTOA\_ENABLED}).

\subsection{Модули и регистрация (\texttt{robosd\_app.hpp}, \texttt{robosd\_system\_module\_reg.hpp})}

Фреймворк поддерживает динамическую загрузку модулей (например, для драйверов устройств). Каждый модуль должен реализовать класс, наследующий \texttt{robo::app::module}, и предоставить функции \texttt{robo\_module\_query()} и \texttt{robo\_module\_release()}.

Макросы в \texttt{robosd\_system\_module\_reg.hpp} автоматически генерируют эти функции с учётом типа компоновки:
\begin{itemize}
    \item При сборке как динамической библиотеки -- экспортируются стандартные C-функции.
    \item При статической компоновке (native) -- используется класс \texttt{robo::native::lib}, который регистрирует модули в глобальном списке.
\end{itemize}

Класс \texttt{robo::native::lib} (из \texttt{robosd\_system\_native.hpp}) позволяет статически объявить библиотеку с набором процедур и затем обращаться к ним через \texttt{system::lib::proc\_get()}, что имитирует динамическую загрузку.

\subsection{Вспомогательные компоненты}

В \texttt{core} также находятся (будут описаны после предоставления соответствующих файлов):
\begin{itemize}
    \item \texttt{robosd\_delegat.hpp} -- делегаты и подписка на события.
    \item \texttt{robosd\_list.hpp} -- интрузивные списки (для эффективного управления объектами).
    \item \texttt{robosd\_log.hpp/cpp} -- система логирования с уровнями важности.
    \item \texttt{robosd\_ring\_buf.hpp} -- кольцевые буферы.
    \item \texttt{robosd\_stateflow.hpp} -- конечные автоматы.
    \item \texttt{robosd\_matrix.hpp} -- матричные операции.
    \item \texttt{robosd\_tree.hpp} -- деревья.
    \item \texttt{robosd\_zconveyer.hpp} -- конвейер для передачи данных между контекстами.
    \item \texttt{robosd\_autonum.hpp} -- автоматическое управление номерами объектов.
    \item \texttt{robosd\_crc.hpp} -- расчёт CRC.
\end{itemize}

\section{Заключение}

Фреймворк \texttt{robosd} представляет собой мощный инструмент для построения встраиваемых систем реального времени, сочетающий гибкость настройки под различные платформы и строгую дисциплину разделения контекстов. Ядро (\texttt{core}) предоставляет все необходимые абстракции, на которых строятся более высокоуровневые модули (управление двигателями, сетевое взаимодействие, пользовательские интерфейсы).

В следующих разделах документации будут подробно рассмотрены подсистемы \texttt{burst}, \texttt{net}, \texttt{servo} и другие, а также приведены примеры использования.

\section{Каталог файлов ядра (\texttt{core})}

Ниже приведён подробный перечень всех файлов, входящих в состав ядра \texttt{robosd}, с указанием определённых в них классов, функций и основных макросов. Описание основано исключительно на предоставленных исходных кодах.

\subsection{\texttt{robosd\_common.hpp}}

Главный заголовочный файл с общими определениями.

\begin{itemize}
    \item Подключает \texttt{robosd\_app\_tuning.hpp} и \texttt{robosd\_target.hpp}.
    \item Определяет типы: \texttt{random\_t}, \texttt{time\_us\_t}, \texttt{time\_ms\_t}, \texttt{char\_t}, \texttt{cstr}.
    \item Макросы \texttt{RT} (строковый литерал с учётом Unicode) и \texttt{RT8}.
    \item Функции: \texttt{hash} (две перегрузки), \texttt{crash}.
    \item Пространство имён \texttt{robo} с математическими константами: \texttt{pi}, \texttt{deg2rad}, \texttt{rad2deg}, \texttt{g}.
    \item Шаблонные функции: \texttt{csqrt}, \texttt{parity}, \texttt{abs}, \texttt{min}, \texttt{max}, \texttt{fma}, \texttt{saturate}.
    \item Классы: \texttt{ostram}, \texttt{istram} -- простые потоки для сериализации.
    \item Пространство \texttt{digit} -- функции для округления при сдвигах: \texttt{rsh}, \texttt{lsh}, \texttt{round}, \texttt{pack}, \texttt{select}.
    \item Шаблон \texttt{generator\_t} -- генератор таблиц для синуса.
    \item Шаблон \texttt{debouncer\_t} -- антидребезг контактов.
    \item Шаблон \texttt{math\_t} -- обёртки над математическими функциями (cos, sin, sqrt) для разных типов.
    \item Макросы отладки: \texttt{ROBO\_APP\_CRASH}, \texttt{ROBO\_APP\_ASSERT}.
    \item Определения типов платформ: \texttt{ROBO\_APP\_TYPE\_NONE}, \texttt{ROBO\_APP\_TYPE\_NATIVE}, \texttt{ROBO\_APP\_TYPE\_STD}, \texttt{ROBO\_APP\_TYPE\_WIN}, \texttt{ROBO\_APP\_TYPE\_LINUX} и др.
\end{itemize}

\subsection{\texttt{robosd\_system.hpp}}

Объявление класса \texttt{system} и вложенных классов.

\begin{itemize}
    \item \texttt{system::guard} -- класс для синхронизации контекстов (frontend/backend).
    \item \texttt{system::critical} -- критическая секция для frontend.
    \item \texttt{system::lazzyboy} -- измерение времени простоя.
    \item \texttt{system::fall} -- маркер входа в backend.
    \item \texttt{system::env} -- абстрактный интерфейс окружения (статические методы).
    \item \texttt{system::ini} -- работа с INI-файлами.
    \item \texttt{system::lib} -- работа с динамическими библиотеками.
    \item \texttt{system::shared} -- разделяемая память.
    \item \texttt{system::consol} -- обработка консольных сигналов.
    \item \texttt{system::timelog\_driver}, \texttt{timelogger} -- профилирование времени.
    \item Статические методы: \texttt{begin}, \texttt{finish}, \texttt{start}, \texttt{stop}, \texttt{backend\_loop}, \texttt{frontend\_loop}, \texttt{time\_us}, \texttt{time\_ms}, \texttt{clock\_period\_us}, \texttt{printf}, \texttt{sprintf}, \texttt{rand}.
\end{itemize}

\subsection{\texttt{robosd\_system.cpp}}

Реализация методов \texttt{system} и вспомогательных функций.

\begin{itemize}
    \item \texttt{hash} -- хеш-функция.
    \item \texttt{crash} -- аварийный останов.
    \item Пространство имён \texttt{sysclock} -- управление системным временем.
    \item Реализация \texttt{system::guard}, \texttt{system::critical}, \texttt{system::lazzyboy}, \texttt{system::fall}.
    \item Внутренний аллокатор \texttt{allocator} (при \texttt{ROBO\_APP\_ALLOC\_ENABLED}).
    \item Методы \texttt{system::mem\_alloc\_}, \texttt{system::mem\_free\_}.
    \item Реализация \texttt{system::printf}, \texttt{system::sprintf}.
    \item При \texttt{ROBO\_APP\_SHARED\_ENABLED} -- реализация \texttt{system::shared}.
    \item При \texttt{ROBO\_APP\_CONSOL\_ENABLED} -- реализация \texttt{system::consol}.
\end{itemize}

\subsection{\texttt{robosd\_system\_weak.cpp}}

Слабые реализации методов \texttt{system::env} и других подсистем, которые могут быть переопределены пользователем. Содержит пустые или заглушечные версии всех функций окружения, аллокатора, INI, библиотек, консоли и т.д., помеченные \texttt{ROBO\_WEAK}.

\subsection{\texttt{robosd\_system\_std.cpp}}

Реализация \texttt{system::env} на основе стандартной библиотеки C++.

\begin{itemize}
    \item \texttt{system::env::begin}, \texttt{finish} -- пустые.
    \item \texttt{system::env::start}, \texttt{stop} -- работа с таймером через \texttt{std::chrono}.
    \item \texttt{system::env::is\_frontend}, \texttt{is\_backend} -- сравнение идентификаторов потоков.
    \item \texttt{system::env::fall}, \texttt{comeback} -- установка идентификатора backend-потока.
    \item \texttt{system::env::realtime\_us} -- через \texttt{std::chrono::steady\_clock}.
    \item \texttt{system::env::rand} -- \texttt{std::rand}.
    \item \texttt{system::env::sleep} -- \texttt{std::this\_thread::yield}.
    \item \texttt{system::env::print} -- \texttt{wprintf} или \texttt{printf}.
    \item \texttt{system::env::abort} -- \texttt{std::abort}.
    \item \texttt{system::env::critical\_enter/leave}, \texttt{enter/leave}, \texttt{lock/unlock} -- через \texttt{std::recursive\_mutex}.
\end{itemize}

\subsection{\texttt{robosd\_system\_win.cpp}}

Реализация \texttt{system::env} для Windows.

\begin{itemize}
    \item Функции установки цвета консоли.
    \item \texttt{system::env::print} с цветами для разных уровней логирования.
    \item Работа с высокой точностью времени через \texttt{QueryPerformanceCounter}.
    \item Управление приоритетом потока (\texttt{SetThreadPriority}).
    \item Критические секции через \texttt{CRITICAL\_SECTION}.
    \item Работа с библиотеками: \texttt{LoadLibrary}, \texttt{GetProcAddress}, \texttt{FreeLibrary}, \texttt{CopyFile}, \texttt{DeleteFile}.
    \item Разделяемая память: \texttt{CreateFileMapping}, \texttt{MapViewOfFile}, \texttt{CreateMutex}.
    \item Обработка консоли: \texttt{SetConsoleCtrlHandler}.
    \item Аллокатор: \texttt{malloc}, \texttt{free}, \texttt{\_msize}.
    \item Таймер для профилирования: \texttt{QueryPerformanceCounter}.
\end{itemize}

\subsection{\texttt{robosd\_system\_linux.cpp} и \texttt{robosd-system-ubuntu.cpp}}

Реализации для Linux. Содержат:
\begin{itemize}
    \item Проверка существования файла через \texttt{stat} или \texttt{fopen}.
    \item Обработка сигналов (\texttt{sigaction}) для консоли.
    \item Работа с библиотеками: \texttt{dlopen}, \texttt{dlsym}, \texttt{dlclose}, копирование через \texttt{std::filesystem}.
    \item Разделяемая память: \texttt{shm\_open}, \texttt{mmap}, \texttt{munmap}, \texttt{shm\_unlink}.
    \item INI-файлы через \texttt{minIni} (в \texttt{robosd\_system\_linux.cpp}) или через собственную реализацию.
    \item Управление приоритетом через \texttt{pthread\_setschedparam} (закомментировано).
\end{itemize}

\subsection{\texttt{robosd\_system\_native.cpp} и \texttt{robosd\_system\_native.hpp}}

Реализация "нативной" подсистемы библиотек и INI для статической компоновки.

\begin{itemize}
    \item Класс \texttt{native::lib} -- представляет статическую библиотеку, регистрирует свои процедуры в глобальном списке.
    \item Структура \texttt{proc} -- пара (имя, указатель на функцию).
    \item Методы \texttt{reg}, \texttt{unreg}, \texttt{find\_by\_name\_}, \texttt{proc\_get}.
    \item Статические методы \texttt{system::lib::exists}, \texttt{load}, \texttt{free}, \texttt{proc\_get} делегируют в \texttt{native::lib}.
    \item INI-реализация через встроенный парсер \texttt{robosd\_ini\_parser.h} (загружает весь файл в память и парсит).
\end{itemize}

\subsection{\texttt{robosd\_string.hpp} и \texttt{robosd\_string.cpp}}

Класс \texttt{string} для работы со строками.

\begin{itemize}
    \item Внутреннее представление -- \texttt{std::basic\_string<char\_t>}.
    \item Конструкторы: по умолчанию, копирования, из форматной строки (с \texttt{va\_list} или \texttt{...}).
    \item Методы: \texttt{format}, \texttt{load} (из INI), \texttt{tryload}, \texttt{to\_number}, \texttt{to\_number\_array}, \texttt{scan\_numbers}, \texttt{ascii} (конвертация в UTF-8), \texttt{length}, \texttt{c\_str}, \texttt{clear}, операторы присваивания и сравнения.
    \item Статические буферы \texttt{string\_buffer\_frontend} и \texttt{string\_buffer\_backend} для безопасного форматирования в разных контекстах.
    \item Статические методы \texttt{sprintf\_backend\_} и \texttt{sprintf\_frontend\_} для форматирования в соответствующие буферы.
    \item \texttt{utoa\_n} -- быстрое преобразование целого в строку (если включено).
\end{itemize}

\subsection{\texttt{robosd\_app.hpp} и \texttt{robosd\_app.cpp}}

Реализация каркаса приложения на основе узлов (\texttt{node}) и поддержка модулей.

\begin{itemize}
    \item \texttt{node} -- базовый класс узла. Содержит:
        \begin{itemize}
            \item Перечисление \texttt{state} (unknown, panic, clean, stopped, startup, execute, shutdown).
            \item Списки дочерних узлов по состояниям: \texttt{disabled\_}, \texttt{stopped\_}, \texttt{startupped\_}, \texttt{active\_}, \texttt{shutdowned\_}.
            \item Методы жизненного цикла: \texttt{load()}, \texttt{start()}, \texttt{startup()}, \texttt{stop()}, \texttt{shutdown()}, \texttt{clean()}, \texttt{panic()}.
            \item Виртуальные методы \texttt{do\_load}, \texttt{do\_clean}, \texttt{do\_start}, \texttt{do\_stop}, \texttt{do\_startup}, \texttt{do\_shutdown}, \texttt{do\_panic}.
            \item Поддержка пути в конфигурации (вложенный класс \texttt{path}).
            \item Алиас (\texttt{alias\_}) из INI.
            \item Статические методы \texttt{find} для поиска узла по хешу.
        \end{itemize}
    \item \texttt{module} -- наследник \texttt{node} с чисто виртуальными методами \texttt{frontend\_loop()} и \texttt{backend\_loop()}.
    \item \texttt{wrapper} -- обёртка для динамически загружаемого модуля. Содержит \texttt{handle\_} (библиотека) и \texttt{module\_} (экземпляр модуля). Методы \texttt{begin\_} и \texttt{finish} загружают/выгружают модуль.
    \item \texttt{machine} -- корневой узел-синглтон. Управляет всем приложением: инициализация системы, загрузка INI, запуск модулей, циклы frontend/backend. Методы: \texttt{begin}, \texttt{start}, \texttt{stop}, \texttt{frontend\_loop}, \texttt{backend\_loop}, \texttt{terminated}.
\end{itemize}

\subsection{\texttt{robosd\_autonum.hpp} и \texttt{robosd\_autonum.cpp}}

Механизм автоматического удаления делегатов (сборка мусора для делегатов).

\begin{itemize}
    \item Класс \texttt{receicledbin::ref} -- базовый класс для объектов с автоматическим удалением. Содержит счётчик использований и флаг контекста (frontend/backend). При обнулении счётчика объект помещается в соответствующую корзину (\texttt{rbin}).
    \item Статические методы \texttt{frontend\_clean()} и \texttt{backend\_clean()} очищают корзины.
    \item Шаблон \texttt{fabric<B,R,Args...>} -- фабрика для создания делегатов, наследующих одновременно от \texttt{B} и \texttt{receicledbin::ref}. Предоставляет вложенные классы \texttt{simple}, \texttt{uni}, \texttt{lambda}, \texttt{member} (аналогично \texttt{delegat}).
    \item Псевдоним \texttt{autonum\_fabric} и функция \texttt{ausimple} для удобства.
\end{itemize}

\subsection{\texttt{robosd\_delegat.hpp}}

Реализация различных видов делегатов.

\begin{itemize}
    \item Базовый шаблон \texttt{ref<R,Args...>} -- интерфейс делегата.
    \item \texttt{dummy} -- заглушка.
    \item \texttt{simple} -- делегат от обычной функции.
    \item \texttt{uni} -- делегат от функции с первым параметром \texttt{void*} (объект).
    \item \texttt{lambda} -- делегат от \texttt{robo::lambda}.
    \item \texttt{member} и \texttt{rmember} -- делегаты от метода класса (по указателю и по ссылке).
    \item Пространство имён \texttt{owned} -- фабрики для создания делегатов с автоматическим управлением памятью (без автоудаления).
\end{itemize}

\subsection{\texttt{robosd\_lambda.hpp}}

Реализация функционального объекта \texttt{lambda}, способного хранить произвольные функторы.

\begin{itemize}
    \item Внутренний класс \texttt{context} хранит скопированный объект и указатель на функцию-запускатель.
    \item \texttt{lambda} содержит указатель на \texttt{context} с подсчётом ссылок.
    \item Поддержка копирования и присваивания.
    \item Оператор вызова \texttt{operator()} запускает сохранённый функтор.
\end{itemize}

\subsection{\texttt{robosd\_list.hpp}}

Библиотека интрузивных списков.

\begin{itemize}
    \item \texttt{base\_ref<T>} -- базовый класс для элемента списка, содержит указатели next/prev и ссылку на владельца.
    \item \texttt{base<T>} -- базовый двунаправленный список, управляет элементами.
    \item \texttt{unsorted<T>} -- неупорядоченный список.
    \item \texttt{sorted<T,K>} -- упорядоченный список с ключом.
    \item \texttt{unique<T,K>} -- список с уникальными ключами (метод \texttt{find}).
    \item \texttt{pool<T,K>} -- неупорядоченный список с ключом.
    \item В пространстве имён \texttt{unidir} -- однонаправленные списки:
        \begin{itemize}
            \item \texttt{base\_t<T,I>} -- база с методами push, pop\_first, extract.
            \item \texttt{store\_t<T,Arg...>} -- хранилище, создающее элементы через new.
            \item \texttt{queue\_t<T,Arg...>} -- очередь (FIFO) на основе \texttt{base\_t}.
            \item \texttt{stack\_t<T,Arg...>} -- стек (LIFO).
        \end{itemize}
    \item В пространстве имён \texttt{queue} -- адаптеры для использования списков в качестве очередей.
\end{itemize}

\subsection{\texttt{robosd\_log.hpp} и \texttt{robosd\_log.cpp}}

Система логирования.

\begin{itemize}
    \item Перечисление \texttt{log::verb} -- уровни важности.
    \item Структура \texttt{log::mask} -- битовые маски.
    \item Глобальные переменные \texttt{verb\_}, \texttt{mask\_}.
    \item Функция \texttt{log::print} -- форматированный вывод с учётом уровня и маски.
    \item \texttt{log::begin}, \texttt{log::finish}.
    \item Множество макросов для логирования и проверок: \texttt{robo\_errlog}, \texttt{robo\_warninglog}, \texttt{robo\_infolog}, \texttt{robo\_detaillog}, \texttt{ROBO\_BREAKN\_F}, \texttt{ROBO\_JAMPN\_F}, \texttt{ROBO\_ASSERT\_F}, \texttt{ROBO\_ALARMN} и др.
\end{itemize}

\subsection{\texttt{robosd\_ini.hpp}}

Шаблонные функции для загрузки параметров из INI.

\begin{itemize}
    \item \texttt{try\_load<T>} -- загружает значение и преобразует в тип T.
    \item \texttt{load<T>} -- аналогично, но с логированием ошибки.
    \item \texttt{try\_load\_arr<T>} -- загружает массив чисел.
    \item \texttt{try\_load\_list<T>} -- загружает список чисел (с подсчётом количества).
\end{itemize}

\subsection{\texttt{robosd\_alocator.cpp}}

Переопределение глобальных операторов \texttt{new} и \texttt{delete}.

\begin{itemize}
    \item \texttt{operator new} вызывает \texttt{robo::system::mem::alloc} (если \texttt{ROBO\_APP\_ALLOC\_ENABLED}) или \texttt{malloc}.
    \item \texttt{operator delete} вызывает \texttt{robo::system::mem::free} или \texttt{free}.
\end{itemize}

\subsection{\texttt{robosd\_system\_module\_reg.hpp}}

Макросы для автоматической генерации функций входа в модуль.

\begin{itemize}
    \item Если \texttt{ROBO\_APP\_LIB\_TYPE == ROBO\_APP\_TYPE\_NATIVE} -- создаётся экземпляр \texttt{native::lib} с функциями \texttt{query} и \texttt{release}.
    \item Иначе -- генерируются extern "C" функции \texttt{robo\_module\_query} и \texttt{robo\_module\_release} с атрибутами экспорта.
    \item Использует макросы \texttt{MODULE\_NAME}, \texttt{MODULE\_NAME\_STR}, \texttt{MODULE\_NAME\_PREFIX}.
\end{itemize}

\section{Макросы конфигурации (\texttt{robosd\_app\_tuning.hpp})}

Для настройки фреймворка пользователь должен определить ряд макросов в файле \texttt{robosd\_app\_tuning.hpp} (или в командной строке компилятора). Ниже приведён полный список макросов, которые встречаются в предоставленных исходных кодах и могут быть определены пользователем.

\subsection{Основные макросы включения функциональности}

\begin{description}
    \item[\texttt{ROBO\_APP\_SYSTEM\_ENABLED}] -- если 1, включает системный слой (класс \texttt{system}). По умолчанию 0.
    \item[\texttt{ROBO\_APP\_MODULE\_ENABLED}] -- если 1, включает поддержку модулей (динамическую загрузку). По умолчанию 0.
    \item[\texttt{ROBO\_APP\_DEBUG\_LOG\_ENABLED}] -- если 1, включает систему логирования. По умолчанию 0.
    \item[\texttt{ROBO\_APP\_DEBUG\_ENABLED}] -- если 1, включает отладочные ассерты. По умолчанию 1.
    \item[\texttt{ROBO\_UNICODE\_ENABLED}] -- если 1, используются широкие символы (\texttt{wchar\_t}). По умолчанию 0.
    \item[\texttt{ROBO\_AUTONUM\_ENABLED}] -- если 1, включает механизм автоудаления делегатов. По умолчанию 1.
    \item[\texttt{ROBO\_APP\_RTTI\_ENABLED}] -- если 1, разрешает использование RTTI (например, для \texttt{dynamic\_cast}). По умолчанию 1.
    \item[\texttt{ROBO\_APP\_ULTRACOMPACT}] -- если 1, включает ультракомпактный режим (не используется в предоставленном коде, но упомянут).
    \item[\texttt{ROBO\_APP\_SERVO\_ENABLED}] -- если 1, включает подсистему servo (используется в \texttt{robosd\_app.cpp}).
    \item[\texttt{ROBO\_APP\_TERMINAL\_ENABLED}] -- если 1, включает терминал.
    \item[\texttt{ROBO\_APP\_TRACE\_ENABLED}] -- если 1, включает трассировку.
    \item[\texttt{ROBO\_APP\_QUAZZY\_REALTIME\_ENABLED}] -- если 1, включает принудительную синхронизацию в \texttt{fall}.
    \item[\texttt{ROBO\_APP\_SYSTEM\_GUARD\_DEBUG\_ENABLED}] -- если 1, ведётся подсчёт вложенности guard'ов для отладки.
\end{description}

\subsection{Макросы выбора типа реализации}

Для каждой подсистемы можно задать конкретную реализацию с помощью макросов вида \texttt{ROBO\_APP\_XXXX\_TYPE}. Допустимые значения:

\begin{itemize}
    \item \texttt{ROBO\_APP\_TYPE\_NONE} (0) -- отключено.
    \item \texttt{ROBO\_APP\_TYPE\_NATIVE} (1) -- встроенная реализация.
    \item \texttt{ROBO\_APP\_TYPE\_STD} (2) -- реализация на стандартной библиотеке C++.
    \item \texttt{ROBO\_APP\_TYPE\_WIN} (3) -- Windows API.
    \item \texttt{ROBO\_APP\_TYPE\_LINUX} (4) -- Linux (POSIX).
    \item \texttt{ROBO\_APP\_TYPE\_RASPBERRY} (5) -- Raspberry Pi.
    \item \texttt{ROBO\_APP\_TYPE\_UBUNTU} (6) -- Ubuntu.
    \item \texttt{ROBO\_APP\_TYPE\_ASTRA} (7) -- Astra Linux.
    \item \texttt{ROBO\_APP\_TYPE\_SPECIFIC} (8) -- пользовательская реализация (требуется определить weak-символы).
    \item \texttt{ROBO\_APP\_TYPE\_DUMMY} (9) -- заглушка.
\end{itemize}

Список макросов выбора:

\begin{description}
    \item[\texttt{ROBO\_APP\_OS\_TYPE}] -- операционная система.
    \item[\texttt{ROBO\_APP\_ENV\_TYPE}] -- окружение (\texttt{system::env}).
    \item[\texttt{ROBO\_APP\_INI\_TYPE}] -- реализация INI-файлов.
    \item[\texttt{ROBO\_APP\_LIB\_TYPE}] -- динамические библиотеки.
    \item[\texttt{ROBO\_APP\_SHARED\_TYPE}] -- разделяемая память.
    \item[\texttt{ROBO\_APP\_ALLOC\_TYPE}] -- аллокатор.
    \item[\texttt{ROBO\_APP\_PRINT\_TYPE}] -- низкоуровневый вывод.
    \item[\texttt{ROBO\_APP\_LOG\_PRINT\_TYPE}] -- вывод для логирования.
    \item[\texttt{ROBO\_APP\_FORMATING\_TYPE}] -- форматированный вывод.
    \item[\texttt{ROBO\_APP\_CONSOL\_TYPE}] -- обработка консоли.
    \item[\texttt{ROBO\_APP\_CRITICAL\_TYPE}] -- критические секции.
    \item[\texttt{ROBO\_APP\_TIMELOG\_TYPE}] -- высокоточный таймер.
    \item[\texttt{ROBO\_APP\_REALTIME\_TYPE}] -- переключение приоритета.
    \item[\texttt{ROBO\_APP\_ABORT\_TYPE}] -- аварийный останов.
\end{description}

\subsection{Макросы для настройки размеров}

\begin{description}
    \item[\texttt{ROBO\_STRING\_BUFFER\_SIZE}] -- размер временных буферов для строк. По умолчанию 255, но в некоторых местах переопределён как 4096.
    \item[\texttt{ROBO\_STRING\_UTOA\_TYPE}] -- тип реализации \texttt{utoa} (NATIVE или NONE).
    \item[\texttt{ROBO\_APP\_NODE\_PATH\_SIZE}] -- максимальная длина пути узла. По умолчанию 1024.
    \item[\texttt{ROBO\_APP\_BACKEND\_PRINT\_BUFFER\_BITS}] -- размер кольцевого буфера для печати из backend (2\textsuperscript{n}). По умолчанию 10 (т.е. 1024).
    \item[\texttt{ROBO\_APP\_TIMELOG\_SIZE}] -- размер буфера временного лога. По умолчанию 100000.
    \item[\texttt{ROBO\_SYSTEM\_ALLOCATOR\_WORD\_TYPE}] -- тип машинного слова аллокатора. По умолчанию \texttt{uint32\_t}.
    \item[\texttt{ROBO\_SYSTEM\_ALLOCATOR\_SIZE\_TYPE}] -- тип для хранения размеров в аллокаторе. По умолчанию \texttt{uint16\_t}.
    \item[\texttt{ROBO\_SYSTEM\_ALLOCATOR\_SIZE}] -- общий размер пула в словах. По умолчанию 1024.
    \item[\texttt{ROBO\_SYSTEM\_ALLOCATOR\_WORD\_SIZE\_BITS}] -- логарифм размера слова в байтах. По умолчанию 2 (т.е. 4 байта).
\end{description}

\subsection{Прочие макросы}

\begin{description}
    \item[\texttt{ROBO\_APP\_DMA\_BUFFER}] -- может быть определён как атрибут памяти для DMA.
    \item[\texttt{ROBO\_APP\_ALIGN\_32BYTES(x)}] -- макрос для выравнивания данных (по умолчанию пустой).
    \item[\texttt{ROBO\_ZEROS\_STRUCT}] -- макрос для обнуления структур (пустой по умолчанию).
    \item[\texttt{ROBO\_EXPORT\_FUNCTION\_PREFIX}] -- префикс для экспортируемых функций (например, пустая строка или \texttt{"\_"}).
    \item[\texttt{ROBO\_APP\_MATH\_SHIFT\_ENABLE}] -- если 1, использовать аппаратные сдвиги для отрицательных чисел. По умолчанию 0.
    \item[\texttt{ROBO\_APP\_CRASH}] -- макрос для аварийного останова (по умолчанию \texttt{robo::crash}).
    \item[\texttt{ROBO\_APP\_ASSERT}] -- макрос для ассертов.
    \item[\texttt{ROBO\_TYPE\_RANDOM}] -- тип для случайных чисел. По умолчанию \texttt{int}.
    \item[\texttt{ROBO\_TYPE\_TIME\_US}] -- тип для микросекунд. По умолчанию \texttt{uint32\_t}.
    \item[\texttt{ROBO\_TYPE\_TIME\_MS}] -- тип для миллисекунд. По умолчанию \texttt{uint32\_t}.
\end{description}

\section{Детальное описание \texttt{robosd\_system.hpp}}

Файл \texttt{robosd\_system.hpp} является центральным заголовочным файлом системного слоя фреймворка \texttt{robosd}. Он определяет класс \texttt{robo::system}, который предоставляет унифицированный интерфейс для работы с аппаратным и программным окружением, а также набор вложенных классов и подсистем, обеспечивающих:

\begin{itemize}
    \item Абстракцию операционной системы и аппаратной платформы.
    \item Разделение контекстов выполнения на \textbf{frontend} (обычные потоки ОС) и \textbf{backend} (жёсткое реальное время).
    \item Примитивы синхронизации, учитывающие это разделение.
    \item Управление временем (системные часы, высокоточное время).
    \item Динамическое выделение памяти с возможностью использования пула для backend.
    \item Работу с конфигурационными INI-файлами.
    \item Загрузку динамических библиотек (модулей).
    \item Разделяемую память.
    \item Обработку консольных сигналов (Ctrl+C и т.п.).
    \item Профилирование времени выполнения.
    \item Низкоуровневый форматированный вывод и печать.
\end{itemize}

\subsection{Макросы конфигурации}

В начале файла определяются макросы, управляющие включением тех или иных подсистем. Они могут быть переопределены пользователем в файле \texttt{robosd\_app\_tuning.hpp}. Если не определены, им присваиваются значения по умолчанию \texttt{ROBO\_APP\_TYPE\_NONE} (подсистема отключена).

\begin{lstlisting}[caption={Макросы выбора реализации}]
#ifndef ROBO_APP_ENV_TYPE 
#define ROBO_APP_ENV_TYPE  ROBO_APP_TYPE_NONE
#endif
// ... аналогично для других подсистем
\end{lstlisting}

На основе этих макросов вычисляются флаги включения, например:
\begin{lstlisting}
#define ROBO_APP_ENV_ENABLED  (ROBO_APP_ENV_TYPE != ROBO_APP_TYPE_NONE)
\end{lstlisting}

Также определяется макрос \texttt{ROBO\_APP\_SYSTEM\_GUARD\_DEBUG\_ENABLED} (по умолчанию 0) для отладочного подсчёта вложенности guard'ов.

В конце файла определяются удобные макросы для использования в коде приложения:
\begin{itemize}
    \item \texttt{guard\_t}, \texttt{guard\_\_}, \texttt{fall\_\_}, \texttt{critical\_\_} -- для создания объектов синхронизации.
    \item \texttt{is\_frontend\_\_}, \texttt{is\_backend\_\_} -- для проверки контекста.
    \item \texttt{switch\_to\_realtime\_()}, \texttt{switch\_to\_normal\_()} -- для переключения приоритета потока (если поддерживается ОС).
\end{itemize}

\subsection{Класс \texttt{system}}

Класс \texttt{robo::system} реализован как синглтон (метод \texttt{instance\_()}) и содержит в себе все системные сервисы. Он объявляет дружественные классы и хранит состояние.

\subsubsection{Вложенная структура \texttt{mem}}

Предназначена для учёта статистики выделения памяти (включается макросом \texttt{ROBO\_APP\_ALLOC\_ENABLED}). Содержит вложенную структуру \texttt{stat} с полями:
\begin{itemize}
    \item \texttt{used} -- используемая память (размер и количество блоков).
    \item \texttt{total} -- общая выделенная память (полезная нагрузка, общий размер с учётом служебных данных, количество вызовов).
\end{itemize}
Методы:
\begin{itemize}
    \item \texttt{static void* alloc(size\_t \_sz)} -- выделяет память через внутренний механизм (\texttt{system::instance\_().mem\_alloc\_(\_sz)}).
    \item \texttt{static void free(void* \_memo)} -- освобождает память.
    \item \texttt{static mem::stat\& get\_mem\_statistic(void)} -- возвращает ссылку на статистику (с использованием guard для потокобезопасности).
\end{itemize}

\subsubsection{Класс \texttt{guard}}

Это основной примитив синхронизации, который автоматически определяет контекст вызова (frontend или backend) и выполняет соответствующие действия:
\begin{itemize}
    \item В контексте frontend: входит в критическую секцию (останавливая backend и другие frontend-потоки), затем блокирует backend (вызывая \texttt{env::enter()}) и переключает поток в режим реального времени (если поддерживается).
    \item В контексте backend: только блокирует доступ к разделяемым данным через \texttt{env::lock()}, не приостанавливая сам backend.
\end{itemize}
В конструкторе выполняются эти действия, в деструкторе -- обратные. Для отладки ведётся подсчёт вложенности (если включено).

Использование:
\begin{lstlisting}
{
    robo::system::guard g__; // or macro guard__
    // protected code
}
\end{lstlisting}

\subsubsection{Класс \texttt{critical}}

Используется только в контексте frontend для защиты от одновременного доступа нескольких frontend-потоков. В конструкторе вызывает \texttt{env::critical\_enter()}, в деструкторе -- \texttt{env::critical\_leave()}. В backend использовать нельзя.

Пример:
\begin{lstlisting}
void frontend_function() {
    robo::system::critical c__;
    // critical section between frontend threads
}
\end{lstlisting}

\subsubsection{Класс \texttt{lazzyboy}}

Позволяет измерять время простоя (или интервал) от момента создания объекта. Хранит начальное время (\texttt{sleep\_us\_}). Метод \texttt{idle\_us()} возвращает количество микросекунд, прошедших с создания. Метод \texttt{reset()} сбрасывает начальное время.

Использование:
\begin{lstlisting}
robo::system::lazzyboy timer;
// ... some code ...
robo::time_us_t elapsed = timer.idle_us();
\end{lstlisting}

\subsubsection{Класс \texttt{fall}}

Маркер, который поток должен создать при входе в цикл backend. В конструкторе:
\begin{itemize}
    \item Если включён режим \texttt{ROBO\_APP\_QUAZZY\_REALTIME\_ENABLED}, ожидает наступления следующего такта (синхронизация по таймеру).
    \item Вызывает \texttt{env::fall()}, который помечает текущий поток как backend.
\end{itemize}
В деструкторе вызывает \texttt{env::comeback()}, снимая пометку.

Использование:
\begin{lstlisting}
void backend_thread() {
    while (!robo::system::terminated()) {
        robo::system::fall f__;
        // backend loop body
        robo::system::backend_loop();
    }
}
\end{lstlisting}

\subsubsection{Класс \texttt{env}}

Это абстрактный интерфейс, который должен быть реализован для каждой конкретной платформы. Все методы статические. Они делятся на несколько групп:

\begin{description}
    \item[Управление контекстом:]
        \begin{itemize}
            \item \texttt{is\_frontend()}, \texttt{is\_backend()} -- определение текущего контекста.
            \item \texttt{fall()}, \texttt{comeback()} -- маркировка потока как backend.
        \end{itemize}
    \item[Синхронизация:]
        \begin{itemize}
            \item \texttt{critical\_enter()}, \texttt{critical\_leave()} -- для frontend.
            \item \texttt{enter()}, \texttt{leave()} -- блокировка backend'а.
            \item \texttt{lock()}, \texttt{unlock()} -- быстрая блокировка для backend.
        \end{itemize}
    \item[Управление приоритетом:] \texttt{switch\_to\_realtime()}, \texttt{switch\_to\_normal()} (если поддерживается ОС).
    \item[Инициализация и завершение:] \texttt{begin()}, \texttt{finish()}, \texttt{start()}, \texttt{stop()}, \texttt{backend\_loop()}, \texttt{frontend\_loop()}, \texttt{startup()}, \texttt{shutdown()} (последние два -- если включены модули).
    \item[Время:] \texttt{realtime\_us()} -- высокоточное время.
    \item[Случайные числа:] \texttt{rand()}, \texttt{rand\_maxd()}.
    \item[Управление потоком:] \texttt{wakeup()}, \texttt{sleep()}.
    \item[Форматирование и печать:] \texttt{sprintf()} (если не используется стандартная библиотека), \texttt{print()}.
    \item[Аварийный останов:] \texttt{abort()}, \texttt{reset()}.
    \item[Аллокация:] \texttt{mem\_alloc()}, \texttt{mem\_free()}, \texttt{mem\_size()} (если включена).
\end{description}

Реализации этих методов находятся в файлах \texttt{robosd\_system\_std.cpp}, \texttt{robosd\_system\_win.cpp}, \texttt{robosd\_system\_linux.cpp} и др., а также предоставлена слабая реализация в \texttt{robosd\_system\_weak.cpp}.

\subsubsection{Подсистема профилирования времени (\texttt{timelog})}

Для измерения временных характеристик кода (профилирования) в фреймворк встроены шаблонные классы, расположенные в заголовочном файле \texttt{prf/timelogger.hpp}. Они позволяют с высокой точностью фиксировать моменты начала и окончания событий, а также сохранять произвольные данные с привязкой ко времени.

Если макрос \texttt{ROBO\_APP\_TIMELOG\_ENABLED} установлен в 1, в \texttt{robosd\_system.hpp} подключается \texttt{prf/timelogger.hpp} и определяется:

\begin{itemize}
    \item Структура \texttt{timelogev} с предопределёнными событиями (begin, start, startup, halt) -- может использоваться как идентификаторы событий.
    \item Структура \texttt{timelog\_driver} -- должна предоставлять метод \texttt{time\_ns()} для получения текущего времени в наносекундах и объявлять \texttt{typedef system::guard guard} для синхронизации.
    \item Псевдоним \texttt{timelog\_s} -- конкретный экземпляр логгера событий: \texttt{prf::timelogger\_t<timelog\_driver, ROBO\_APP\_TIMELOG\_SIZE>}. Размер буфера по умолчанию 100000 (задаётся макросом \texttt{ROBO\_APP\_TIMELOG\_SIZE}).
    \item Макросы \texttt{robo\_logger\_raise(e)} и \texttt{robo\_logger\_fail(e)} для отметки начала и конца события.
    \item Класс \texttt{timelogger} -- RAII-обёртка: в конструкторе вызывает \texttt{raise}, в деструкторе -- \texttt{fail}.
\end{itemize}

Если подсистема отключена, все макросы превращаются в пустые операции.

\paragraph{Шаблон \texttt{timelogger\_t<P, N>}}

Предназначен для логирования событий с двумя состояниями (начало/конец). Внутреннее устройство:

\begin{itemize}
    \item Два буфера (\texttt{pool[2][N]}) типа \texttt{record\_s}, организованных по принципу двойной буферизации (double buffering). Это позволяет безопасно собирать данные в реальном времени (например, из прерываний или высокоприоритетного потока) без длительных блокировок.
    \item Структура \texttt{record\_s} содержит:
        \begin{itemize}
            \item \texttt{uint32\_t ns} -- время, прошедшее с предыдущего события (в наносекундах).
            \item битовое поле: \texttt{event:31} -- идентификатор события (до 2\textsuperscript{31}-1), \texttt{dir:1} -- направление (1 -- начало, 0 -- конец).
        \end{itemize}
    \item Структура \texttt{records\_s} объединяет счётчик записей и указатель на массив записей.
    \item Методы:
        \begin{itemize}
            \item \texttt{static void raise(uint32\_t event)} -- фиксирует начало события.
            \item \texttt{static void fail(uint32\_t event)} -- фиксирует конец события.
            \item \texttt{static const records\_s* get(void)} -- возвращает указатель на заполненный буфер и переключает на другой буфер для продолжения записи. Вызов \texttt{get()} должен быть защищён мьютексом (в реализации используется \texttt{P::guard}) и предполагается, что вызывающий (обычно frontend) обработает данные и освободит буфер до следующего переключения.
        \end{itemize}
\end{itemize}

Использование в коде:

\begin{lstlisting}
#include "core/robosd_system.hpp"

void some_function() {
    robo_logger(123); // macro creates timelogger object, which calls raise(123) in constructor and fail(123) in destructor
    // ... code ...
}

// Elsewhere, e.g., in frontend loop, periodically fetch data:
const auto records = robo::system::timelog_s::get();
for (int i = 0; i < records->count; ++i) {
    const auto& r = records->items[i];
    printf("Event %u dir %u time %u\n", r.event, r.dir, r.ns);
}
\end{lstlisting}

\paragraph{Шаблон \texttt{datalogger\_t<D, P, N>}}

Аналогичен \texttt{timelogger\_t}, но предназначен для сохранения произвольных данных типа \texttt{D} (например, мгновенных значений сигналов, положений, токов) с меткой времени. Структура \texttt{record\_s} содержит:

\begin{itemize}
    \item \texttt{uint32\_t ns} -- время от предыдущего сохранения.
    \item \texttt{D value} -- сохранённое значение.
\end{itemize}

Методы:
\begin{itemize}
    \item \texttt{static void save(const D\& value)} -- добавляет запись с текущим значением.
    \item \texttt{static const records\_s* get(void)} -- аналогично \texttt{timelogger\_t}.
\end{itemize}

Этот класс удобен для сбора осциллограмм работы регуляторов или других быстрых процессов.

Оба класса используют шаблонный параметр \texttt{P} (драйвер), который должен предоставлять:
\begin{itemize}
    \item \texttt{typedef ... guard} -- тип примитива синхронизации (обычно \texttt{system::guard}).
    \item \texttt{static uint32\_t time\_ns()} -- функция получения текущего времени в наносекундах.
\end{itemize}

Таким образом, подсистема профилирования легко адаптируется под разные платформы и требования точности.

\subsubsection{Подсистема INI-файлов (\texttt{ini})}

Структура \texttt{ini} предоставляет статические методы для работы с конфигурационными файлами:
\begin{itemize}
    \item \texttt{begin(cstr \_ini)} -- инициализация (открытие файла).
    \item \texttt{source()} -- возвращает имя текущего файла.
    \item \texttt{finish()} -- завершение работы.
    \item \texttt{load\_str()} -- загружает строковое значение.
    \item \texttt{load\_data()} -- загружает данные (возможно, бинарные) или список ключей (если \texttt{\_key == nullptr}).
\end{itemize}

Реализации зависят от платформы (Windows: \texttt{GetPrivateProfileString}, Linux: \texttt{minIni}, Native: встроенный парсер).

\subsubsection{Подсистема динамических библиотек (\texttt{lib})}

Структура \texttt{lib} предоставляет методы для работы с разделяемыми библиотеками:
\begin{itemize}
    \item \texttt{exists(cstr \_lib\_name)} -- проверяет существование файла библиотеки.
    \item \texttt{load(cstr \_lib\_name)} -- загружает библиотеку, возвращает opaque-указатель.
    \item \texttt{free(void* \_instance)} -- выгружает библиотеку.
    \item \texttt{proc\_get(void* \_handle, cstr \_proc\_name)} -- получает адрес функции по имени.
    \item \texttt{copy(cstr \_src, cstr \_dst)} -- копирует файл библиотеки.
    \item \texttt{remove(cstr \_lib\_name)} -- удаляет файл библиотеки.
\end{itemize}

Реализации под Windows используют \texttt{LoadLibrary} / \texttt{GetProcAddress}, под Linux -- \texttt{dlopen} / \texttt{dlsym}. Также есть "нативная" реализация для статической компоновки (см. \texttt{robosd\_system\_native.hpp}).

\subsubsection{Подсистема разделяемой памяти (\texttt{shared})}

Класс \texttt{shared} предоставляет доступ к разделяемой памяти между процессами. Содержит:
\begin{itemize}
    \item Внутренний класс \texttt{driver} (скрытая реализация).
    \item Методы \texttt{open(cstr \_path, size\_t \_sz)} -- создаёт или открывает существующий объект разделяемой памяти.
    \item \texttt{close()} -- закрывает.
    \item \texttt{memo()} -- возвращает указатель на отображённую память.
    \item \texttt{size()} -- возвращает размер.
    \item Вложенный класс \texttt{guard} для синхронизации доступа (блокирует мьютекс, связанный с объектом).
    \item Статические методы \texttt{find(cstr \_name)} и \texttt{find(int \_id)} для поиска уже открытых объектов по имени или хешу.
\end{itemize}

Реализации под Windows используют \texttt{CreateFileMapping} / \texttt{MapViewOfFile} + \texttt{CreateMutex}, под Linux -- \texttt{shm\_open} / \texttt{mmap} + семафоры.

\subsubsection{Подсистема консоли (\texttt{consol})}

Класс \texttt{consol} обрабатывает сигналы завершения (Ctrl+C, закрытие окна и т.п.):
\begin{itemize}
    \item Перечисление \texttt{event} -- тип события (\texttt{app}, \texttt{keypbrd}, \texttt{other}).
    \item Тип \texttt{on\_break\_f} -- функтор (\texttt{lambda<void(event)>}) для обратного вызова.
    \item Методы \texttt{driver\_begin()} и \texttt{driver\_finish()} -- платформозависимая установка/снятие обработчиков.
    \item \texttt{stop(event \_ev)} -- вызывает зарегистрированный обработчик.
    \item \texttt{begin(const on\_break\_f\& \_on\_break)} -- регистрирует обработчик и инициализирует драйвер.
    \item \texttt{finish()} -- завершает работу.
\end{itemize}

\subsubsection{Методы управления временем и циклами}

\begin{itemize}
    \item \texttt{clock\_period\_us()} -- возвращает период системного таймера (в микросекундах), установленный в \texttt{start()}.
    \item \texttt{time\_us()}, \texttt{time\_ms()} -- возвращают текущее системное время (счётчик тактов с момента запуска).
    \item \texttt{begin()} -- инициализирует окружение (\texttt{env::begin()}).
    \item \texttt{finish()} -- завершает окружение.
    \item \texttt{start(time\_us\_t \_period\_us = 0)} -- запускает системный таймер с заданным периодом. Если период не указан (0), пытается загрузить его из INI. Вызывает \texttt{env::start()}.
    \item \texttt{stop()} -- останавливает таймер (\texttt{env::stop()}).
    \item \texttt{backend\_loop()} -- вызывается в каждом такте backend. Обновляет системное время (\texttt{sysclock::tick()}) и вызывает \texttt{env::backend\_loop()}, а также \texttt{autonum::backend\_clean()}.
    \item \texttt{frontend\_loop()} -- вызывается в цикле frontend. Выводит накопленные сообщения из кольцевого буфера (если есть) и вызывает \texttt{env::frontend\_loop()} и \texttt{autonum::frontend\_clean()}.
    \item \texttt{startup()}, \texttt{shutdown()} -- если включена поддержка модулей, делегируют в \texttt{env}.
\end{itemize}

\subsubsection{Методы печати и форматирования}

\begin{itemize}
    \item \texttt{printf(cstr \_format, ...)} -- форматирует строку и выводит её. В зависимости от контекста:
        \begin{itemize}
            \item Если вызван из frontend -- сразу печатает через \texttt{env::print()}.
            \item Если из backend -- помещает символы в кольцевой буфер \texttt{print\_buffer\_} для последующего вывода в frontend.
        \end{itemize}
    \item \texttt{sprintf()} -- форматирует строку в предоставленный буфер. Использует либо \texttt{env::sprintf()} (если определён), либо стандартный \texttt{vsnprintf} / \texttt{vswprintf}.
\end{itemize}

\subsubsection{Методы генерации случайных чисел}

Шаблонные методы \texttt{rand} позволяют получать случайные числа в заданном диапазоне, используя либо \texttt{env::rand()}, либо стандартный \texttt{std::rand()}.

\subsection{Идеи, заложенные в дизайн}

\begin{enumerate}
    \item \textbf{Разделение контекстов frontend/backend} -- основа для построения систем жёсткого реального времени. Backend работает с минимальными задержками и без блокировок, а весь "тяжёлый" код выносится во frontend.
    \item \textbf{Автоматическое определение контекста в \texttt{guard}} -- программисту не нужно думать, из какого потока вызывается защищённый код; синхронизация выбирается оптимальная.
    \item \textbf{Слабая связанность с платформой} -- через класс \texttt{env} и слабые символы (\texttt{ROBO\_WEAK}) можно легко портировать фреймворк на новую платформу, переопределив только необходимые функции.
    \item \textbf{Безопасность при форматировании} -- использование отдельных буферов для frontend и backend исключает конфликты при параллельном вызове \texttt{printf}.
    \item \textbf{Встроенные средства отладки} -- статистика памяти, счётчики guard'ов, профилирование времени.
\end{enumerate}

\subsection{Что можно доработать}

\begin{itemize}
    \item \textbf{Улучшить обработку ошибок} -- сейчас многие функции просто ассертируются или возвращают false. Можно добавить коды ошибок или исключения (для frontend).
    \item \textbf{Сделать более гибкой настройку размеров буферов} -- например, размер кольцевого буфера для печати задаётся через \texttt{ROBO\_APP\_BACKEND\_PRINT\_BUFFER\_BITS}, но его можно было бы сделать настраиваемым через шаблон.
    \item \textbf{Добавить поддержку большего числа ОС} -- например, FreeRTOS, QNX.
    \item \textbf{Унифицировать работу с разделяемой памятью} -- сейчас в Linux используется \texttt{shm\_open}, но нет автоматической очистки; можно добавить RAII-обёртки.
    \item \textbf{Улучшить документацию} -- добавить примеры использования для каждой подсистемы.
    \item \textbf{Рефакторинг макросов} -- многие макросы дублируются в разных файлах; можно вынести все настройки в один заголовочный файл.
\end{itemize}

\subsection{Система логирования (\texttt{robosd\_log.hpp}, \texttt{robosd\_log.cpp})}

Файлы \texttt{robosd\_log.hpp} и \texttt{robosd\_log.cpp} реализуют гибкую систему логирования, позволяющую выводить сообщения с различными уровнями важности и фильтровать их по битовым маскам. Система тесно интегрирована с окружением (\texttt{system::env}) для платформозависимого вывода (например, с цветовым выделением в консоли) и поддерживает условную компиляцию: если макрос \texttt{ROBO\_APP\_DEBUG\_LOG\_ENABLED} установлен в 0, весь код логирования исключается, что позволяет полностью убрать накладные расходы в релизных сборках.

\subsubsection{Уровни логирования и маски}

В пространстве имён \texttt{robo::log} определены:

\begin{itemize}
    \item Перечисление \texttt{verb} -- уровни важности сообщения:
    \begin{lstlisting}
enum class verb { 
    skip = -3,   // skip
    error = -2,  // error
    warning = -1,// warning
    info = 0,    // info
    detail_1 = 1,// detail level 1
    detail_2 = 2,// detail level 2
    detail_3 = 3,// detail level 3
    detail_4 = 4,// detail level 4
    detail_5 = 5,// detail level 5
    detail_6 = 6,// detail level 6
    detail_7 = 7 // detail level 7
};
    \end{lstlisting}
    Чем выше значение, тем более детальная информация. Уровни \texttt{error}, \texttt{warning}, \texttt{info} предназначены для основных сообщений, а \texttt{detail\_N} -- для отладочной печати различной степени подробности.

    \item Структура \texttt{mask} -- битовая маска для категоризации сообщений:
    \begin{lstlisting}
struct mask {
    enum {
        disabled = 0, // output all (mask not used)
        app = 1,        // messages from application modules
        net = 2         // messages from network modules
    };
};
    \end{lstlisting}
    Пользователь может определять собственные маски (биты) для фильтрации вывода по категориям.
\end{itemize}

Глобальные переменные (скрытые в файле реализации):
\begin{itemize}
    \item \texttt{verb verb\_} -- текущий установленный уровень детализации.
    \item \texttt{unsigned int mask\_} -- текущая маска (биты, которые разрешены).
\end{itemize}

\subsubsection{Функции управления}

\begin{itemize}
    \item \texttt{bool begin(verb \_verb, unsigned int \_mask)} -- инициализирует систему логирования, устанавливая уровень и маску. Возвращает всегда \texttt{true}.
    \item \texttt{void finish(void)} -- завершает работу (в текущей реализации пустая).
    \item \texttt{void print(verb \_verb, unsigned int \_mask, cstr \_format, ...)} -- основная функция вывода. Проверяет, нужно ли печатать данное сообщение на основе текущего уровня и маски, и если да -- форматирует строку и передаёт её в \texttt{system::env::print()} вместе с уровнем для возможного цветового оформления. Реализация находится в \texttt{robosd\_log.cpp}.
\end{itemize}

Условие печати:
\begin{lstlisting}
if ( (_verb < verb::info) || ((verb_ >= _verb) && ((mask_ & _mask) == _mask)) )
\end{lstlisting}
Это означает, что сообщения уровня \texttt{error} и \texttt{warning} (которые меньше \texttt{info}) печатаются всегда, а для остальных проверяется, что текущий уровень не ниже уровня сообщения и все биты маски сообщения присутствуют в текущей маске.

\subsubsection{Макросы для логирования}

Для удобства использования определено множество макросов, которые скрывают вызов \texttt{log::print} и автоматически добавляют контекстную информацию (имя функции, файл, строку) для ошибок.

Основные макросы:

\begin{description}
    \item[\texttt{robo\_errlog(f, ...)}] -- выводит сообщение уровня \texttt{error} с добавлением имени функции, файла и номера строки. В Unicode-версии использует \texttt{\%S} для широких строк.
    \item[\texttt{robo\_warninglog(f, ...)}] -- уровень \texttt{warning}.
    \item[\texttt{robo\_infolog(f, ...)}] -- уровень \texttt{info}.
    \item[\texttt{robo\_detaillog(lvl, mask, f, ...)}] -- уровень \texttt{detail\_lvl} с указанной маской. Раскрывается в вызов \texttt{robo\_detaillog\_}, который подставляет нужное имя уровня.
    \item[\texttt{robo\_applog(f, ...)}] -- макрос для сообщений приложения (уровень \texttt{detail\_4}, маска \texttt{app}).
    \item[\texttt{robo\_netlog(f, ...)}] -- для сетевых сообщений (уровень \texttt{detail\_4}, маска \texttt{net}).
\end{description}

Все эти макросы используют механизм \texttt{VA\_OPT\_SUPPORTED} для правильной обработки пустого списка аргументов (добавление запятой только если есть аргументы). Если \texttt{ROBO\_APP\_DEBUG\_LOG\_ENABLED == 0}, макросы раскрываются в пустые выражения, и код логирования полностью исключается.

\subsubsection{Макросы для проверок с логированием}

Особую ценность представляют макросы, которые объединяют проверку условия и вывод сообщения об ошибке с последующим возвратом или переходом. Они значительно сокращают код обработки ошибок.

Категории макросов:

\begin{itemize}
    \item \textbf{Безусловные} -- всегда выводят сообщение и выполняют действие:
        \begin{itemize}
            \item \texttt{ROBO\_ALARM()} -- выводит сообщение "nested error".
            \item \texttt{ROBO\_ALARM\_F(f, ...)} -- выводит форматированное сообщение.
            \item \texttt{ROBO\_BREAK(fault)} -- выводит сообщение и возвращает \texttt{fault}.
            \item \texttt{ROBO\_BREAK\_F(fault, f, ...)} -- выводит форматированное сообщение и возвращает \texttt{fault}.
            \item \texttt{ROBO\_JAMP\_F(lbl, f, ...)} -- выводит сообщение и переходит на метку \texttt{lbl}.
        \end{itemize}
    \item \textbf{С проверкой условия} -- выполняют действие только если условие ложно:
        \begin{itemize}
            \item \texttt{ROBO\_ALARMN(x)} -- если \texttt{x} ложно, выводит сообщение.
            \item \texttt{ROBO\_ALARMN\_F(x, f, ...)} -- аналогично с форматированием.
            \item \texttt{ROBO\_BREAKN(x, fault)} -- если \texttt{x} ложно, выводит сообщение и возвращает \texttt{fault}.
            \item \texttt{ROBO\_BREAKN\_F(x, fault, f, ...)} -- с форматированием.
            \item \texttt{ROBO\_JAMPN(x, lbl)} -- если \texttt{x} ложно, выводит сообщение и переходит на \texttt{lbl}.
            \item \texttt{ROBO\_JAMPN\_F(x, lbl, f, ...)} -- с форматированием.
            \item \texttt{ROBO\_RET(x, success, fault)} -- если \texttt{x} ложно, выводит сообщение и возвращает \texttt{fault}, иначе \texttt{success}.
            \item \texttt{ROBO\_RET\_F(x, success, fault, f, ...)} -- с форматированием.
            \item \texttt{ROBO\_ASSERT(x)} -- если \texttt{x} ложно, выводит сообщение и вызывает \texttt{ROBO\_APP\_CRASH()}.
            \item \texttt{ROBO\_ASSERT\_F(x, f, ...)} -- с форматированием.
        \end{itemize}
\end{itemize}

Для удобства определены сокращения с префиксами \texttt{L} (возврат \texttt{false}) и \texttt{V} (возврат \texttt{ROBO\_EMPTY\_PARAM}, который используется для обозначения пустого параметра в шаблонах). Например:
\begin{itemize}
    \item \texttt{ROBO\_LBREAK()} -- \texttt{ROBO\_BREAK(false)}.
    \item \texttt{ROBO\_LBREAKN(x)} -- \texttt{ROBO\_BREAKN(x, false)}.
    \item \texttt{ROBO\_LRET(x)} -- \texttt{ROBO\_RET(x, true, false)}.
    \item \texttt{ROBO\_LRET\_F(x, f, ...)} -- \texttt{ROBO\_RET\_F(x, true, false, f, ...)}.
\end{itemize}

Эти макросы значительно упрощают написание защищённого кода, особенно при работе с ресурсами и проверкой возвращаемых значений.

\subsubsection{Реализация в \texttt{robosd\_log.cpp}}

Файл реализации содержит:

\begin{itemize}
    \item Определение глобальных переменных \texttt{verb\_} (по умолчанию \texttt{verb::info}) и \texttt{mask\_} (по умолчанию 0).
    \item Функцию \texttt{print}, которая:
        \begin{enumerate}
            \item Проверяет условие печати (описано выше).
            \item Если условие истинно, инициализирует \texttt{va\_list}, извлекает аргументы и вызывает \texttt{system::env::print(\_verb, \_format, args)}.
            \item Освобождает \texttt{va\_list}.
        \end{enumerate}
    \item Функции \texttt{begin} и \texttt{finish} -- тривиальные.
\end{itemize}

Важно отметить, что фактический вывод делегируется платформозависимому методу \texttt{system::env::print}, который может, например, раскрашивать сообщения в консоли в зависимости от уровня (как это сделано в реализациях для Windows и Linux). Таким образом, система логирования остаётся переносимой.

\subsubsection{Пример использования}

\begin{lstlisting}
#include "core/robosd_log.hpp"

void example() {
    robo::log::begin(robo::log::verb::detail_3, robo::log::mask::app | robo::log::mask::net);

    robo_infolog("Application started");
    robo_detaillog(2, robo::log::mask::app, "Initializing module %s", "motor");

    int result = some_function();
    ROBO_BREAKN_F(result == 0, "some_function failed with code %d", result);

    // ...
    robo::log::finish();
}
\end{lstlisting}

\subsubsection{Идеи и особенности}

\begin{itemize}
    \item \textbf{Условная компиляция} -- позволяет полностью исключить код логирования из финального бинарного файла, экономя память и процессорное время.
    \item \textbf{Централизованная фильтрация} -- через установку уровня и маски можно гибко управлять объёмом вывода без изменения исходного кода.
    \item \textbf{Интеграция с системой окружения} -- вывод может быть перенаправлен в любое устройство (консоль, файл, сеть) путём переопределения \texttt{system::env::print}.
    \item \textbf{Богатый набор макросов} -- сокращает объём кода обработки ошибок и делает его более читаемым.
    \item \textbf{Поддержка Unicode} -- через макросы \texttt{RT} и условную компиляцию.
\end{itemize}

\subsubsection{Возможные улучшения}

\begin{itemize}
    \item Добавить поддержка асинхронного логирования (например, через отдельный поток и очередь), чтобы не блокировать выполнение при медленном выводе.
    \item Расширить набор предопределённых масок.
    \item Добавить возможность перенаправления вывода в несколько целей одновременно (телеграм, файл, консоль).
    \item Реализовать автоматическое добавление временных меток к сообщениям.
\end{itemize}

\subsection{Утилиты для работы с INI-файлами (\texttt{robosd\_ini.hpp})}

Файл \texttt{robosd\_ini.hpp} предоставляет набор шаблонных функций, упрощающих загрузку параметров из INI-файлов и их преобразование в различные типы данных. Эти функции являются надстройкой над классом \texttt{string} и используют его методы \texttt{tryload} и \texttt{to\_number}, \texttt{to\_number\_array}, \texttt{to\_number\_list}. Все функции находятся в пространстве имён \texttt{robo::ini}.

Основная цель — минимизировать повторяющийся код при чтении конфигурации и автоматически обрабатывать ошибки с помощью макросов логирования (например, \texttt{ROBO\_LRET\_F}).

\subsubsection{Загрузка одиночных значений}

\begin{itemize}
    \item \texttt{template<typename T> bool try\_load(cstr \_sect, cstr \_key, T\& \_value)} — пытается загрузить значение параметра \texttt{\_key} из секции \texttt{\_sect} и преобразовать его в тип \texttt{T}. Возвращает \texttt{true} при успехе, \texttt{false} в противном случае. Использует промежуточный объект \texttt{string} для чтения строки и метод \texttt{to\_number} для преобразования.
    
    \item \texttt{template<typename T> bool try\_load(cstr \_sect\_first, cstr \_sect\_second, cstr \_key, T\& \_value)} — аналогично, но пробует две секции последовательно. Возвращает \texttt{true}, если значение найдено хотя бы в одной из них.
    
    \item \texttt{template<typename T> bool load(cstr \_sect, cstr \_key, T\& \_value)} — то же, что и \texttt{try\_load}, но в случае неудачи генерирует сообщение об ошибке через \texttt{ROBO\_LRET\_F} и возвращает \texttt{false}. Подходит для обязательных параметров.
    
    \item \texttt{template<typename T> bool load(cstr \_sect\_first, cstr \_sect\_second, cstr \_key, T\& \_value)} — загружает из первой секции, если не получается — из второй, иначе логирует ошибку.
\end{itemize}

\subsubsection{Загрузка массивов фиксированной длины}

\begin{itemize}
    \item \texttt{template<typename T> bool try\_load\_arr(cstr \_section, cstr \_key, T* \_values, size\_t \_count)} — загружает строку из INI, которая должна содержать ровно \texttt{\_count} чисел, разделённых пробелами или другими пробельными символами, и заполняет массив \texttt{\_values}. Возвращает \texttt{true} при успехе.
    
    \item \texttt{template<typename T> bool load\_arr(...)} — аналогично, но с логированием ошибки.
\end{itemize}

\subsubsection{Загрузка списков переменной длины}

\begin{itemize}
    \item \texttt{template<typename T> bool try\_load\_list(cstr \_section, cstr \_key, size\_t \_max\_count, T* \_values, size\_t\& \_count)} — загружает строку, содержащую список чисел (не более \texttt{\_max\_count}), и сохраняет их в массив \texttt{\_values}. Фактическое количество чисел возвращается в \texttt{\_count}. Возвращает \texttt{true}, если строка успешно прочитана и все числа корректны (даже если их меньше максимума).
    
    \item \texttt{template<typename T> bool load\_list(...)} — с логированием ошибки.
\end{itemize}

\subsubsection{Примеры использования}

Рассмотрим типичный INI-файл \texttt{config.ini}:

\begin{lstlisting}
[CONTROL]
gain = 1.5
kp = 2.0 0.5 0.1
limits = -10 10
\end{lstlisting}

Пример кода для загрузки параметров (комментарии на английском):

\begin{lstlisting}
#include "core/robosd_ini.hpp"

void load_config() {
    float gain;
    if (robo::ini::load("CONTROL", "gain", gain)) {
        // gain loaded successfully
    }

    double kp[3];
    if (robo::ini::load_arr("CONTROL", "kp", kp, 3)) {
        // coefficient array loaded
    }

    int16_t limits[2];
    size_t count;
    if (robo::ini::try_load_list("CONTROL", "limits", 2, limits, count)) {
        if (count == 2) {
            // two numbers loaded
        } else {
            // something is wrong, but no error logged
        }
    }
}
\end{lstlisting}

\subsubsection{Примечания}

\begin{itemize}
    \item Все функции полагаются на правильную работу класса \texttt{string} и его методов преобразования. Поддерживаются любые типы, для которых определена специализация \texttt{string::to\_number} (встроенные арифметические типы).
    \item Функции с префиксом \texttt{try\_} не генерируют сообщений в лог и могут использоваться для проверки наличия необязательных параметров.
    \item Макрос \texttt{ROBO\_LRET\_F} используется для логирования ошибок и возврата \texttt{false}. Он определён в \texttt{robosd\_log.hpp}.
    \item Для работы с широкими символами (Unicode) убедитесь, что \texttt{ROBO\_UNICODE\_ENABLED} установлен соответствующим образом.
\end{itemize}

\subsection{Класс \texttt{string} (\texttt{robosd\_string.hpp}, \texttt{robosd\_string.cpp})}

Класс \texttt{robo::string} предоставляет удобную обёртку над \texttt{std::basic\_string} с учётом особенностей фреймворка: поддержка Unicode (широких символов), безопасное форматирование в разных контекстах (frontend/backend), загрузка из INI-файлов, преобразование в числа и обратно. Это один из ключевых компонентов ядра, используемый повсеместно для работы со строками.

\subsubsection{Внутреннее устройство}

\begin{itemize}
    \item \texttt{value\_} — указатель на внутренний объект типа \texttt{base\_string\_}, который является наследником \texttt{std::basic\_string<char\_t>} (где \texttt{char\_t} зависит от \texttt{ROBO\_UNICODE\_ENABLED}: \texttt{char} или \texttt{wchar\_t}).
    \item Два статических буфера: \texttt{string\_buffer\_frontend} и \texttt{string\_buffer\_backend} размера \texttt{ROBO\_STRING\_BUFFER\_SIZE} (по умолчанию 255). Они используются для временного форматирования в соответствующем контексте без конфликтов.
    \item Вспомогательная структура \texttt{stream\_s} хранит указатель на заполненный буфер и размер.
\end{itemize}

\subsubsection{Конструкторы и операторы присваивания}

\begin{itemize}
    \item \texttt{string()} — создаёт пустую строку.
    \item \texttt{string(const string\&)} — копирование (глубокое).
    \item \texttt{string(cstr \_format, ...)} и \texttt{string(cstr \_format, va\_list)} — форматированное создание (доступно при \texttt{ROBO\_APP\_FORMATING\_TYPE != ROBO\_APP\_TYPE\_NONE}).
    \item \texttt{string(const T\&)} — шаблонный конструктор, использующий метод \texttt{from()} для преобразования арифметических типов в строку.
    \item Операторы присваивания от \texttt{string} и C-строки.
    \item Операторы сравнения (\texttt{==}, \texttt{!=}) как с \texttt{string}, так и с C-строкой.
\end{itemize}

\subsubsection{Загрузка из INI}

\begin{itemize}
    \item \texttt{bool load(cstr \_section, cstr \_key)} — загружает значение из указанной секции. В случае неудачи логирует ошибку.
    \item \texttt{bool load(cstr \_first, cstr \_second, cstr \_key)} — пробует две секции.
    \item \texttt{bool tryload(...)} — аналогично, но без логирования ошибок.
    \item \texttt{bool load(delegat::ref\& \_converter)} — загружает через пользовательский конвертер, который заполняет буфер.
\end{itemize}

При загрузке используется системная подсистема INI (\texttt{system::ini::load\_str}), а контекст (frontend/backend) определяет, какой буфер будет задействован.

\subsubsection{Преобразование в числа}

Класс предоставляет набор статических и нестатических методов для преобразования строк в числа и обратно.

\begin{itemize}
    \item \texttt{static bool to\_number(cstr begc, char\_t* \&endc, T\& \_value)} — базовый метод, использующий \texttt{wcstod} или \texttt{strtod} (с временной установкой локали "C"). Проверяет выход за пределы типа.
    \item \texttt{bool to\_number(T\& \_value)} — применяет \texttt{to\_number} к содержимому строки.
    \item \texttt{static bool scan\_numbers(cstr, size\_t max, T* values, size\_t\& count)} — последовательно извлекает числа из строки.
    \item \texttt{static bool to\_number\_array(cstr, T* values, size\_t count)} — требует ровно \texttt{count} чисел.
    \item \texttt{static bool to\_number\_list(...)} — извлекает до \texttt{max\_count} чисел, возвращает реальное количество.
\end{itemize}

Эти методы используются в утилитах \texttt{robosd\_ini.hpp} для загрузки массивов и списков.

\subsubsection{Конвертация в ASCII (UTF-8)}

Для вывода строк в окружения, не поддерживающие Unicode (например, старые консоли или файловые системы), предусмотрены методы:

\begin{itemize}
    \item \texttt{const char* ascii()} — возвращает указатель на временный буфер, содержащий строку в UTF-8 (если исходная строка в Unicode). Буфер выбирается в зависимости от контекста.
    \item \texttt{void ascii(char* buf, size\_t len) const} — копирует UTF-8 представление в предоставленный буфер.
    \item \texttt{void ascii(robo::lambda<void(const char*)> d) const} — передаёт UTF-8 строку в лямбду (с выделением временной памяти).
    \item \texttt{void asciib(robo::lambda<void(const uint8\_t*, size\_t)> d) const} — аналогично, но передаёт указатель на байты и длину.
\end{itemize}

\textbf{Важно:} В реализации \texttt{ascii} при \texttt{ROBO\_UNICODE\_ENABLED == 1} используется устаревший класс \texttt{std::wstring\_convert} и \texttt{std::codecvt\_utf8\_utf16}. Эти компоненты объявлены устаревшими в C++17 и могут быть удалены в C++26. Их использование может привести к предупреждениям или ошибкам компиляции в современных стандартах. Рекомендуется заменить на альтернативные методы конвертации, например, с помощью библиотеки \texttt{<charconv>} (для чисел) или специализированных библиотек (ICU, Boost.Locale). В будущем эта часть кода требует доработки.

\subsubsection{Форматирование}

При включённой опции \texttt{ROBO\_APP\_FORMATING\_TYPE} доступны методы:

\begin{itemize}
    \item \texttt{bool format(cstr \_format, va\_list \_args)} и \texttt{format(cstr \_format, ...)} — формируют строку по шаблону, используя контекстно-зависимые буферы (\texttt{sprintf\_backend\_} или \texttt{sprintf\_frontend\_}).
    \item \texttt{template<typename B> static bool format\_stream(B\& \_b, cstr \_format, va\_list \_args)} — форматирует и помещает результат в поток (объект с методом \texttt{put}), также учитывая контекст.
    \item Методы \texttt{from(...)} для арифметических типов — просто вызывают \texttt{format} с соответствующим спецификатором.
\end{itemize}

Внутренние функции \texttt{sprintf\_backend\_} и \texttt{sprintf\_frontend\_} используют \texttt{system::sprintf}, который в свою очередь может вызывать \texttt{vsnprintf} / \texttt{vswprintf} или платформозависимые реализации. Это стандартный и безопасный подход.

\subsubsection{Вспомогательные методы}

\begin{itemize}
    \item \texttt{size\_t length()} — длина строки.
    \item \texttt{cstr c\_str()} — доступ к внутренней C-строке.
    \item \texttt{void clear()} — очищает содержимое.
    \item \texttt{static size\_t strlen(cstr)} — обёртка над \texttt{strlen} / \texttt{wcslen}.
    \item \texttt{static uint8\_t utoa\_n(uint32\_t value, uint8\_t n, char\_t* r, char\_t space)} — быстрое безформатное преобразование целого числа в строку с фиксированной шириной (используется при \texttt{ROBO\_STRING\_UTOA\_ENABLED}).
\end{itemize}

\subsubsection{Примеры использования}

\begin{lstlisting}[caption={Демонстрация возможностей string}]
#include "core/robosd_string.hpp"

void example() {
    robo::string s1("Hello");
    robo::string s2(RT("World"));  // Unicode-aware

    s1.format(RT("%s %d"), s2.c_str(), 42);
    robo_infolog("s1 = %s", s1.ascii());  // console output

    double val;
    size_t count;
    if (s1.scan_numbers(1, &val, count)) {
        // extract number from string
    }

    // load from INI
    robo::string path;
    if (path.load("PATHS", "config")) {
        // path contains the string from config
    }
}
\end{lstlisting}

\subsubsection{Устаревшие и проблемные места}

\begin{enumerate}
    \item \textbf{Конвертация Unicode в UTF-8} (\texttt{ascii}): использование \texttt{std::wstring\_convert} и \texttt{std::codecvt\_utf8\_utf16}. Эти компоненты устарели в C++17 и будут удалены в C++26. Компиляция с флагом \texttt{-std=c++17} или выше может выдавать предупреждения \texttt{‘std::wstring\_convert’ is deprecated}. Рекомендуется заменить на библиотеку \texttt{<charconv>} (только для чисел), либо использовать внешние зависимости (ICU, Boost.Locale), либо написать собственную простую конвертацию для ограниченного набора символов. В контексте встраиваемых систем, где Unicode может не требоваться, можно отключить \texttt{ROBO\_UNICODE\_ENABLED}.

    \item \textbf{Использование \texttt{setlocale} в \texttt{to\_number}}: временная смена локали влияет на весь процесс и небезопасна в многопоточной среде. Вместо этого следует использовать функции, не зависящие от локали, например, \texttt{std::from\_chars} (C++17) для чисел. Однако текущий код рассчитан на поддержку старых компиляторов, поэтому остаётся как есть, но это потенциальная проблема.

    \item \textbf{TODO в коде}: "todo а ну как округлится?" указывает на неопределённое поведение при приведении double к целому типу. Хотя проверка диапазона выполняется, само приведение может округлить значение не так, как ожидается. Лучше использовать явное округление (например, \texttt{std::llround}) и проверять потерю точности.

    \item \textbf{Глобальные буферы}: \texttt{string\_buffer\_frontend} и \texttt{string\_buffer\_backend} имеют фиксированный размер. При форматировании длинных строк возможен обрезка. В коде нет проверки на полное заполнение (используется \texttt{system::sprintf}, который возвращает необходимое количество символов, но записывает только до предела). Однако риск потери данных существует.

    \item \textbf{Скрытое наследование от \texttt{std::basic\_string}}: класс \texttt{base\_string\_} публично наследует \texttt{stds} (typedef). Это не является ошибкой, но такое наследование от стандартного контейнера без виртуального деструктора может быть опасным, если использовать указатель на базовый класс. В данном случае деструктор \texttt{string} виртуален, а \texttt{base\_string\_} не имеет виртуальных методов, и удаление через указатель на \texttt{stds} не происходит. Однако потенциально это может привести к проблемам, если кто-то попытается использовать \texttt{base\_string\_} полиморфно. Лучше было бы использовать композицию (хранить \texttt{stds} как член), но текущий подход работает, пока нет внешнего вмешательства.
\end{enumerate}

\subsubsection{Макросы конфигурации, влияющие на \texttt{string}}

\begin{itemize}
    \item \texttt{ROBO\_UNICODE\_ENABLED} — если 1, используется \texttt{wchar\_t} и макрос \texttt{RT}.
    \item \texttt{ROBO\_STRING\_BUFFER\_SIZE} — размер временных буферов.
    \item \texttt{ROBO\_STRING\_UTOA\_TYPE} — тип реализации быстрого преобразования целых в строку (NATIVE или NONE).
    \item \texttt{ROBO\_APP\_FORMATING\_TYPE} — включает/выключает методы форматирования.
    \item \texttt{ROBO\_APP\_ENV\_ENABLED}, \texttt{ROBO\_APP\_SYSTEM\_ENABLED} — влияют на доступ к контекстным буферам.
\end{itemize}

\subsection{Делегаты и автоматическое управление памятью (\texttt{robosd\_delegat.hpp}, \texttt{robosd\_autonum.hpp}, \texttt{robosd\_autonum.cpp})}

Фреймворк \texttt{robosd} предоставляет две взаимосвязанные подсистемы для работы с вызываемыми объектами (делегатами):
\begin{itemize}
    \item Базовые классы делегатов (\texttt{robosd\_delegat.hpp}) — реализуют различные виды обёрток над функциями, методами и лямбдами.
    \item Механизм автоматического удаления (\texttt{robosd\_autonum.hpp/cpp}) — позволяет создавать делегаты, которые после использования не удаляются немедленно, а помещаются в специальные корзины (\texttt{receicledbin}) и очищаются в безопасные моменты (например, в циклах frontend/backend). Это позволяет избегать динамического выделения памяти в критических контекстах.
\end{itemize}

Обе подсистемы широко используются в ядре и модулях для реализации обратных вызовов, подписок на события и асинхронных операций.

\subsubsection{Базовые классы делегатов (\texttt{robosd\_delegat.hpp})}

Заголовочный файл определяет иерархию шаблонных классов, реализующих концепцию делегата — объекта, который можно вызвать как функцию. Все делегаты наследуются от абстрактного базового шаблона:

\begin{lstlisting}[caption={Базовый интерфейс делегата}]
template <typename R, typename ... Args> class ref {
protected:
    virtual ~ref(void) {}
public:
    virtual R operator ()(Args... args) = 0;
};
\end{lstlisting}

Производные классы реализуют оператор вызова, вызывая соответствующий целевой объект. Предусмотрены следующие типы делегатов (все они дополнительно параметризованы шаблонным параметром \texttt{B} — базовым классом, что позволяет комбинировать их с механизмами автоматического удаления):

\begin{itemize}
    \item \texttt{simple<B, R, Args...>} — делегат от обычной функции.
    \item \texttt{uni<B, R, Args...>} — делегат от функции, принимающей первым параметром \texttt{void*} (например, статический метод с передачей экземпляра).
    \item \texttt{lambda<B, R, Args...>} — делегат от объекта \texttt{robo::lambda} (см. \texttt{robosd\_lambda.hpp}).
    \item \texttt{member<B, C, R, Args...>} — делегат от метода класса, хранящий указатель на экземпляр.
    \item \texttt{rmember<B, C, R, Args...>} — делегат от метода класса, хранящий ссылку на экземпляр.
\end{itemize}

Для удобства создания делегатов без автоматического управления памятью определено пространство имён \texttt{owned::fabric}, которое предоставляет вложенные классы \texttt{simple}, \texttt{uni}, \texttt{lambda}, \texttt{member}, но с запретом на операторы \texttt{new} (чтобы объекты создавались только в статической памяти или на стеке). Однако эти классы редко используются напрямую; обычно применяются фабрики из \texttt{autonum}.

\subsubsection{Механизм автоматического удаления (\texttt{robosd\_autonum})}

Основная идея: объекты-делегаты, создаваемые динамически, после потери всех ссылок не удаляются сразу, а помещаются в одну из двух корзин (для frontend и backend). Периодически (в соответствующих циклах) эти корзины очищаются, и объекты действительно удаляются. Это позволяет избегать операций \texttt{delete} в контексте реального времени (backend) и снижает фрагментацию памяти.

\begin{itemize}
    \item \texttt{receicledbin::ref} — базовый класс для объектов, подлежащих автоудалению. Содержит:
        \begin{itemize}
            \item Счётчик использований \texttt{used\_}.
            \item Флаг \texttt{isfrontend\_}, определяющий контекст создания (frontend/backend).
            \item Ссылку на корзину (\texttt{rbref\_}) через интрузивный список.
            \item Методы \texttt{use()} и \texttt{unuse()}, увеличивающие/уменьшающие счётчик. При обнулении счётчика вызывается \texttt{release\_()}, которая перемещает объект в соответствующую корзину (а не удаляет сразу).
        \end{itemize}
    \item Статические методы \texttt{frontend\_clean()} и \texttt{backend\_clean()} очищают корзины, удаляя все находящиеся там объекты. Эти методы должны вызываться в соответствующих циклах (например, из \texttt{system::frontend\_loop} и \texttt{system::backend\_loop}).
    \item Для каждого контекста существует отдельный экземпляр структуры \texttt{core}, хранящий корзину (\texttt{rbin}) и ведущий счётчик созданных объектов (для отладки). Эти ядра инициализируются как статические объекты внутри функций \texttt{frontend\_core\_()} и \texttt{backend\_core\_()}.
\end{itemize}

\subsubsection{Фабрика \texttt{fabric} для создания автоудаляемых делегатов}

Шаблон \texttt{fabric<B, R, Args...>} объединяет возможности базовых делегатов и механизма автоудаления. Он создаёт классы-делегаты, которые одновременно наследуют от заданного базового класса \texttt{B} (который должен быть наследником \texttt{receicledbin::ref}) и от соответствующего типа делегата (например, \texttt{simple}, \texttt{uni} и т.д.). Внутри \texttt{fabric} определены вложенные классы:

\begin{itemize}
    \item \texttt{simple} — делегат от функции.
    \item \texttt{uni} — делегат от функции с \texttt{void*}.
    \item \texttt{lambda} — делегат от \texttt{robo::lambda}.
    \item \texttt{member<C>} — делегат от метода класса.
\end{itemize}

Каждый из них имеет статический метод \texttt{create}, который создаёт новый экземпляр делегата в куче (через \texttt{new}) и возвращает указатель на него. Сам делегат уже включает счётчик ссылок и логику автоудаления.

Кроме того, \texttt{fabric} предоставляет несколько перегрузок статического метода \texttt{create}, которые автоматически выбирают нужный тип делегата на основе переданных аргументов (простая функция, функция с \texttt{void*}, лямбда, метод класса). Это упрощает создание делегатов:

\begin{lstlisting}[caption={Примеры создания делегатов через autonum::fabric}]
// for a simple function
auto* d1 = robo::delegat::autonum_fabric<void(int)>::create(my_function);

// for a lambda
auto* d2 = robo::delegat::autonum_fabric<void(int)>::create([](int x) { ... });

// for a member function
SomeClass obj;
auto* d3 = robo::delegat::autonum_fabric<void(int)>::create(obj, &SomeClass::method);
\end{lstlisting}

Для удобства определён псевдоним \texttt{autonum\_fabric} и функция \texttt{ausimple}, которая создаёт простой делегат от функции.

\subsubsection{Пример использования}

Рассмотрим типичный сценарий: подписка на событие с использованием делегата, который будет автоматически удалён после отписки.

\begin{lstlisting}[caption={Использование autonum-делегата}]
#include "core/robosd_autonum.hpp"

class EventSource {
    using Callback = robo::delegat::autonum::performer<void, int>;
    robo::list::unsorted<Callback> callbacks_;
public:
    void subscribe(Callback* cb) {
        cb->attach();  // increases use count
        callbacks_.push(cb);
    }
    void unsubscribe(Callback* cb) {
        cb->dettach(); // decreases use count, may move to recycle bin
    }
    void fire(int value) {
        for (auto* cb = callbacks_.first(); cb; cb = cb->next()) {
            (*cb)(value);
        }
    }
};

void handler(int x) {
    robo_infolog("Handler called with %d", x);
}

int main() {
    EventSource src;
    auto* cb = robo::delegat::ausimple(handler);
    src.subscribe(cb);
    src.fire(42);
    src.unsubscribe(cb);
    // cb is now in recycle bin; will be deleted later during clean
    robo::delegat::autonum::frontend_clean(); // actually deletes cb
}
\end{lstlisting}

В этом примере делегат \texttt{cb} после вызова \texttt{dettach()} перемещается в корзину, но не удаляется немедленно, что безопасно для контекста, в котором могла происходить отписка (например, из обработчика события). Позже, при вызове \texttt{frontend\_clean()}, он будет окончательно удалён.

\subsubsection{Внутреннее устройство и важные детали}

\begin{itemize}
    \item Счётчик использований \texttt{used\_} управляется методами \texttt{use()} и \texttt{unuse()}. Обычно делегат имеет начальное значение 1 после создания. Метод \texttt{attach()} (определённый в производных классах) увеличивает счётчик, а \texttt{dettach()} — уменьшает. Когда счётчик достигает нуля, объект перемещается в корзину.
    \item Корзина реализована как интрузивный список (\texttt{rblist}), что позволяет хранить сами объекты без дополнительных выделений памяти.
    \item Очистка корзины (\texttt{clean\_}) просто извлекает элементы из списка и удаляет их через \texttt{delete}. Это делается под защитой \texttt{system::guard} для потокобезопасности.
    \item При создании делегата фиксируется контекст (\texttt{isfrontend\_}) с помощью \texttt{system::env::is\_frontend()}. Это гарантирует, что объект попадёт в правильную корзину.
\end{itemize}

\subsubsection{Устаревшие и проблемные места}

\begin{enumerate}
    \item \textbf{Орфографическая ошибка}: в названии класса \texttt{receicledbin} пропущена буква (должно быть \texttt{recycledbin}). Это мелочь, но может затруднять поиск.
    \item \textbf{Логирование в конструкторе/деструкторе} \texttt{ref} использует \texttt{robo\_infolog}, что может быть избыточно в релизных сборках и влиять на производительность. Рекомендуется обернуть в макрос отладки.
    \item \textbf{Потенциальная проблема с двойным удалением}: если объект уже находится в корзине, но кто-то ещё хранит указатель на него и вызывает \texttt{use()}/\texttt{unuse()}, это может привести к undefined behavior. В текущей реализации после помещения в корзину объект больше не доступен через обычные ссылки, но если программист сохранил сырой указатель и попытается его использовать, последствия непредсказуемы. Документация должна предупреждать об этом.
    \item \textbf{Запрет на копирование}: делегаты не копируются (нет конструктора копирования), но это неявно подразумевается. Однако в шаблонах \texttt{fabric} нет явного запрета, и при неправильном использовании можно случайно создать копию (если передать делегат по значению). Рекомендуется удалить конструкторы копирования и операторы присваивания в базовом классе \texttt{ref}.
    \item \textbf{Использование \texttt{system::guard} в критических местах}: в конструкторе и деструкторе \texttt{ref} используется \texttt{system::guard}, что может быть избыточно, так как счётчики ядер (\texttt{counter\_}) инкрементируются под тем же guard'ом. Возможно, стоит использовать более лёгкую синхронизацию или атомарные операции.
\end{enumerate}

\subsubsection{Макросы конфигурации}

\begin{itemize}
    \item \texttt{ROBO\_AUTONUM\_ENABLED} — включает/выключает весь механизм автоудаления. По умолчанию 1.
    \item \texttt{ROBO\_APP\_ENV\_ENABLED}, \texttt{ROBO\_APP\_SYSTEM\_ENABLED} — влияют на доступ к функциям определения контекста и синхронизации.
\end{itemize}

\subsection{Конвертер единиц измерения (\texttt{robosd\_convert.hpp}, \texttt{robosd\_convert.cpp})}

Модуль \texttt{robosd\_convert} предоставляет класс \texttt{converter}, предназначенный для линейного преобразования физических величин (например, напряжения, тока, угла) в целочисленные коды, используемые в цифровых системах (АЦП, ЦАП, регистры управления), и обратно. Класс является узлом приложения (\texttt{app::node}) и загружает свои параметры из INI-файла, что позволяет гибко настраивать масштабирование и ограничения для каждого канала.

\subsubsection{Назначение и область применения}

В робототехнических системах часто возникает необходимость пересчитывать физические величины (например, ток в амперах, напряжение в вольтах, угол в градусах) в целые числа, понятные микроконтроллеру или драйверу, и наоборот. Этот пересчёт обычно линеен:

\[
\text{code} = (\text{physical} - \text{offset}) \times \text{scale}
\]

\[
\text{physical} = \frac{\text{code}}{\text{scale}} + \text{offset}
\]

Кроме того, обычно задаются минимальное и максимальное физическое значение для ограничения (saturation). Класс \texttt{converter} инкапсулирует эти параметры и предоставляет удобные методы для преобразования в различные целочисленные типы (8, 16, 24, 32 бита со знаком и без знака) и обратно во float/double.

\subsubsection{Структура класса}

\begin{itemize}
    \item \textbf{Наследование}: \texttt{public app::node} — класс является узлом дерева приложения, что позволяет организовывать иерархию конвертеров, загружать параметры из INI и обращаться к ним по имени.
    \item \textbf{Приватные поля}:
        \begin{itemize}
            \item \texttt{float offset\_} — смещение (значение физической величины, соответствующее коду 0).
            \item \texttt{float scale\_} — масштабный коэффициент (количество единиц кода на единицу физической величины).
            \item \texttt{float min\_}, \texttt{float max\_} — минимальное и максимальное допустимое физическое значение (для насыщения).
            \item \texttt{float eps\_} — вычисляемая точность (половина шага квантования: \(0.5 / \text{scale}\)).
        \end{itemize}
    \item \textbf{Приватный метод}:
        \begin{itemize}
            \item \texttt{template<typename T> T sut\_(T value) const} — применяет насыщение (saturation) к значению, ограничивая его диапазоном \texttt{[min\_, max\_]}. Название, вероятно, сокращение от "saturate".
        \end{itemize}
\end{itemize}

\subsubsection{Загрузка параметров из INI}

При загрузке узла (\texttt{do\_load()}) из конфигурационного файла считываются следующие параметры (в текущей секции или в секции defaults):

\begin{itemize}
    \item \texttt{offset} — значение смещения.
    \item \texttt{scale} — масштаб.
    \item \texttt{min} — минимальная физическая величина.
    \item \texttt{max} — максимальная физическая величина.
\end{itemize}

После загрузки вычисляется \texttt{eps\_ = abs(0.5f / scale\_)} — половина шага квантования, которая может использоваться для оценки погрешности.

\subsubsection{Методы преобразования физическая величина → код}

Все методы преобразования принимают физическое значение типа \texttt{float} и возвращают целое число соответствующего типа. Общий алгоритм:

\begin{enumerate}
    \item Насыщение входного значения в диапазоне \texttt{[min\_, max\_]} (вызов \texttt{sut\_}).
    \item Вычисление промежуточного значения: \texttt{tmp = (value - offset\_) * scale\_}.
    \item Коррекция округления: если \texttt{tmp > 0}, прибавляется 0.5; если \texttt{tmp < 0}, вычитается 0.5 (для перехода к ближайшему целому).
    \item Ограничение результата диапазоном целевого типа (например, для int8: \texttt{[-127, 127]}, для uint8: \texttt{[0, 255]} и т.д.).
    \item Приведение к целевому типу.
\end{enumerate}

Предусмотрены методы:

\begin{itemize}
    \item \texttt{int8\_t to\_i8(float value) const}
    \item \texttt{int16\_t to\_i16(float value) const}
    \item \texttt{int32\_t to\_i24(float value) const} — для 24-битных знаковых целых (диапазон \texttt{-8388607 ... 8388607}).
    \item \texttt{int32\_t to\_i32(float value) const} (шаблонный, работает с любым типом, но фактически для float).
    \item \texttt{uint8\_t to\_u8(float value) const}
    \item \texttt{uint16\_t to\_u16(float value) const}
    \item \texttt{uint32\_t to\_u24(float value) const} — для 24-битных беззнаковых (до 16777215).
    \item \texttt{uint32\_t to\_u32(float value) const}
\end{itemize}

Обратите внимание: в методах для беззнаковых типов отрицательные промежуточные значения просто обнуляются, а не вызывают ошибку.

\subsubsection{Методы обратного преобразования код → физическая величина}

\begin{itemize}
    \item \texttt{float to\_float(int value) const}
    \item \texttt{double to\_double(int value) const}
\end{itemize}

Вычисление: \texttt{return value / scale\_ + offset\_} (если \texttt{scale\_} не равен нулю, иначе 0). Здесь нет насыщения, так как предполагается, что входной код уже лежит в допустимых пределах.

\subsubsection{Поиск конвертера по имени}

Статический метод \texttt{converter* find(cstr \_name)} использует \texttt{app::node::find} и динамическое приведение к \texttt{converter*}. Это позволяет получить доступ к конвертеру из любого места приложения, зная его иерархический путь.

\subsubsection{Пример использования}

Предположим, в INI-файле описаны два конвертера:

\begin{verbatim}
[SENSORS]
current.offset = 0.0
current.scale = 0.001   ; 1 мА на единицу АЦП
current.min = -10.0
current.max = 10.0

voltage.offset = 0.0
voltage.scale = 0.01    ; 10 мВ на единицу
voltage.min = 0.0
voltage.max = 50.0
\end{verbatim}

Код приложения:

\begin{lstlisting}[caption={Использование конвертера}]
#include "core/robosd_convert.hpp"

void process() {
    // find converter by path
    auto* current_conv = robo::converter::find("SENSORS.current");
    auto* voltage_conv = robo::converter::find("SENSORS.voltage");

    float phys_current = 2.5; // 2.5 A
    int16_t adc_code = current_conv->to_i16(phys_current);
    // adc_code = (2.5 - 0) / 0.001 = 2500, after rounding 2500

    uint16_t raw_voltage = 1234; // from ADC
    float phys_voltage = voltage_conv->to_float(raw_voltage);
    // phys_voltage = 1234 * 0.01 = 12.34 V
}
\end{lstlisting}

\subsubsection{Потенциальные проблемы и устаревшие места}

\begin{enumerate}
    \item \textbf{Массив \texttt{convert\_names} в .cpp-файле содержит русские строки}: 
\begin{lstlisting}
const char * convert_names[] = {
    "Смещение",
    "Разр.",
    "Мин.",
    "Макс.",
    0
};
\end{lstlisting}
    Этот массив нигде не используется в предоставленном коде. Вероятно, это остаток отладочной печати или заготовка для интерфейса. Его наличие не влияет на работу, но может сбивать с толку. Рекомендуется либо удалить, либо перевести на английский и использовать для логирования.

    \item \textbf{Использование \texttt{float} для всех вычислений}: в методах \texttt{to\_i24}, \texttt{to\_u24} и т.д. происходят операции с плавающей точкой, что может быть нежелательно в контексте реального времени (backend). Однако сам класс \texttt{converter} обычно используется во frontend при настройке или в медленных циклах. В критических местах лучше предварительно вычислить целочисленные коэффициенты.

    \item \textbf{Отсутствие проверки \texttt{scale\_ == 0}} в методах преобразования в код: если \texttt{scale\_} окажется нулевым, произойдёт деление на ноль в обратных методах, но в прямых методах умножение на ноль даст ноль, что может быть неожиданным. Желательно добавить проверку и обработку ошибки.

    \item \textbf{Округление}: используется "половина" прибавлением 0.5. Это стандартный способ округления к ближайшему целому, но он не соответствует IEEE-754 round-half-to-even. Для большинства инженерных задач это приемлемо, но стоит иметь в виду.

    \item \textbf{Ограничения для знаковых типов}: в методах \texttt{to\_i8}, \texttt{to\_i16} используются границы \texttt{-127 ... 127} и \texttt{-32767 ... 32767}, что исключает минимальное значение (\texttt{-128} и \texttt{-32768}). Это соответствует дополнительному коду с симметричным диапазоном, что типично для некоторых датчиков, но может не совпадать с ожиданиями. Возможно, это историческая особенность.

    \item \textbf{Шаблонный метод \texttt{to\_i32}} определён в заголовочном файле, но не используется для других типов. Он содержит похожую логику, но ограничивает результат диапазоном 32-битного знакового целого. Однако его реализация дублирует код из специализированных методов, что затрудняет поддержку. Лучше было бы сделать единый шаблон, параметризованный целевым типом и пределами.

    \item \textbf{Метод \texttt{sut\_}} (saturate) использует \texttt{float} сравнения и может быть медленным. Для целочисленных преобразований иногда выгоднее насыщать уже после масштабирования, но здесь это делается до умножения, что правильно для предотвращения переполнения при больших коэффициентах.
\end{enumerate}

\subsubsection{Макросы конфигурации}

\begin{itemize}
    \item \texttt{ROBO\_APP\_MODULE\_ENABLED} — если 1, класс \texttt{converter} включается в сборку. Поскольку он наследует \texttt{app::node}, требуется поддержка модулей.
\end{itemize}

\subsection{Функциональный объект \texttt{lambda} (\texttt{robosd\_lambda.hpp})}

Заголовочный файл \texttt{robosd\_lambda.hpp} предоставляет реализацию шаблонного класса \texttt{robo::lambda}, который представляет собой функциональный объект, способный хранить произвольные вызываемые объекты (функторы, лямбда-выражения) с захватом контекста. Этот класс является ключевым строительным блоком для системы делегатов (см. \texttt{robosd\_delegat.hpp}) и позволяет передавать лямбды и другие функторы как объекты с динамическим временем жизни, поддерживая копирование и присваивание.

\subsubsection{Назначение}

В C++ лямбда-выражения создают объекты неуказанного типа, которые нельзя напрямую хранить в контейнерах или передавать между функциями без шаблонов. Класс \texttt{robo::lambda} решает эту проблему, стирая тип (type erasure): он сохраняет копию переданного функтора в динамической памяти и предоставляет единый интерфейс вызова. Это позволяет, например, хранить список различных лямбд в контейнере или возвращать лямбду из функции.

\subsubsection{Внутреннее устройство}

Класс \texttt{lambda} является специализацией первичного шаблона \texttt{lambda<T>} для сигнатуры функции \texttt{Out(In...)}. Он содержит единственное поле — указатель на внутренний класс \texttt{context}. Каждый экземпляр \texttt{context} хранит:

\begin{itemize}
    \item Указатель на скопированные данные функтора (\texttt{p\_}) в виде сырых байтов.
    \item Размер скопированных данных (\texttt{sz\_}).
    \item Счётчик ссылок (\texttt{used\_}), позволяющий безопасно копировать объекты \texttt{lambda} без глубокого копирования функтора.
    \item Указатель на функцию-запускатель (\texttt{run\_}), которая принимает указатель на \texttt{context} и аргументы, приводит \texttt{p\_} обратно к исходному типу и вызывает оператор вызова функтора.
\end{itemize}

Статический метод \texttt{context::query<T>(const T\&)} создаёт новый контекст, копирует в него переданный объект \texttt{T} и устанавливает \texttt{run\_} как лямбду, выполняющую приведение и вызов. Счётчик ссылок инициализируется единицей.

Методы \texttt{use()} и \texttt{release(context*)} управляют счётчиком ссылок. При обнулении счётчика контекст удаляется.

\subsubsection{Публичный интерфейс класса \texttt{lambda}}

\begin{itemize}
    \item \texttt{lambda()} — создаёт пустой объект (без контекста).
    \item \texttt{lambda(const lambda\&)} — копирующий конструктор: увеличивает счётчик ссылок контекста источника.
    \item \texttt{template<typename T> lambda(const T\&)} — создаёт объект из произвольного функтора \texttt{T}. Вызывает \texttt{context::query(T)}.
    \item \texttt{operator=(const lambda\&)} — присваивание: освобождает старый контекст и присоединяется к новому с увеличением счётчика.
    \item \texttt{operator bool()} — проверяет, содержит ли объект действительный контекст.
    \item \texttt{Out operator()(In...)} — вызывает сохранённый функтор с переданными аргументами. Требует, чтобы контекст существовал.
\end{itemize}

\subsubsection{Пример использования}

\begin{lstlisting}[caption={Создание и использование lambda}]
#include "core/robosd_lambda.hpp"
#include <iostream>

void example() {
    // Store a lambda with capture
    int factor = 2;
    robo::lambda<int(int)> doubler = [factor](int x) { return x * factor; };

    // Call it
    int result = doubler(5); // result = 10

    // Copy lambda (shares the same context)
    auto copy = doubler;

    // Store different lambdas in a container
    robo::lambda<void()> tasks[2];
    tasks[0] = []() { std::cout << "Task 1\n"; };
    tasks[1] = []() { std::cout << "Task 2\n"; };

    for (auto& task : tasks) {
        if (task) task(); // check if not empty
    }
}
\end{lstlisting}

\subsubsection{Потенциальные проблемы и ограничения}

\begin{enumerate}
    \item \textbf{Динамическая аллокация}: каждый вызов конструктора от функтора выделяет память для копии объекта. В контексте реального времени (backend) это может быть недопустимо. Рекомендуется создавать объекты \texttt{lambda} только во frontend или использовать пулы памяти.
    \item \textbf{Исключения}: код не использует исключения, но при нехватке памяти \texttt{new} может бросить исключение (если не переопределён глобальный оператор). В проекте \texttt{robosd} оператор \texttt{new} переопределён и использует \texttt{system::mem::alloc}, который при ошибке вызывает ассерт. Это приемлемо для встраиваемых систем.
    \item \textbf{Производительность вызова}: вызов оператора \texttt{()} включает косвенность через указатель на функцию \texttt{run\_}. Это добавляет небольшие накладные расходы по сравнению с прямым вызовом, но обычно приемлемо.
    \item \textbf{Поддержка move-семантики}: класс не имеет конструктора перемещения, что может привести к лишним копированиям. Однако благодаря счётчику ссылок копирование дешёвое (только увеличение счётчика). Всё же для полноты можно добавить перемещение.
    \item \textbf{Потокобезопасность}: класс не является потокобезопасным, так как счётчик ссылок изменяется без синхронизации. Если объекты \texttt{lambda} используются в нескольких потоках, необходима внешняя синхронизация (например, через \texttt{system::guard}).
\end{enumerate}

\subsubsection{Связь с другими компонентами}

Класс \texttt{lambda} используется в \texttt{robosd\_delegat.hpp} как основа для делегата типа \texttt{lambda}. Он также может применяться самостоятельно для передачи лямбд в асинхронные обработчики или для хранения колбэков в модулях.

\subsection{Расчёт контрольных сумм (CRC) (\texttt{robosd\_crc.hpp}, \texttt{robosd\_crc.cpp})}

Модуль \texttt{robosd\_crc} предоставляет набор функций для вычисления циклических избыточных кодов (CRC), используемых для контроля целостности данных при передаче или хранении. В состав входят реализации нескольких популярных алгоритмов: CRC-8 ETSI EN 302 307, табличный CRC-8 с полиномом 0x31, табличный CRC-7 и табличный CRC-16 Modbus. Каждая функция включается в сборку отдельным макросом, что позволяет гибко управлять размером кода.

\subsubsection{Конфигурационные макросы}

\begin{itemize}
    \item \texttt{ROBO\_APP\_CRC8\_ETSI\_EN\_302\_307\_ENABLED} — если 1, включает функцию \texttt{crc8\_ETSIEN302307}.
    \item \texttt{ROBO\_APP\_CRC8\_BY\_TABLE\_ENABLED} — если 1, включает табличную реализацию CRC-8.
    \item \texttt{ROBO\_APP\_CRC7\_BY\_TABLE\_ENABLED} — если 1, включает табличную реализацию CRC-7.
    \item \texttt{ROBO\_APP\_CRC16\_MODBUS\_BY\_TABLE\_ENABLED} — если 1, включает табличную реализацию CRC-16 Modbus.
\end{itemize}

Все макросы по умолчанию равны 0 и должны быть явно определены пользователем.

\subsubsection{Функция \texttt{crc8\_ETSIEN302307}}

Вычисляет 8-битную контрольную сумму по алгоритму, описанному в стандарте ETSI EN 302 307 (используется в спутниковой связи, DVB-S2). Полином: 0xD5 (в виде многочлена: x8 + x7 + x6 + x4 + x2 + 1). Реализация без таблицы, битовыми операциями.

\begin{itemize}
    \item Прототип: \texttt{uint8\_t crc8\_ETSIEN302307(uint8\_t * \_pdata, size\_t \_size)}.
    \item Параметры: указатель на буфер данных, размер в байтах.
    \item Возвращает: 8-битную CRC.
    \item Алгоритм: начальное значение 0x00, для каждого байта выполняется XOR с текущим CRC, затем 8 циклов сдвига и XOR с полиномом при установленном старшем бите.
\end{itemize}

\textbf{Применение}: используется в системах спутниковой связи, но может быть полезен и в других задачах, где требуется быстрый бестабличный CRC.

\subsubsection{Табличная реализация CRC-8 (полином 0x31)}

Функция \texttt{crc8\_by\_table} реализует CRC-8 с полиномом 0x31 (x8 + x5 + x4 + 1), начальным значением 0xFF, без отражения и без финального XOR. Этот вариант часто используется в протоколах 1-Wire, Dallas/Maxim.

\begin{itemize}
    \item Прототип: \texttt{uint8\_t crc8\_by\_table(uint8\_t * \_pdata, size\_t \_len)}.
    \item Параметры: указатель на данные, длина.
    \item Возвращает: 8-битную CRC.
    \item Алгоритм: использует предвычисленную таблицу \texttt{crc8\_table} размером 256 байт. CRC инициализируется 0xFF, затем для каждого байта выполняется операция \texttt{crc = crc8\_table[crc XOR *\_pdata++]} (здесь XOR обозначает исключающее ИЛИ).
\end{itemize}

\textbf{Таблица CRC-8} (приведена в исходном коде) представляет собой стандартную таблицу для данного полинома. В комментариях к таблице указаны параметры: полином, инициализация, отсутствие отражения и XOR, а также проверочная сумма для строки "123456789" (0xF7). При включении функции эти комментарии остаются в коде, но не влияют на работу.

\subsubsection{Табличная реализация CRC-7}

Функция \texttt{crc7\_by\_table} вычисляет 7-битную CRC (результат помещается в 8-битную переменную, старший бит всегда 0). Используется, например, в протоколах SD/MMC. Параметры: полином 0x31 (аналогично CRC-8, но 7-битный), начальное значение 0xFF.

\begin{itemize}
    \item Прототип: \texttt{uint8\_t crc7\_by\_table(uint8\_t * \_pdata, size\_t \_len)}.
    \item Алгоритм: используется таблица \texttt{crc7\_table} размером 128 байт (поскольку CRC 7-битная, индекс 7 бит). Для каждого байта выполняется \texttt{crc = crc7\_table[(crc XOR *\_pdata++) >> 1]}. Сдвиг вправо на 1 обусловлен тем, что таблица построена для 7-битного индекса, а реальный индекс получается из 8 бит, но старший бит отбрасывается.
\end{itemize}

\textbf{Особенность реализации}: в исходном коде таблица \texttt{crc7\_table} определена только для первых 128 значений, хотя в комментариях присутствует ещё 128 закомментированных строк. Это может быть ошибкой или недоделкой. Пользователю следует проверить полноту таблицы и, при необходимости, раскомментировать остальную часть. В текущем виде функция работает только с одной половиной таблицы, что может приводить к неверным результатам для некоторых комбинаций данных. Рекомендуется исправить таблицу или использовать полную версию.

\subsubsection{Табличная реализация CRC-16 Modbus}

Функция \texttt{crc16\_modbus\_by\_table} реализует CRC-16, используемую в протоколе Modbus. Полином: 0x8005 (x16 + x15 + x2 + 1), начальное значение 0xFFFF, с отражением байт и финальным XOR с 0x0000.

\begin{itemize}
    \item Прототип: \texttt{uint16\_t crc16\_modbus\_by\_table(const uint8\_t * \_pdata, size\_t \_len)}.
    \item Алгоритм: использует две таблицы \texttt{aucCRCHi} и \texttt{aucCRCLo} по 256 байт каждая, что является классической реализацией Modbus CRC. Переменные \texttt{ucCRCHi} и \texttt{ucCRCLo} инициализируются 0xFF. Для каждого байта вычисляется индекс, затем обновляются старший и младший байты CRC.
    \item Результат возвращается как 16-битное число со старшим байтом в старшей позиции.
\end{itemize}

Эта функция широко применяется в промышленных протоколах и проверена временем.

\subsubsection{Пример использования}

Предположим, необходимо проверить целостность пакета данных с помощью CRC-16 Modbus:

\begin{lstlisting}[caption={Вычисление CRC-16 Modbus}]
#include "core/robosd_crc.hpp"

// Example data buffer
uint8_t data[] = {0x01, 0x03, 0x00, 0x00, 0x00, 0x01};
size_t len = sizeof(data);

#if ROBO_APP_CRC16_MODBUS_BY_TABLE_ENABLED == 1
uint16_t crc = robo::crc16_modbus_by_table(data, len);
// crc can now be appended to the message (low byte first)
uint8_t crc_lo = crc & 0xFF;
uint8_t crc_hi = (crc >> 8) & 0xFF;
#endif
\end{lstlisting}

Для CRC-8:

\begin{lstlisting}[caption={Вычисление CRC-8}]
uint8_t data[] = "123456789";
uint8_t crc8 = robo::crc8_by_table(data, 9);
// Expected value: 0xF7
\end{lstlisting}

\subsubsection{Потенциальные проблемы и устаревшие места}

\begin{enumerate}
    \item \textbf{Неполная таблица для CRC-7} — как отмечено выше, таблица \texttt{crc7\_table} содержит только 128 элементов, хотя в закомментированной части есть ещё 128. Это может привести к неверным расчётам. Пользователям следует проверить и при необходимости дополнить таблицу.

    \item \textbf{Отсутствие проверки на нулевой указатель} — функции не проверяют, что \texttt{\_pdata} не равен nullptr. Это может вызвать аварийное завершение, если передать некорректный указатель. Желательно добавить проверки (по крайней мере, в отладочных сборках).

    \item \textbf{Использование макросов для включения} — такой подход позволяет экономить память, но требует аккуратного управления макросами в проекте. Возможны конфликты, если макросы не определены.

    \item \textbf{Отсутствие C++-обёрток} — функции оформлены как обычные C-функции в пространстве имён \texttt{robo}. Для удобства можно было бы добавить шаблонные обёртки, работающие с контейнерами, но это не критично.
\end{enumerate}
\subsection{Механизм событий (\texttt{robosd\_event.hpp})}

Файл \texttt{robosd\_event.hpp} реализует гибкую систему событий, построенную на основе делегатов (см. \texttt{robosd\_delegat.hpp}) и механизма автоматического удаления (\texttt{robosd\_autonum.hpp}). Основной шаблонный класс \texttt{event\_t<R, Args...>} управляет списком подписчиков (\texttt{performer}), которые вызываются при наступлении события. Подписчики могут иметь приоритет, что определяет порядок их вызова, а также могут быть одноразовыми (удаляются после первого срабатывания). Система позволяет как получать результат от подписчиков (аккумулируя его), так и работать с событиями, не возвращающими значения.

\subsubsection{Класс \texttt{event\_t<R, Args...>}}

Шаблон \texttt{event\_t} параметризуется типом возвращаемого значения \texttt{R} и списком типов аргументов \texttt{Args...}, которые передаются подписчикам при вызове события.

Основные компоненты:

\begin{itemize}
    \item Вложенный класс \texttt{performer} — базовый класс для всех подписчиков. Он определяет интерфейс, который должны реализовать конкретные обработчики.
    \item Приватное поле \texttt{performers\_} — упорядоченный список подписчиков (тип \texttt{typename performer::list}), сортированный по приоритету. Используется интрузивный список из \texttt{robosd\_list.hpp}.
    \item Вспомогательный шаблон \texttt{result<T>} для аккумуляции результатов вызовов, когда \texttt{R} не равен \texttt{void}. Для \texttt{void} определена специализация, которая просто вызывает подписчики без накопления.
    \item Метод \texttt{raise(Args...)} — итерирует по списку подписчиков в порядке приоритета, вызывает каждого, передавая аргументы, и аккумулирует результат (если нужно). После вызова одноразовые подписчики удаляются из списка.
    \item Публичные typedef'ы \texttt{owned} и \texttt{autonum} — фабрики для создания подписчиков с автоматическим управлением памятью (см. \texttt{robosd\_delegat.hpp} и \texttt{robosd\_autonum.hpp}).
\end{itemize}

\subsubsection{Класс \texttt{performer}}

Абстрактный базовый класс для всех подписчиков. Содержит:

\begin{itemize}
    \item Перечисление \texttt{priority} — возможные приоритеты: \texttt{lo = -1}, \texttt{normal = 0}, \texttt{hi = 1}. Приоритет определяет порядок вызова: сначала более высокий приоритет (численно больше).
    \item Поле \texttt{ref\_} типа \texttt{typename list::ref} — интрузивная ссылка для включения в список \texttt{performers\_}. Ключом служит приоритет.
    \item Флаг \texttt{once\_} — если \texttt{true}, после первого вызова подписчик автоматически отписывается.
    \item Методы:
        \begin{itemize}
            \item \texttt{virtual void attach()} — вызывается после успешного добавления в список (может быть переопределён).
            \item \texttt{virtual void dettach()} — отписывает подписчика (удаляет из списка).
            \item \texttt{bool once()} и \texttt{void set\_once(bool)} — доступ к флагу одноразовости.
            \item \texttt{bool cancel()} — если подписчик ещё в списке, отписывает его и возвращает \texttt{false}; если уже отписан, возвращает \texttt{true}.
            \item \texttt{virtual R operator()(Args...)} — чисто виртуальный оператор вызова, который должен реализовать конкретный обработчик.
            \item \texttt{bool attach\_to(event\_t* \_event, priority \_priority = priority::lo)} — добавляет подписчика в событие с указанным приоритетом. Устанавливает ключ ссылки, вставляет в список и вызывает \texttt{attach()}.
        \end{itemize}
    \item Конструктор принимает флаг \texttt{once} (по умолчанию \texttt{false}) и инициализирует ссылку с приоритетом \texttt{normal} (временным, но потом при вызове \texttt{attach\_to} ключ меняется).
\end{itemize}

\subsubsection{Аккумуляция результатов}

Для обработки возвращаемых значений используется вспомогательный шаблон \texttt{result<T>}. Он определён внутри \texttt{event\_t} и имеет две специализации:

\begin{itemize}
    \item Для произвольного типа \texttt{T} (не void) содержит поле \texttt{t} типа \texttt{int} (странное решение, см. "Потенциальные проблемы") и методы:
        \begin{itemize}
            \item \texttt{void run(performer* p, Args... \_arg)} — вызывает подписчика и выполняет побитовое ИЛИ результата с \texttt{t} (приводя результат к \texttt{int}).
            \item \texttt{T value()} — возвращает \texttt{t}, приведённое к \texttt{T}.
        \end{itemize}
    \item Для \texttt{void} специализация не хранит результат, а просто вызывает подписчика.
\end{itemize}

Этот механизм позволяет, например, собирать флаги ошибок от нескольких подписчиков: если каждый возвращает целое число с битовыми флагами, результатом события будет объединение всех флагов.

\subsubsection{Типовые события}

В конце файла определены два часто используемых типа событий:

\begin{lstlisting}[caption={Predefined event types}]
namespace events {
    typedef event_t<void, const uint8_t*, size_t > on_receive;
    typedef event_t<void> on_panic;
}
\end{lstlisting}

\begin{itemize}
    \item \texttt{on\_receive} — event fired when data is received (pointer to buffer and size). Returns nothing.
    \item \texttt{on\_panic} — emergency event, no arguments.
\end{itemize}

\subsubsection{Пример использования}

Создадим простое событие, которое уведомляет подписчиков об изменении значения, и подпишем несколько обработчиков с разными приоритетами.

\begin{lstlisting}[caption={event\_t usage example}]
#include "core/robosd_event.hpp"
#include <iostream>

// Event that passes an integer and returns int (e.g., flags)
using MyEvent = robo::event_t<int, int>;

// Concrete subscriber
class MyHandler : public MyEvent::performer {
    int id_;
public:
    MyHandler(int id, bool once = false) : performer(once), id_(id) {}
    virtual int operator()(int value) override {
        std::cout << "Handler " << id_ << " got " << value << std::endl;
        return 1 << id_;  // returns a flag
    }
};

int main() {
    MyEvent ev;

    // Create subscribers with different priorities
    MyHandler h1(1, false);
    MyHandler h2(2, false);
    MyHandler h3(3, true); // one-shot

    h1.attach_to(&ev, MyEvent::performer::priority::hi);
    h2.attach_to(&ev, MyEvent::performer::priority::normal);
    h3.attach_to(&ev, MyEvent::performer::priority::lo);

    // Raise the event
    int result = ev.raise(42);
    std::cout << "Result flags: " << result << std::endl;

    // Second call: h3 is already unsubscribed
    result = ev.raise(100);
    std::cout << "Result flags: " << result << std::endl;

    return 0;
}
\end{lstlisting}

Expected output:
\begin{verbatim}
Handler 1 got 42
Handler 2 got 42
Handler 3 got 42
Result flags: 14   (bits 1,2,3: 2+4+8=14)
Handler 1 got 100
Handler 2 got 100
Result flags: 6    (bits 1,2: 2+4=6)
\end{verbatim}

\subsubsection{Потенциальные проблемы и устаревшие места}

\begin{enumerate}
       \item \textbf{Приведение результата к \texttt{int} в шаблоне \texttt{result<T>} для не-void типов.} Это серьёзное ограничение: предполагается, что все возвращаемые значения могут быть неявно преобразованы в \texttt{int} и что побитовое ИЛИ имеет смысл. Для типов, не являющихся целыми числами (например, указатели, структуры), такое приведение приведёт к ошибкам компиляции или неопределённому поведению. Более правильным было бы использовать операцию накопления, задаваемую пользователем, или требовать, чтобы тип \texttt{R} поддерживал оператор \texttt{|=}.

    \item \textbf{Порядок вызова и приоритеты.} Приоритет задаётся перечислением с тремя значениями, но технически можно использовать любые целые числа, так как \texttt{sorted}-список использует сравнение ключей. Однако в текущей реализации ключ устанавливается только через \texttt{priority}, что ограничивает гибкость.

    \item \textbf{Потокобезопасность.} Класс \texttt{event\_t} не содержит никакой синхронизации. Предполагается, что все операции (подписка, отписка, вызов) происходят в одном контексте (например, во frontend). Если событие может вызываться из разных потоков, необходима внешняя синхронизация.

    \item \textbf{Управление памятью.} В примере подписчики созданы на стеке, но они должны жить до тех пор, пока не будут отписаны. При использовании динамического создания через фабрики \texttt{owned} или \texttt{autonum} необходимо следить за временем жизни, чтобы избежать висячих указателей.

    \item \textbf{Орфография.} Метод \texttt{dettach()} написан с двумя 't' (правильно: \texttt{detach}). Это не ошибка компиляции, но может затруднить поиск.

    \item \textbf{Метод \texttt{attach()}} в \texttt{performer} по умолчанию пуст. Возможно, он задумывался для уведомления о том, что подписчик добавлен, но не используется. Его можно было бы сделать чисто виртуальным, если требуется обязательная реализация.
\end{enumerate}

\subsubsection{Конфигурационные макросы}

Для работы \texttt{event\_t} требуется включение \texttt{robosd\_delegat.hpp} и \texttt{robosd\_autonum.hpp}, которые, в свою очередь, зависят от макросов \texttt{ROBO\_AUTONUM\_ENABLED} и других. Сам \texttt{event\_t} не имеет собственных макросов включения, поэтому он доступен всегда, если подключены соответствующие заголовки.

\subsection{Интрузивные списки (\texttt{robosd\_list.hpp})}

Файл \texttt{robosd\_list.hpp} предоставляет реализацию различных типов интрузивных списков — структур данных, в которых ссылки на элементы хранятся непосредственно в самих объектах. Это позволяет избежать дополнительных выделений памяти и повышает эффективность в системах реального времени. Все списки находятся в пространстве имён \texttt{robo::list}.

Основные возможности:
\begin{itemize}
    \item Двунаправленные списки (\texttt{base}, \texttt{unsorted}, \texttt{sorted}, \texttt{unique}, \texttt{pool}).
    \item Однонаправленные списки (\texttt{unidir::base\_t}, \texttt{store\_t}, \texttt{queue\_t}, \texttt{stack\_t}).
    \item Очереди на основе списков (\texttt{queue::fifo}, \texttt{queue::priority}).
\end{itemize}

Все списки используют интрузивный подход: объект, который может быть помещён в список, должен содержать в себе объект-ссылку (например, \texttt{list::unsorted<T>::ref}), который хранит указатели на следующий и предыдущий элементы.

\subsubsection{Базовые классы}

\texttt{base\_ref<T>} — базовый класс для элемента списка. Содержит:
\begin{itemize}
    \item Указатели \texttt{next} и \texttt{prev} на соседние элементы.
    \item Ссылку на владельца \texttt{owner\_} (объект типа \texttt{T}).
    \item Указатель на список \texttt{list\_}, в котором находится элемент.
    \item Методы \texttt{attach\_to}, \texttt{attach\_before}, \texttt{attach\_after} для вставки в список.
    \item Метод \texttt{dettach()} для удаления из списка.
    \item Методы \texttt{next\_ptr()}, \texttt{prev\_ptr()} для доступа к соседним объектам.
\end{itemize}

\texttt{base<T>} — базовый класс двунаправленного списка. Управляет головой (\texttt{first}) и хвостом (\texttt{last}), поддерживает вставку до и после заданного элемента, удаление, подсчёт количества элементов. Методы:
\begin{itemize}
    \item \texttt{int count()} — количество элементов.
    \item \texttt{T* pop()} — извлекает первый элемент из списка.
    \item \texttt{void clean()} — удаляет все элементы.
    \item \texttt{base\_ref<T>* locate(const T*)} — поиск ссылки по объекту.
    \item \texttt{bool contains(const T\&)} — проверка наличия объекта.
\end{itemize}

\subsubsection{Неупорядоченный список \texttt{unsorted<T>}}

Предназначен для хранения элементов без определённого порядка. Добавление происходит в конец списка (если не указано иное). Предоставляет вложенный класс \texttt{ref}, который наследует \texttt{base\_ref<T>} и добавляет удобные методы для вставки:

\begin{lstlisting}[caption={Интерфейс unsorted::ref}]
class ref : public base_ref<T> {
public:
    ref(T& _owner);
    void attach_to(unsorted& _L);          // add to end
    void attach_before(unsorted& _L, ref* _prev); // insert before given ref
    void attach_after(unsorted& _L, ref* _next);  // insert after given ref
    void attach_before(unsorted& _L, T* _prev);   // insert before object
    void attach_after(unsorted& _L, T* _next);    // insert after object
    ref* prev();
    ref* next();
};
\end{lstlisting}

Пример использования:

\begin{lstlisting}[caption={Неупорядоченный список}]
#include "core/robosd_list.hpp"
#include <iostream>

struct MyData {
    int value;
    robo::list::unsorted<MyData>::ref list_ref;
    MyData(int v) : value(v), list_ref(*this) {}
};

int main() {
    robo::list::unsorted<MyData> list;
    MyData a(1), b(2), c(3);

    a.list_ref.attach_to(list);
    b.list_ref.attach_to(list);
    c.list_ref.attach_before(list, &b); // insert c before b

    for (auto* ref = list.first(); ref; ref = ref->next()) {
        std::cout << ref->owner().value << " ";
    }
    // Output: 1 3 2
}
\end{lstlisting}

\subsubsection{Упорядоченные списки с ключами}

Для сортированных списков используется шаблон \texttt{pair<T, K, L, unique>}, который добавляет ключ типа \texttt{K}. Классы \texttt{sorted<T,K>}, \texttt{unique<T,K>} и \texttt{pool<T,K>} используют этот базовый класс.

\begin{itemize}
    \item \texttt{sorted<T,K>} — список, отсортированный по ключу (по возрастанию). Поддерживает вставку с сохранением порядка. Имеет флаг \texttt{sort\_}, который всегда \texttt{true}. Метод \texttt{inc\_key(const K\&)} позволяет увеличить ключи всех элементов на заданную величину (использует \texttt{hack\_key\_}).
    \item \texttt{unique<T,K>} — список с уникальными ключами. При попытке вставить элемент с уже существующим ключом вставка не производится. Предоставляет метод \texttt{find(const K\&)} для поиска объекта по ключу.
    \item \texttt{pool<T,K>} — неупорядоченный список, но с ключами (наследует \texttt{sorted}, но \texttt{sort\_} установлен в \texttt{false}). Фактически работает как \texttt{unsorted} с ключами.
\end{itemize}

Пример использования \texttt{unique}:

\begin{lstlisting}[caption={Список с уникальными ключами}]
struct Item {
    int id;
    robo::list::unique<Item, int>::ref list_ref;
    Item(int i) : id(i), list_ref(*this, i) {}
};

robo::list::unique<Item, int> items;
Item a(1), b(2), c(1); // c has same key as a

a.list_ref.attach_to(items); // ok
b.list_ref.attach_to(items); // ok
c.list_ref.attach_to(items); // fails, returns false

Item* found = items.find(2); // returns &b
\end{lstlisting}

\subsubsection{Однонаправленные списки (\texttt{unidir})}

Пространство имён \texttt{unidir} содержит реализации однонаправленных списков, полезных для очередей и стеков.

\begin{itemize}
    \item \texttt{base\_t<T, I>} — базовый класс однонаправленного списка. Параметр \texttt{I} — тип элемента, который должен иметь поле \texttt{next\_}. Методы: \texttt{push(I\&)}, \texttt{pop\_first()}, \texttt{extract(const T\&)}.
    \item \texttt{store\_t<T, Arg...>} — хранилище, которое само создаёт элементы через \texttt{new} при вызове \texttt{push(Arg...)}. Возвращает ссылку на созданный объект. Имеет метод \texttt{free()} для удаления всех элементов.
    \item \texttt{queue\_t<T, Arg...>} — очередь (FIFO) на основе \texttt{base\_t}. Позволяет добавлять уже существующие элементы (\texttt{push(item*)}) и извлекать первый (\texttt{pop()}).
    \item \texttt{stack\_t<T, Arg...>} — стек (LIFO) (не полностью реализован, метод \texttt{pop\_last} отсутствует в \texttt{base\_t}, но заявлен).
\end{itemize}

Пример использования \texttt{store\_t}:

\begin{lstlisting}[caption={Хранилище объектов}]
struct Point {
    int x, y;
    Point(int a, int b) : x(a), y(b) {}
};

robo::list::unidir::store_t<Point, int, int> points;
Point& p1 = points.push(10, 20);
Point& p2 = points.push(30, 40);

// p1 и p2 живут в списке, будут удалены при уничтожении points
\end{lstlisting}

\subsubsection{Очереди на основе списков}

Пространство имён \texttt{queue} содержит адаптеры, превращающие списки в очереди. Шаблон \texttt{base<T, L>} берёт список типа \texttt{L} (например, \texttt{unsorted<T>}) и добавляет метод \texttt{push(T*)}. Предопределены две очереди:

\begin{itemize}
    \item \texttt{fifo<T>} — очередь на основе \texttt{unsorted<T>} (добавление в конец, извлечение из начала через \texttt{pop()}).
    \item \texttt{priority<T>} — очередь с приоритетом на основе \texttt{sorted<T, typename T::priority\_t>}. Требует, чтобы тип \texttt{T} имел вложенный тип \texttt{priority\_t} и соответствующий ключ.
\end{itemize}

\subsubsection{Потенциальные проблемы и замечания}

\begin{enumerate}
    \item \textbf{Интрузивность}: объекты, помещаемые в список, должны содержать в себе объект-ссылку соответствующего типа. Это может быть неудобно, если класс уже спроектирован иначе, но даёт максимальную производительность.
    \item \textbf{Ручное управление ссылками}: необходимо явно вызывать \texttt{attach\_to} и \texttt{dettach()}. При копировании объекта ссылка не копируется автоматически, что может привести к некорректному состоянию.
    \item \textbf{Потокобезопасность}: ни один из списков не синхронизирован. В многопоточной среде необходима внешняя синхронизация (например, через \texttt{system::guard}).
    \item \textbf{Орфография}: метод \texttt{dettach()} написан с двумя 't' (правильно \texttt{detach}), но это не влияет на компиляцию.
    \item \textbf{Неполная реализация стека}: в \texttt{stack\_t} используется метод \texttt{pop\_last()}, который не определён в \texttt{base\_t}. Вероятно, это недоработка; для стека нужен доступ к последнему элементу, что в однонаправленном списке требует обхода.
    \item \textbf{Использование \texttt{hack\_key\_}} для изменения ключа без перевставки может нарушить упорядоченность, если ключ меняется после вставки. Метод \texttt{inc\_key} в \texttt{sorted} использует это, но предполагает, что изменение ключа на одну и ту же величину для всех элементов сохраняет порядок. Это верно только для линейного преобразования.
\end{enumerate}

\subsubsection{Макросы конфигурации}

Сам файл не использует специальных макросов, кроме стандартных \texttt{ROBO\_EXPORT} и \texttt{ROBO\_APP\_ASSERT}, определённых в \texttt{robosd\_common.hpp}. Для корректной работы требуется включение \texttt{robosd\_common.hpp}.

\subsection{Кольцевые буферы (\texttt{robosd\_ring.hpp}, \texttt{robosd\_ring\_buf.hpp}, \texttt{robosd\_ring\_safe\_buf.hpp})}

В составе фреймворка \texttt{robosd} представлено несколько реализаций кольцевых буферов (FIFO), различающихся способом управления памятью (статический vs динамический) и уровнем потокобезопасности. Они широко используются для организации очередей данных между контекстами (frontend/backend), буферизации потокового ввода/вывода и временного хранения сообщений.

\subsubsection{Общие принципы}

Все реализации основаны на классическом кольцевом буфере с двумя счётчиками: количество записанных элементов (\texttt{writeCount}) и количество прочитанных (\texttt{readCount}). Размер буфера задаётся как степень двойки ($2^{\texttt{BITS}}$), что позволяет использовать битовую маску для вычисления индекса без операции деления. Это даёт высокую производительность и детерминированное поведение.

Основные операции:
\begin{itemize}
    \item \texttt{put(value)} — запись элемента.
    \item \texttt{get()} — чтение элемента.
    \item \texttt{available()}, \texttt{count()} — проверка наличия данных.
    \item \texttt{full()}, \texttt{space()} — проверка заполненности и свободного места.
    \item \texttt{clear()} — сброс буфера.
\end{itemize}

\subsubsection{Класс \texttt{ring} (\texttt{robosd\_ring.hpp})}

Представляет собой динамическую версию кольцевого буфера для байтов (\texttt{uint8\_t}). Память под буфер выделяется в куче при вызове \texttt{setup()} или в конструкторе.

\begin{itemize}
    \item Поля:
        \begin{itemize}
            \item \texttt{uint8\_t* data\_} — указатель на выделенный буфер.
            \item \texttt{unsigned int readCount\_}, \texttt{writeCount\_} — счётчики чтения и записи.
            \item \texttt{unsigned int mask\_}, \texttt{masksz\_} — маски для индексов.
            \item \texttt{uint8\_t bits\_} — логарифм размера.
            \item \texttt{unsigned int size\_} — фактический размер.
        \end{itemize}
    \item Методы:
        \begin{itemize}
            \item \texttt{bool setup(uint8\_t \_bits)} — выделяет буфер размером $2^{\_bits}$ байт. Возвращает \texttt{true} при успехе.
            \item \texttt{void cleanup()} — освобождает память.
            \item \texttt{void put(uint8\_t value)} — записывает один байт.
            \item \texttt{uint8\_t get()} — читает один байт.
            \item \texttt{bool put(uint8\_t* \_data, unsigned int \_count)} — записывает массив, если достаточно места.
            \item \texttt{unsigned int get(uint8\_t* \_data, unsigned int \_max\_count)} — читает до \texttt{\_max\_count} байт, возвращает реальное количество.
            \item \texttt{unsigned int count()}, \texttt{unsigned int space()}, \texttt{bool available()}, \texttt{bool full()}, \texttt{void clear()} — аналогично.
        \end{itemize}
\end{itemize}

\textbf{Примечание:} Класс не является потокобезопасным. Использование в многопоточной среде требует внешней синхронизации (например, \texttt{system::guard}).

\subsubsection{Шаблон \texttt{ring\_t<BITS, DATA\_T>} (\texttt{robosd\_ring\_buf.hpp})}

Статическая версия кольцевого буфера, где массив элементов встроен в объект. Размер фиксирован и задаётся шаблонным параметром \texttt{BITS} (количество бит). Тип элементов может быть любым (\texttt{DATA\_T} по умолчанию \texttt{uint8\_t}).

\begin{itemize}
    \item Поля:
        \begin{itemize}
            \item \texttt{DATA\_T \_data[SIZE]} — статический массив размера $2^{\texttt{BITS}}$.
            \item \texttt{size\_t \_readCount}, \texttt{\_writeCount} — счётчики.
            \item Статические константы \texttt{\_mask} и \texttt{\_masksz} для индексации.
        \end{itemize}
    \item Методы аналогичны \texttt{ring}, но оперируют типом \texttt{DATA\_T} и не требуют выделения памяти.
    \item Присутствует неиспользуемое поле \texttt{int test}, вероятно, остаток отладки.
\end{itemize}

Преимущество — отсутствие динамической аллокации, что делает его пригодным для использования в контексте реального времени (backend).

\subsubsection{Шаблон \texttt{struct\_ring\_t<BITS, STRUCT\_T>} (\texttt{robosd\_ring\_buf.hpp})}

Специализированная версия для работы со структурами. Вместо непосредственной записи/чтения предоставляет доступ к пустому элементу и текущему первому элементу, а затем отдельными вызовами \texttt{put()} и \texttt{get()} увеличивает счётчики. Это удобно, когда нужно заполнить структуру поэтапно, прежде чем считать её готовой.

\begin{itemize}
    \item Методы:
        \begin{itemize}
            \item \texttt{STRUCT\_T\& empty()} — возвращает ссылку на следующий свободный элемент (по индексу \texttt{\_writeCount}). После заполнения нужно вызвать \texttt{put()}.
            \item \texttt{STRUCT\_T\& first()} — возвращает ссылку на первый элемент в очереди (по индексу \texttt{\_readCount}). После обработки нужно вызвать \texttt{get()}.
            \item \texttt{void put()} — увеличивает счётчик записи, фиксируя, что элемент готов.
            \item \texttt{void get()} — увеличивает счётчик чтения, удаляя элемент из очереди.
            \item Остальные методы (\texttt{available()}, \texttt{count()}, \texttt{full()}, \texttt{clear()}) работают аналогично.
        \end{itemize}
\end{itemize}

Этот подход позволяет избежать копирования данных при передаче крупных объектов.

\subsubsection{Шаблон \texttt{safe\_ring\_t<G, BITS, DATA\_T>} (\texttt{robosd\_ring\_safe\_buf.hpp})}

Потокобезопасная версия кольцевого буфера. Шаблонный параметр \texttt{G} должен быть классом-охранником (guard), который обеспечивает синхронизацию (например, \texttt{system::guard}). Все публичные методы создают локальный объект \texttt{G g\_\_}, который блокирует доступ на время выполнения операции.

\begin{itemize}
    \item Внутренняя реализация аналогична \texttt{ring\_t}, но все операции с данными вынесены в приватные методы \texttt{put\_()}, \texttt{get\_()}, \texttt{available\_()} и т.д., которые не синхронизированы.
    \item Публичные методы (\texttt{put()}, \texttt{get()}, \texttt{count()}, \texttt{space()} и т.д.) создают объект \texttt{G} и вызывают соответствующие приватные методы.
    \item Благодаря этому буфер можно безопасно использовать из нескольких потоков (например, один писатель в backend, один читатель в frontend).
\end{itemize}

\textbf{Важно:} Класс предполагает, что \texttt{G} реализует блокировку в конструкторе и разблокировку в деструкторе (RAII). В фреймворке этому соответствует \texttt{system::guard}.

\subsubsection{Примеры использования}

\textbf{Пример 1:} Передача байтов из прерывания в основной цикл (с использованием \texttt{safe\_ring\_t}).

\begin{lstlisting}[caption={Thread-safe ring buffer}]
#include "core/robosd_ring_safe_buf.hpp"
#include "core/robosd_system.hpp"

using SafeRing = robo::safe_ring_t<robo::system::guard, 8>; // 256 bytes
SafeRing rx_buffer;

void UART_IRQHandler() {
    uint8_t byte = UART->DR;
    rx_buffer.try_put(byte); // atomic, protected by guard
}

void main_loop() {
    uint8_t data[16];
    size_t got = rx_buffer.get(data, 16);
    if (got) process(data, got);
}
\end{lstlisting}

\textbf{Пример 2:} Очередь сообщений с помощью \texttt{struct\_ring\_t}.

\begin{lstlisting}[caption={Message queue using struct\_ring\_t}]
struct Message {
    uint32_t id;
    float value;
};

robo::struct_ring_t<4, Message> msg_queue; // 16 messages

void produce() {
    if (!msg_queue.full()) {
        Message& m = msg_queue.empty();
        m.id = 1;
        m.value = 3.14;
        msg_queue.put(); // message is ready
    }
}

void consume() {
    if (msg_queue.available()) {
        Message& m = msg_queue.first();
        process(m);
        msg_queue.get(); // remove from queue
    }
}
\end{lstlisting}

\subsubsection{Потенциальные проблемы и замечания}

\begin{enumerate}
    \item \textbf{Неиспользуемая переменная \texttt{test}} в \texttt{ring\_t} и \texttt{safe\_ring\_t} — вероятно, артефакт отладки. Не влияет на функциональность, но может быть удалена.
    \item \textbf{Отсутствие проверки на переполнение} в методах \texttt{put(value)} (одиночная запись). В \texttt{ring\_t} и \texttt{safe\_ring\_t} нет защиты от записи в полный буфер — это приведёт к перезаписи старых данных. Разработчик должен сам проверять \texttt{space()} перед вызовом. В \texttt{try\_put} такая проверка есть.
    \item \textbf{Потокобезопасность \texttt{safe\_ring\_t}} зависит от правильной реализации guard'а. В фреймворке \texttt{system::guard} автоматически определяет контекст и выбирает подходящую синхронизацию, что делает этот буфер удобным для взаимодействия frontend/backend.
    \item \textbf{Динамическое выделение памяти} в классе \texttt{ring} может быть проблемой в системах с ограниченными ресурсами или в контексте реального времени. Предпочтение следует отдавать статическим версиям (\texttt{ring\_t}).
    \item \textbf{Неполная инкапсуляция}: в \texttt{ring\_t} и \texttt{safe\_ring\_t} счётчики \texttt{\_readCount} и \texttt{\_writeCount} могут переполниться (теоретически через очень долгое время). Использование маски \texttt{\_masksz} для разницы снимает проблему, но сами счётчики растут неограниченно. На практике это несущественно, так как для переполнения 64-битного счётчика потребуются миллиарды операций.
    \item \textbf{Отсутствие метода \texttt{peek()}} — нет возможности посмотреть первый элемент без его извлечения. При необходимости можно использовать \texttt{first()} в \texttt{struct\_ring\_t} или комбинацию \texttt{count()} и индексов.
\end{enumerate}

\subsubsection{Макросы конфигурации}

Эти файлы не используют специальных макросов, кроме стандартных \texttt{ROBO\_EXPORT} и включения \texttt{robosd\_common.hpp}. Однако для использования \texttt{safe\_ring\_t} требуется, чтобы класс-охранник был определён (обычно \texttt{system::guard} из \texttt{robosd\_system.hpp}).

\subsection{Статический массив объектов (\texttt{robosd\_static\_array.hpp})}

Файл \texttt{robosd\_static\_array.hpp} предоставляет шаблон \texttt{array\_t<C, N>}, который реализует контейнер для хранения ограниченного числа объектов типа \texttt{C}, создаваемых статически или глобально. Основная идея — каждый экземпляр вложенного класса \texttt{item} при своём создании автоматически регистрируется в статическом массиве указателей, что позволяет впоследствии обращаться ко всем созданным объектам по индексу. Это полезно для реализации конечных автоматов, таблиц, реестров и других структур, где требуется фиксированное множество экземпляров, известных на этапе компиляции.

\subsubsection{Устройство шаблона \texttt{array\_t<C, N>}}

\begin{itemize}
    \item \texttt{template<class C, unsigned int N> class array\_t} — шаблон, параметризованный типом хранимых объектов \texttt{C} и максимальным количеством \texttt{N}.
    \item Статическая переменная \texttt{count\_} (inline, C++17) — счётчик уже созданных объектов.
    \item Вложенный класс \texttt{item : public C} — наследник \texttt{C}. В своём конструкторе проверяет, не превышен ли лимит \texttt{N}; если нет, сохраняет указатель на себя в статическом массиве \texttt{index\_} и увеличивает \texttt{count\_}. В противном случае вызывает аварийное завершение (\texttt{ROBO\_APP\_CRASH()}).
    \item Статический массив \texttt{index\_} размера \texttt{N}, содержащий указатели на созданные объекты (инициализируется нулями).
    \item Метод \texttt{unsigned int count()} возвращает текущее количество созданных объектов.
    \item Оператор \texttt{C\& operator[](unsigned int \_no)} возвращает ссылку на объект по индексу (с проверкой выхода за границы и аварийным завершением при ошибке).
\end{itemize}

Все созданные объекты должны быть глобальными или статическими, чтобы их конструкторы вызывались до \texttt{main()}. После этого массив \texttt{index\_} содержит упорядоченный список всех экземпляров в порядке их создания.

\subsubsection{Пример использования}

Предположим, мы хотим создать конечный автомат с фиксированным набором состояний, каждое из которых представлено отдельным классом-наследником.

\begin{lstlisting}[caption={State registry using array\_t}]
#include "core/robosd_static_array.hpp"
#include <iostream>

// Base state class
struct State {
    virtual void enter() = 0;
    virtual const char* name() = 0;
};

// Container for storing all states (max 5)
using StateRegistry = robo::array_t<State, 5>;

// Concrete states
class IdleState : public StateRegistry::item {
public:
    void enter() override { std::cout << "Enter Idle\n"; }
    const char* name() override { return "Idle"; }
} idle; // global object

class RunState : public StateRegistry::item {
public:
    void enter() override { std::cout << "Enter Run\n"; }
    const char* name() override { return "Run"; }
} run;

class ErrorState : public StateRegistry::item {
public:
    void enter() override { std::cout << "Enter Error\n"; }
    const char* name() override { return "Error"; }
} error;

int main() {
    StateRegistry states;
    std::cout << "Total states: " << states.count() << "\n";
    for (unsigned i = 0; i < states.count(); ++i) {
        std::cout << i << ": " << states[i].name() << "\n";
    }
    // Output:
    // Total states: 3
    // 0: Idle
    // 1: Run
    // 2: Error
}
\end{lstlisting}

\subsubsection{Потенциальные проблемы и устаревшие места}

\begin{enumerate}
    \item \textbf{Зависимость от порядка инициализации}: объекты создаются в порядке, определённом компоновщиком (обычно в порядке объявления в единице трансляции). Это может быть непредсказуемо при использовании нескольких файлов. Если порядок важен, его нужно контролировать явно.
    \item \textbf{Отсутствие защиты от дублирования}: ничто не мешает создать несколько экземпляров одного и того же класса; они все попадут в массив. Иногда это желательно, иногда — нет.
    \item \textbf{Аварийное завершение при переполнении}: если создано больше \texttt{N} объектов, вызывается \texttt{ROBO\_APP\_CRASH()}. Это приемлемо для встраиваемых систем, где количество объектов известно статически, но может быть неожиданным при динамическом создании (которое, впрочем, не предполагается).
    \item \textbf{Использование C++17 inline static}: переменная \texttt{count\_} объявлена как \texttt{inline static unsigned int count\_ = 0;}. Это корректно, но требует поддержки C++17. В более старых стандартах пришлось бы выносить определение в .cpp файл. Фреймворк, вероятно, предполагает использование современного компилятора.
    \item \textbf{Обработка ошибок в операторе \texttt{[]}}: при выходе за границы вызывается \texttt{ROBO\_APP\_CRASH()}, но после этого идёт \texttt{return *index\_[0]}; — это dead code (никогда не выполняется). Можно удалить или заменить на ассерт.
    \item \textbf{Ограниченная применимость}: шаблон полезен только для глобальных/статических объектов. Для динамически создаваемых объектов он не подходит, так как требует знания максимального числа на этапе компиляции и не освобождает память.
    \item \textbf{Отсутствие итераторов}: нет стандартных методов begin/end, что затрудняет использование в алгоритмах STL. Впрочем, для встраиваемых систем это не критично.
\end{enumerate}

\subsubsection{Макросы конфигурации}

Файл использует \texttt{ROBO\_APP\_CRASH()} из \texttt{robosd\_common.hpp}. Никаких специальных макросов для включения/отключения не предусмотрено.

\subsection{Конечные автоматы (\texttt{robosd\_stateflow.hpp}, \texttt{robosd\_stateflow.cpp})}

Модуль \texttt{robosd\_stateflow} предоставляет реализацию конечных автоматов двух уровней сложности: простой переключатель состояний (\texttt{stateflow}) и более мощный контроллер задач (\texttt{controller}) с поддержкой вложенных состояний и процессов. Эти механизмы широко используются в фреймворке для управления поведением модулей, обработки событий и реализации асинхронных операций.

\subsubsection{Класс \texttt{stateflow} и узел \texttt{node}}

\texttt{stateflow} — простейший переключатель, который хранит текущий активный узел (наследник \texttt{node}) и позволяет переключаться между узлами. Каждый узел должен реализовать три виртуальных метода:

\begin{itemize}
    \item \texttt{onEnter()} — вызывается при активации узла.
    \item \texttt{doRun()} — вызывается при каждом вызове \texttt{run()} у переключателя, пока узел активен.
    \item \texttt{onLeave()} — вызывается при деактивации узла.
\end{itemize}

Метод \texttt{switchto(node*)} деактивирует текущий узел (если есть), устанавливает новый и активирует его. Метод \texttt{run()} просто делегирует вызов текущему узлу.

Этот примитив удобен для реализации конечных автоматов с небольшим числом состояний, где каждое состояние представлено отдельным объектом.

\subsubsection{Класс \texttt{controller} и иерархия задач}

\texttt{controller} реализует диспетчер задач, который управляет переключением между несколькими процессами или узлами. Он содержит два указателя: \texttt{selected\_} (задача, которую нужно запустить) и \texttt{runned\_} (текущая выполняемая задача). Методы:

\begin{itemize}
    \item \texttt{switchto(base*)} — выбирает новую задачу. Если текущая задача совпадает с выбранной, она просто перезапускается (вызовом \texttt{start\_()}). Иначе текущая задача останавливается (\texttt{stop()}), а выбранная запоминается.
    \item \texttt{stop()} — останавливает текущую задачу и сбрасывает выбранную.
    \item \texttt{run()} — вызывается в цикле. Если нет выполняемой задачи, но есть выбранная, запускает её. Затем вызывает \texttt{run()} у текущей задачи; если та возвращает \texttt{true} (завершение), задача снимается.
    \item \texttt{terminate()} — принудительно завершает текущую задачу (вызовом \texttt{terminate()}) и вызывает собственный \texttt{doTerminate()}.
\end{itemize}

Базовый класс \texttt{base} определяет интерфейс для всех задач:

\begin{itemize}
    \item \texttt{void start\_()} — устанавливает внутреннюю команду \texttt{start}.
    \item \texttt{void stop()} — устанавливает команду \texttt{stop}.
    \item \texttt{commands command()} — возвращает текущую команду.
    \item \texttt{virtual bool run()} — основной метод, который должен быть реализован в наследниках. Возвращает \texttt{true}, если задача завершена и может быть снята.
    \item \texttt{virtual void terminate()} — принудительное завершение.
\end{itemize}

\subsubsection{Класс \texttt{node} (вложенный в \texttt{controller})}

Этот класс реализует типовой узел с четырьмя состояниями:

\begin{description}
    \item[stopped] — начальное состояние.
    \item[entered] — узел вошёл в активную фазу.
    \item[execute] — выполнение.
    \item[leaved] — завершение.
\end{description}

При получении команды \texttt{start} автомат последовательно переходит через состояния, вызывая соответствующие виртуальные методы:

\begin{itemize}
    \item \texttt{onEnter()} — при входе в состояние \texttt{entered}.
    \item \texttt{doEnter()} — опрос готовности (возвращает \texttt{result::success} или \texttt{wait}).
    \item \texttt{onExecute()} — при переходе в \texttt{execute}.
    \item \texttt{doExecute()} — основной цикл выполнения (вызывается многократно).
    \item \texttt{onLeave()} — при переходе в \texttt{leaved}.
    \item \texttt{doLeave()} — опрос завершения.
    \item \texttt{onFinish()} — по окончании.
    \item \texttt{onTerminate()} — при принудительном завершении.
    \item \texttt{onIdle()} — вызывается, когда узел находится в состоянии \texttt{stopped} и нет команды \texttt{start}.
\end{itemize}

Если команда \texttt{stop} поступает во время выполнения, автомат переходит в состояние \texttt{leaved} и пытается завершиться.

Метод \texttt{run()} реализует всю логику переходов и возвращает \texttt{true}, когда узел окончательно останавливается (после успешного \texttt{doLeave} и \texttt{onFinish}).

\subsubsection{Класс \texttt{process} (вложенный в \texttt{controller})}

\texttt{process} представляет более сложный процесс с шестью состояниями:

\begin{description}
    \item[stopped] — остановлен.
    \item[prepare] — подготовка.
    \item[startup] — запуск.
    \item[execute] — выполнение.
    \item[shutdown] — завершение.
    \item[relax] — отдых (пауза перед возможным перезапуском).
\end{description}

Управление аналогично \texttt{node}, но с большим числом фаз. Виртуальные методы:

\begin{itemize}
    \item \texttt{onPrepare()}, \texttt{onStartup()}, \texttt{onExecute()}, \texttt{onShutdown()}, \texttt{onRelax()} — уведомления о входе в состояние.
    \item \texttt{doPrepare()}, \texttt{doStartup()}, \texttt{doExecute()}, \texttt{doShutdown()}, \texttt{doRelax()} — опросы готовности (возвращают \texttt{result}).
    \item \texttt{onFinish()} — вызывается после завершения relax, может вернуть \texttt{repeate::go} для автоматического перезапуска.
    \item \texttt{doIdle()} — вызывается в состоянии \texttt{stopped} при отсутствии команды.
    \item \texttt{doTerminate()} — при принудительном завершении.
\end{itemize}

Метод \texttt{restart()} позволяет сбросить состояние из \texttt{relax} в \texttt{stopped} для возможности повторного запуска без полной остановки.

Класс \texttt{independed} — простой наследник \texttt{process}, который открывает метод \texttt{run()} как публичный, что позволяет использовать процесс независимо от контроллера.

\subsubsection{Пример использования}

Рассмотрим типичный сценарий: контроллер управляет двумя процессами, переключаясь между ними.

\begin{lstlisting}[caption={Controller usage example}]
#include "core/robosd_stateflow.hpp"
#include <iostream>

class MyProcess : public robo::controller::process {
protected:
    result doPrepare() override {
        std::cout << "Preparing\n";
        return result::success;
    }
    result doStartup() override {
        std::cout << "Starting\n";
        return result::success;
    }
    result doExecute() override {
        std::cout << "Executing\n";
        // simulate some work, then finish
        static int count = 0;
        if (++count >= 3) return result::success;
        return result::wait;
    }
    result doShutdown() override {
        std::cout << "Shutting down\n";
        return result::success;
    }
    result doRelax() override {
        std::cout << "Relaxing\n";
        return result::success;
    }
    repeate onFinish() override {
        std::cout << "Finished\n";
        return repeate::no; // do not repeat
    }
};

int main() {
    MyProcess proc1, proc2;
    robo::controller ctrl;

    ctrl.switchto(&proc1);
    for (int i = 0; i < 10; ++i) {
        ctrl.run();
        if (i == 4) ctrl.switchto(&proc2);
    }
    ctrl.terminate();
    return 0;
}
\end{lstlisting}

\subsubsection{Потенциальные проблемы и замечания}

\begin{enumerate}
    \item \textbf{Сложность отладки}: из-за большого числа состояний и виртуальных вызовов трассировка может быть затруднена. Рекомендуется использовать логирование внутри переопределяемых методов.
    \item \textbf{Потокобезопасность}: ни один из классов не синхронизирован. Предполагается, что все вызовы происходят из одного контекста (например, из backend-цикла). Если контроллер используется из нескольких потоков, необходима внешняя синхронизация.
    \item \textbf{Управление памятью}: объекты процессов и узлов обычно создаются статически или глобально, и контроллер хранит только указатели. Ответственность за время жизни лежит на пользователе.
    \item \textbf{Неполная обработка ошибок}: методы \texttt{do...} возвращают лишь \texttt{success} или \texttt{wait}. Нет возможности сообщить об ошибке. В случае сбоя процесс может зависнуть.
    \item \textbf{Отсутствие приоритетов}: контроллер не поддерживает приоритеты задач; переключение происходит только по команде \texttt{switchto}.
    \item \textbf{Устаревший код}: в \texttt{controller::process::run()} присутствует оператор \texttt{goto}-подобное поведение через \texttt{case} без \texttt{break} (fall-through). Это намеренное использование для реализации последовательных переходов, но может сбивать с толку.
    \item \textbf{Двойное назначение \texttt{result}}: в \texttt{process} определён свой \texttt{result}, конфликтующий с глобальным \texttt{robo::result}. Это может вызвать путаницу, но они находятся в разных пространствах имён.
\end{enumerate}

\subsubsection{Макросы конфигурации}

Модуль не использует специальных макросов, кроме стандартных включений. Для работы требуется \texttt{robosd\_common.hpp}.

\subsection{Древовидная структура элементов (\texttt{robosd\_tree.hpp}, \texttt{robosd\_tree.cpp})}

Модуль \texttt{robosd\_tree} предоставляет реализацию иерархической структуры узлов, где каждый узел имеет имя и может содержать дочерние узлы. Это позволяет организовать данные или объекты в виде дерева с доступом по составному пути. Класс \texttt{item} является базовым для построения таких деревьев и используется, например, для представления конфигурационных деревьев, меню, иерархий состояний и т.п.

\subsubsection{Класс \texttt{item}}

Класс \texttt{item} представляет узел дерева. Его основные поля:

\begin{itemize}
    \item \texttt{item* branch\_} — указатель на родительский узел (ветвь). Для корневого узла равен \texttt{nullptr}.
    \item \texttt{list* childs\_} — указатель на список дочерних узлов (тип \texttt{tree::item::list}, который является \texttt{robo::list::unique<item, int>}). Создаётся динамически только при добавлении первого потомка.
    \item \texttt{list::ref ref\_} — интрузивная ссылка для включения в список дочерних узлов родителя. Ключом служит хеш имени узла (\texttt{hash(name)}).
    \item \texttt{string name} — публичное поле, хранящее имя узла.
    \item \texttt{int key() const} — возвращает ключ (хеш имени), под которым узел хранится в списке родителя.
\end{itemize}

Конструктор \texttt{item(cstr \_name, item* \_branch)} инициализирует имя, родителя и ссылку. Если родитель задан, он создаёт у родителя список \texttt{childs\_} (если его ещё нет) и вставляет себя в этот список.

Деструктор удаляет список дочерних узлов (но не сами дочерние узлы, так как они удаляются отдельно; однако список \texttt{childs\_} владеет ссылками, а не объектами — объекты должны удаляться явно или через \texttt{clean}).

\subsubsection{Виртуальные методы для наследников}

\begin{itemize}
    \item \texttt{virtual bool do\_load(cstr \_path)} — вызывается во время загрузки узла. Ему передаётся полный путь к данному узлу (начиная с корня). По умолчанию возвращает \texttt{true}. Наследники могут переопределить этот метод для загрузки своих данных из конфигурации, используя переданный путь.
    \item \texttt{virtual void do\_clean(void)} — вызывается во время очистки дерева. По умолчанию пуст. Предназначен для освобождения ресурсов, специфичных для узла.
\end{itemize}

\subsubsection{Методы для работы с деревом}

\begin{itemize}
    \item \texttt{size\_t path\_len(size\_t \_current, size\_t \_max)} — рекурсивно вычисляет максимальную длину пути (в символах) от корня до самого глубокого листа, учитывая текущую длину \texttt{\_current}. Используется для выделения буфера достаточного размера при построении полных путей.
    \item \texttt{bool load(char\_t* \_beg\_path, char\_t* \_end\_path, size\_t \_path\_size)} — внутренний рекурсивный метод загрузки. \texttt{\_beg\_path} указывает на начало буфера с полным путём к корню, \texttt{\_end\_path} — на текущую позицию в буфере (куда дописывается имя текущего узла). Метод дописывает имя текущего узла (с точкой-разделителем), вызывает \texttt{do\_load(\_beg\_path)} для текущего узла, затем рекурсивно обходит всех потомков.
    \item \texttt{bool load(cstr \_path)} — публичный метод, инициирующий загрузку всего дерева от корня. Он вычисляет необходимый размер буфера (\texttt{path\_len}), выделяет его динамически, формирует начальный путь из \texttt{\_path} и вызывает рекурсивную загрузку.
    \item \texttt{void clean(void)} — рекурсивно вызывает \texttt{do\_clean()} для всех узлов дерева.
    \item \texttt{template <typename T> bool forall(char\_t* \_beg\_path, char\_t* \_end\_path, size\_t \_path\_size, T \_fun)} — шаблонный метод для обхода всех узлов с вызовом функтора \texttt{\_fun}. Функтор должен иметь сигнатуру \texttt{bool (item\&, char\_t*)} или аналогичную. Ему передаётся ссылка на текущий узел и начало пути (полный путь к узлу). Метод строит путь аналогично \texttt{load} и рекурсивно обходит потомков. Возвращает \texttt{false}, если какой-либо вызов функтора вернул \texttt{false}.
\end{itemize}

\subsubsection{Пример использования}

Предположим, мы хотим построить дерево конфигурации, где каждый узел соответствует разделу, и загрузить значения из INI-файла.

\begin{lstlisting}[caption={Configuration tree example}]
#include "core/robosd_tree.hpp"
#include "core/robosd_ini.hpp"

class ConfigNode : public robo::tree::item {
public:
    int value;
    ConfigNode(const char* name, ConfigNode* parent)
        : robo::tree::item(name, parent), value(0) {}
protected:
    bool do_load(cstr path) override {
        // path contains the full path to the node, e.g., "root.section.subsection"
        return robo::ini::load(path, "value", value);
    }
    void do_clean() override {
        // free resources if needed
    }
};

int main() {
    // Create root node (no parent)
    ConfigNode root("root", nullptr);
    // Create child nodes
    ConfigNode section1("section1", &root);
    ConfigNode section2("section2", &root);
    ConfigNode subsection("subsection", &section1);

    // Load the whole tree from INI file (root path is section name)
    if (root.load("")) { // empty string - root
        // values loaded
    }

    // Traverse all nodes and print their values
    auto print = [](robo::tree::item& node, char_t* path) -> bool {
        ConfigNode& cfg = static_cast<ConfigNode&>(node);
        robo_infolog("Path: %S, value: %d", path, cfg.value);
        return true;
    };
    char_t buffer[256];
    root.forall(buffer, buffer, 256, print);
}
\end{lstlisting}

\subsubsection{Потенциальные проблемы и замечания}

\begin{enumerate}
    \item \textbf{Управление памятью}: список дочерних узлов (\texttt{childs\_}) создаётся динамически и хранит только ссылки на узлы. Сами узлы обычно создаются статически или глобально, и их время жизни должно превышать время жизни дерева. Деструктор \texttt{item} удаляет список, но не сами дочерние узлы — это может привести к утечке, если узлы были созданы динамически. Рекомендуется либо создавать узлы статически, либо использовать умные указатели, либо явно удалять их перед уничтожением корня.
    \item \textbf{Использование хеша имени в качестве ключа}: это означает, что в пределах одного родителя имена узлов должны быть уникальны, иначе они не смогут быть вставлены в \texttt{unique}-список (вставка вернёт \texttt{false}). В конструкторе это не проверяется, что может привести к тому, что узел не будет добавлен. Следует контролировать уникальность имён на этапе создания.
    \item \textbf{Рекурсивные методы}: \texttt{path\_len}, \texttt{load}, \texttt{forall} используют рекурсию, что может привести к переполнению стека при очень глубоких деревьях. Для встраиваемых систем с ограниченным стеком это стоит учитывать.
    \item \textbf{Выделение буфера в \texttt{load(cstr)}}: буфер выделяется оператором \texttt{new[]} и не освобождается в случае ошибки (если \texttt{ROBO\_LRET} приводит к досрочному выходу). Это потенциальная утечка. Рекомендуется использовать RAII-обёртку или явно освобождать память.
    \item \textbf{Закомментированный код}: в \texttt{robosd\_tree.cpp} присутствует закомментированная версия метода \texttt{forall}, которая ранее была функцией с указателем. Это не влияет на работу, но может быть удалено для чистоты.
    \item \textbf{Отсутствие виртуального деструктора в \texttt{item}?} В классе \texttt{item} объявлен виртуальный деструктор (\texttt{virtual ~item(void)}), что корректно для полиморфного использования.
\end{enumerate}

\subsubsection{Макросы конфигурации}

Модуль использует \texttt{ROBO\_UNICODE\_ENABLED} для определения длины строки в \texttt{load(cstr)} и \texttt{ROBO\_APP\_ASSERT} для проверок. Никаких специальных макросов включения не требуется.

\subsection{Каркас приложения и поддержка модулей (\texttt{robosd\_app.hpp}, \texttt{robosd\_app.cpp})}

Файлы \texttt{robosd\_app.hpp} и \texttt{robosd\_app.cpp} реализуют основу приложения на базе узлов (\texttt{node}) и механизм динамической загрузки модулей. Классы \texttt{node}, \texttt{module}, \texttt{wrapper} и \texttt{machine} образуют иерархию, позволяющую организовать приложение как дерево компонентов с единым жизненным циклом и возможностью подгружать функциональность из внешних динамических библиотек.

\subsubsection{Класс \texttt{node} — базовый узел иерархии}

\texttt{node} является базовым классом для всех компонентов приложения. Он инкапсулирует:

\begin{itemize}
    \item Состояние узла (\texttt{state}): \texttt{unknown}, \texttt{panic}, \texttt{clean}, \texttt{stopped}, \texttt{startup}, \texttt{execute}, \texttt{shutdown}.
    \item Связь с родителем (\texttt{owner\_}) и списки дочерних узлов, разделённые по состояниям: \texttt{disabled\_}, \texttt{stopped\_}, \texttt{startupped\_}, \texttt{active\_}, \texttt{shutdowned\_}. Это позволяет эффективно управлять группами узлов.
    \item Имя узла (\texttt{name\_}) и алиас (\texttt{alias\_}), загружаемый из INI.
    \item Методы жизненного цикла: \texttt{load()}, \texttt{start()}, \texttt{startup()}, \texttt{stop()}, \texttt{shutdown()}, \texttt{clean()}, \texttt{panic()}. Каждый метод обновляет состояние и делегирует вызовы дочерним узлам.
    \item Виртуальные методы \texttt{do\_load()}, \texttt{do\_clean()}, \texttt{do\_start()}, \texttt{do\_stop()}, \texttt{do\_startup()}, \texttt{do\_shutdown()}, \texttt{do\_panic()}, которые переопределяются в производных классах.
    \item Поддержка пути в конфигурации через вложенные классы \texttt{path\_root} и \texttt{path}. Это позволяет строить иерархический путь (например, \texttt{"root.child.grandchild"}) и использовать его для поиска секций в INI-файлах.
    \item Статические методы \texttt{find()} для поиска узла по полному пути или хешу.
\end{itemize}

Узлы автоматически регистрируются в глобальном индексе (\texttt{index\_()}) по хешу полного пути или алиаса, что позволяет быстро находить их по имени.

\subsubsection{Класс \texttt{module} — основа для модулей}

\texttt{module} наследует \texttt{node} и добавляет два чистых виртуальных метода:

\begin{itemize}
    \item \texttt{virtual void frontend\_loop() = 0} — вызывается в цикле frontend (обычно в контексте ОС) для обработки событий, взаимодействия с пользователем и т.п.
    \item \texttt{virtual void backend\_loop() = 0} — вызывается в цикле backend (в контексте реального времени) для выполнения критических операций.
\end{itemize}

Все динамически загружаемые модули должны быть наследниками \texttt{module} и реализовывать эти методы.

\subsubsection{Класс \texttt{wrapper} — обёртка для динамической загрузки модулей}

\texttt{wrapper} отвечает за загрузку разделяемой библиотеки, создание экземпляра модуля и управление его временем жизни. Он использует системные функции \texttt{system::lib} для работы с библиотеками.

Основные поля:
\begin{itemize}
    \item \texttt{void* handle\_} — дескриптор загруженной библиотеки.
    \item \texttt{module* module\_} — указатель на созданный экземпляр модуля.
    \item \texttt{string lib\_} — имя библиотеки (или путь), загружаемое из INI.
    \item \texttt{ref ref\_} — ссылка для включения в список загруженных модулей машины (\texttt{machine::wrappers\_}).
\end{itemize}

Методы:
\begin{itemize}
    \item \texttt{bool begin\_(cstr \_key)} — загружает имя библиотеки из INI по ключу \texttt{\_key} (секция \texttt{MODULES}), проверяет существование файла, загружает библиотеку, получает адреса функций \texttt{robo\_module\_query} и \texttt{robo\_module\_release}, вызывает \texttt{query} для получения указателя на модуль, инициализирует модуль вызовом \texttt{init(nullptr, \&machine::root())}. В случае ошибки освобождает ресурсы и возвращает \texttt{false}.
    \item \texttt{static bool begin(cstr \_key)} — создаёт новый объект \texttt{wrapper} и вызывает \texttt{begin\_}. Используется для загрузки модулей по одному.
    \item \texttt{void finish()} — завершает работу модуля: вызывает \texttt{module\_->init(nullptr, nullptr)} (отсоединение от дерева), получает функцию \texttt{release} и освобождает модуль, затем выгружает библиотеку.
\end{itemize}

Важно, что модуль после создания автоматически становится дочерним узлом корневой машины, что включает его в общий жизненный цикл приложения.

\subsubsection{Класс \texttt{machine} — корневой узел и диспетчер приложения}

\texttt{machine} реализован как синглтон и является корнем дерева узлов. Он управляет всем приложением: инициализацией системы, загрузкой конфигурации, запуском и остановкой модулей, а также циклическими вызовами frontend/backend.

Основные компоненты:
\begin{itemize}
    \item \texttt{wrapper::map wrappers\_} — список загруженных модулей (обёрток).
    \item Флаги состояния: \texttt{terminated\_}, \texttt{req\_state\_} (запрос на запуск или остановку).
    \item Методы жизненного цикла приложения:
        \begin{itemize}
            \item \texttt{bool begin(cstr \_ini)} — инициализация: создаёт корневой узел, вызывает \texttt{system::begin()}, открывает INI-файл, настраивает логирование (если включено), загружает дерево узлов (\texttt{node::load()}), загружает модули согласно списку из INI.
            \item \texttt{bool start()} — переводит приложение в режим запуска (устанавливает \texttt{req\_state\_ = start}).
            \item \texttt{void stop()} — запрашивает остановку.
            \item \texttt{void frontend\_loop()} — должен вызываться в цикле frontend; обновляет состояние машины, вызывает \texttt{frontend\_loop} всех модулей и системные обработчики.
            \item \texttt{void backend\_loop()} — должен вызываться в цикле backend; обновляет состояние машины, вызывает \texttt{backend\_loop} всех модулей, а также обрабатывает очереди и задачи реального времени.
            \item \texttt{bool terminated()} — проверяет, завершилось ли приложение.
            \item \texttt{void finish()} — завершает работу: очищает дерево узлов, закрывает INI, логирование и систему.
        \end{itemize}
\end{itemize}

При загрузке модулей \texttt{machine} читает из секции \texttt{MODULES} параметр \texttt{COUNT}, а затем для каждого индекса загружает имя библиотеки из ключа \texttt{M\_<i>} и создаёт \texttt{wrapper}. Все загруженные модули автоматически включаются в дерево и получают вызовы циклов.

\subsubsection{Процесс загрузки модуля (по шагам)}

\begin{enumerate}
    \item В INI-файле указывается количество модулей и имена ключей:
\begin{verbatim}
[MODULES]
COUNT = 2
M_1 = "motor_control"
M_2 = "network_interface"
\end{verbatim}
    \item \texttt{machine::do\_load()} читает \texttt{COUNT} и для каждого \texttt{i} вызывает \texttt{wrapper::begin(string(RT("M\_\%d"), i+1))}.
    \item \texttt{wrapper::begin\_} загружает имя библиотеки из INI, вызывает \texttt{system::lib::exists()}, затем \texttt{system::lib::load()}.
    \item Из библиотеки извлекаются функции \texttt{robo\_module\_query} и \texttt{robo\_module\_release} (с учётом \texttt{ROBO\_EXPORT\_FUNCTION\_PREFIX}).
    \item Вызов \texttt{query()} возвращает указатель на объект \texttt{module*}.
    \item У модуля вызывается \texttt{init(nullptr, \&machine::root())}, что подключает его к дереву (родитель — корневая машина).
    \item Обёртка регистрируется в списке \texttt{wrappers\_} корневой машины.
    \item При вызове \texttt{machine::frontend\_loop()} и \texttt{backend\_loop()} происходит итерация по всем загруженным модулям и вызов соответствующих методов.
\end{enumerate}

\subsubsection{Пример создания простого модуля}

\begin{lstlisting}[caption={Пример модуля}]
#include "core/robosd_app.hpp"

class MyModule : public robo::app::module {
public:
    MyModule() : module("MyModule") {}
protected:
    void frontend_loop() override {
        // Called in frontend context (e.g., every 100 ms)
        robo_infolog("MyModule frontend tick");
    }
    void backend_loop() override {
        // Called in backend context (e.g., every 1 ms)
        // Perform time-critical operations
    }
};

// Entry points required for dynamic loading
extern "C" {
    robo::app::module* robo_module_query() {
        static MyModule instance;
        return &instance;
    }
    void robo_module_release(robo::app::module*) {
        // Nothing to do for static instance
    }
}
\end{lstlisting}

Для статической компоновки можно использовать макросы из \texttt{robosd\_system\_module\_reg.hpp}, которые автоматически генерируют эти функции.

\subsubsection{Особенности и рекомендации}

\begin{itemize}
    \item Модули должны быть спроектированы так, чтобы их \texttt{backend\_loop} был максимально быстрым и неблокирующим, так как он выполняется в контексте реального времени.
    \item Для обмена данными между модулями и между frontend/backend рекомендуется использовать очереди (\texttt{ring\_safe\_buf}) или разделяемую память с синхронизацией.
    \item Время жизни модулей управляется \texttt{wrapper} и \texttt{machine}; при завершении приложения все модули корректно останавливаются и выгружаются.
    \item Возможна также статическая компоновка модулей (без динамической загрузки) путём прямой регистрации через \texttt{robo::native::lib} (см. \texttt{robosd\_system\_native.hpp}).
\end{itemize}

\subsubsection{Потенциальные проблемы и устаревшие места}

\begin{enumerate}
    \item \textbf{Отсутствие проверки версий}: при загрузке динамической библиотеки не проверяется совместимость интерфейсов. Это может привести к аварийному завершению, если модуль собран с другой версией фреймворка.
    \item \textbf{Утечка памяти при ошибке загрузки}: если \texttt{wrapper::begin\_} завершается неудачей после выделения памяти под обёртку, объект \texttt{wrapper} удаляется, но возможно, не все ресурсы освобождены (например, уже загруженная библиотека не выгружается). В коде это обрабатывается, но стоит проверить все пути.
    \item \textbf{Использование макросов \texttt{ROBO\_LBREAKN\_F} и т.п.} может скрывать логику и затруднять отладку. Рекомендуется заменять на явные проверки с логированием.
    \item \textbf{Жёсткая связь с INI-форматом}: имена модулей всегда берутся из INI. Для некоторых приложений может потребоваться другой источник конфигурации.
\end{enumerate}

\subsubsection{Макросы конфигурации}

Для работы модулей необходимы следующие макросы (определяемые в \texttt{robosd\_app\_tuning.hpp}):

\begin{itemize}
    \item \texttt{ROBO\_APP\_MODULE\_ENABLED} — должен быть равен 1 для включения поддержки модулей.
    \item \texttt{ROBO\_APP\_LIB\_ENABLED} — включает функции работы с динамическими библиотеками.
    \item \texttt{ROBO\_APP\_SERVO\_ENABLED}, \texttt{ROBO\_APP\_TERMINAL\_ENABLED}, \texttt{ROBO\_APP\_TRACE\_ENABLED} — опционально включают дополнительные подсистемы, вызываемые в циклах.
\end{itemize}

\subsection{Матричные операции (\texttt{robosd\_matrix.hpp})}

Файл \texttt{robosd\_matrix.hpp} предоставляет шаблонный класс \texttt{matrix\_t<math>} для работы с матрицами произвольного размера. Он использует внешний шаблонный параметр \texttt{math}, который должен определять тип чисел с плавающей точкой (\texttt{real}) и базовые математические функции: \texttt{fill}, \texttt{assert}, \texttt{sqrt}. Это позволяет адаптировать матричные операции под разные численные типы (float, double) и платформы.

\subsubsection{Класс \texttt{matrix\_t<math>}}

Класс хранит данные в виде двумерного динамического массива \texttt{real** memo}. Размеры хранятся в полях \texttt{rows} и \texttt{cols} (тип \texttt{int}). Основные методы:

\begin{itemize}
    \item Конструкторы: \texttt{matrix\_t(int rows, int cols)} выделяет память; \texttt{matrix\_t()} создаёт пустую матрицу; \texttt{matrix\_t(const matrix\_t\&)} — копирующий.
    \item Деструктор освобождает память.
    \item \texttt{void reinit(int rows, int cols)} — изменяет размер, при необходимости перераспределяя память.
    \item \texttt{void eye()} — заполняет единичную матрицу (предполагается, что матрица квадратная).
    \item \texttt{void zeros()} — обнуляет все элементы.
    \item \texttt{void invto(matrix\_t\& IM)} — вычисляет обратную матрицу и сохраняет в \texttt{IM}. Предполагает квадратную матрицу. Использует метод Гаусса с выбором главного элемента (строки). Код содержит несколько потенциальных проблем (см. ниже).
    \item \texttt{void multto(const matrix\_t\& M2, matrix\_t\& R)} — умножение текущей матрицы на \texttt{M2}, результат в \texttt{R}.
    \item \texttt{void trans\_sqrto(matrix\_t\& R)} — вычисляет произведение текущей матрицы на её транспонированную ($A \cdot A^T$) и сохраняет в \texttt{R} (требуется, чтобы результат был квадратным).
    \item \texttt{void pinvto(matrix\_t\& IM, matrix\_t\& TS, matrix\_t\& ITS)} — псевдообращение (псевдоинверсия) через транспонирование и обращение $A^+ = (A^T A)^{-1} A^T$. Использует временные матрицы \texttt{TS} и \texttt{ITS}.
    \item Оператор присваивания \texttt{operator=}.
\end{itemize}

Все методы, кроме конструкторов и деструктора, предполагают, что память уже выделена соответствующего размера. Проверки на нулевые указатели или выход за границы отсутствуют, полагаясь на \texttt{math::assert} (который может быть пустым).

\subsubsection{Пример использования}

\begin{lstlisting}[caption={Matrix operations example}]
#include "core/robosd_matrix.hpp"

struct MyMath {
    typedef double real;
    static void fill(real* ptr, int n, real val) {
        for (int i = 0; i < n; ++i) ptr[i] = val;
    }
    static void assert(bool cond) { if (!cond) __builtin_trap(); }
    static real sqrt(real x) { return ::sqrt(x); }
};

int main() {
    robo::matrix_t<MyMath> A(2, 2);
    A.memo[0][0] = 4; A.memo[0][1] = 3;
    A.memo[1][0] = 3; A.memo[1][1] = 2;

    robo::matrix_t<MyMath> Ainv(2, 2);
    A.invto(Ainv);
    // Ainv should be [-2, 3; 3, -4]

    robo::matrix_t<MyMath> B(2, 3);
    // fill B...
    robo::matrix_t<MyMath> C(2, 2);
    A.multto(B, C); // C = A * B
}
\end{lstlisting}

\subsubsection{Потенциальные проблемы и замечания}

\begin{enumerate}
    \item \textbf{Отсутствие проверок}: многие методы не проверяют совместимость размеров (кроме \texttt{math::assert} в нескольких местах). Например, \texttt{invto} полагается, что матрица квадратная, но не проверяет это явно.
    \item \textbf{Утечки памяти}: при повторном вызове \texttt{reinit} с теми же размерами память не перераспределяется, но если вызвать \texttt{reinit} с другими размерами, старая память освобождается. Однако если после \texttt{reinit} последует исключение, может произойти утечка. В коде нет обработки исключений.
    \item \textbf{Эффективность}: все операции выполняются с вложенными циклами и динамической памятью. Для встраиваемых систем с ограниченными ресурсами это может быть неприемлемо. Лучше использовать статические матрицы фиксированного размера.
    \item \textbf{Ошибка в \texttt{invto}}: в цикле обмена строк используется некорректный алгоритм. Переменные \texttt{I} и \texttt{K} определены как \texttt{int I, K;}, затем в цикле \texttt{for (I = 0; I < cols; I++)} проверяется \texttt{memo[I][I] == 0}. Если это так, ищется строка \texttt{K} с ненулевым элементом в столбце \texttt{I}. Далее происходит обмен всех элементов строк \texttt{I} и \texttt{K} для обеих матриц (текущей и обратной). Однако вложенные циклы по \texttt{i} и \texttt{j} перебирают все строки и столбцы, что неверно — нужно обменивать только строки \texttt{I} и \texttt{K}. Фактически код перебирает все строки \texttt{i} и обменивает их с \texttt{K}, что разрушает матрицу. Это серьёзная ошибка.
    \item \textbf{Использование \texttt{int} для размеров}: в некоторых архитектурах \texttt{int} может быть 16-битным, что ограничивает максимальный размер матрицы.
    \item \textbf{Отсутствие move-семантики}: копирование матриц выполняется глубоко, что неэффективно для больших матриц.
    \item \textbf{Метод \texttt{pinvto}} вычисляет псевдообращение через $(A^T A)^{-1} A^T$, что работает только для матриц полного ранга и может быть численно неустойчивым. Нет проверки обусловленности.
\end{enumerate}

\subsubsection{Макросы конфигурации}

Файл не использует специальных макросов, только включает \texttt{robosd\_common.hpp} (неявно через зависимости). Однако он требует определения \texttt{math::real}, \texttt{math::fill}, \texttt{math::assert}, \texttt{math::sqrt} пользователем.

\subsection{Конвейер цифровой обработки сигналов (\texttt{robosd\_zconveyer.hpp})}

Файл \texttt{robosd\_zconveyer.hpp} реализует сложную иерархию классов для организации конвейерной обработки дискретных сигналов с различными частотами дискретизации. По своей сути это гибкий инструмент для построения систем сбора и обработки данных, где каждый блок (актор) может иметь свой темп выполнения, а сигналы передаются между блоками через кольцевые буферы.

\subsubsection{Общая архитектура}

В основе лежит несколько ключевых классов:

\begin{itemize}
    \item \texttt{signal\_s} — представляет собой буфер дискретного сигнала. Содержит кольцевой буфер значений типа \texttt{real}, указатели на первый, последний и текущий элементы, счётчики. Методы:
        \begin{itemize}
            \item \texttt{put(value)} — добавляет новое значение в буфер.
            \item \texttt{value()} — возвращает последнее значение.
            \item \texttt{update\_interval(count)} — вычисляет интервалы для чтения последовательности из \texttt{count} последних элементов с учётом кольцевой природы буфера.
            \item Виртуальные \texttt{clean()}, \texttt{doclean()} для сброса.
        \end{itemize}
    \item \texttt{actor\_s} — наследник \texttt{signal\_s}, представляет обрабатывающий блок. Содержит ссылку на входной сигнал \texttt{input} и чисто виртуальный метод \texttt{doit()}, который должен быть реализован в конкретных акторах. Акторы могут быть объединены в список.
    \item \texttt{counter\_s} — счётчик тактов. Связан с определённой частотой дискретизации (периодом \texttt{presc}). Хранит список акторов, которые должны выполняться с этой частотой. Метод \texttt{run()} уменьшает счётчик и при достижении нуля вызывает \texttt{doit()} всех своих акторов.
    \item \texttt{sample\_s} — группирует несколько счётчиков с одинаковым базовым периодом (но с разными предделителями). Содержит список \texttt{counter\_s} и метод \texttt{run()}, который запускает все свои счётчики.
    \item \texttt{zconveyer\_t<M>} — главный класс-контейнер, содержащий список сэмплов (\texttt{samples}) и фабричные методы для создания акторов. Также содержит множество встроенных типов акторов (см. ниже).
\end{itemize}

Вся система работает следующим образом: внешний код периодически вызывает \texttt{run()} у корневого конвейера (или у нужных сэмплов). Те, в свою очередь, запускают счётчики, которые по своим тактам активируют акторы. Акторы читают значения из входных сигналов (возможно, с задержкой) и записывают результаты в свои собственные буферы, которые могут быть входами для других акторов.

\subsubsection{Встроенные типы акторов}

Внутри \texttt{zconveyer\_t} определены многочисленные классы-акторы, реализующие различные функции обработки сигналов:

\begin{itemize}
    \item \texttt{scale\_s} — линейное масштабирование: $y = gain \cdot x(t - delay) + offset$.
    \item \texttt{delay\_s} — задержка: $y = x(t - delay)$.
    \item \texttt{diff\_s} — разность с другим сигналом: $y = x - y_{ref}$.
    \item \texttt{sqr\_s} — квадрат: $y = x^2$.
    \item \texttt{sqrt\_s} — квадратный корень: $y = \sqrt{x}$.
    \item \texttt{slide\_mean\_s} — скользящее среднее по окну длины \texttt{len}.
    \item \texttt{slide\_sum\_s} — скользящая сумма.
    \item \texttt{var0\_s} — квадрат отклонения от среднего: $(x - mean)^2$ (используется для последующего вычисления дисперсии).
    \item \texttt{disp\_s} — дисперсия по скользящему окну (несмещённая оценка, делитель \texttt{len-1}).
    \item \texttt{filter1\_s} — фильтр первого порядка (неполная реализация, использует поля \texttt{tau} и \texttt{itau}, не определённые в коде).
    \item \texttt{zero\_ord\_holder\_s} — фиксатор нулевого порядка (просто передаёт входной сигнал).
    \item \texttt{cubic\_approxx\_s} — аппроксимация кубическим полиномом методом наименьших квадратов по окну. Использует матрицы для вычисления коэффициентов. Содержит вложенные сигналы для коэффициентов \texttt{A}, \texttt{B}, \texttt{C}, \texttt{D} и выходной сигнал ошибки.
    \item \texttt{cubic\_interp\_s} — интерполяция кубическим полиномом, используя коэффициенты из \texttt{cubic\_approxx\_s}.
\end{itemize}

Каждый актор имеет свою структуру конфигурации (\texttt{config\_s}), которая должна быть передана при создании. Для доступа к конфигурации внутри метода \texttt{doit()} используется макрос \texttt{Z\_CONFIG\_S(s)}, который извлекает ссылку на конфигурацию нужного типа. Это опасный приём, так как основан на \texttt{reinterpret\_cast} без проверки типа.

\subsubsection{Пример использования (гипотетический)}

Предположим, мы хотим создать простой конвейер, который принимает входной сигнал, масштабирует его и вычисляет скользящее среднее.

\begin{lstlisting}[caption={Simple signal processing pipeline}]
#include "core/robosd_zconveyer.hpp"

struct MyMath {
    typedef float real;
    static void fill(real*, int, real) {}
    static void assert(bool cond) {}
    static real sqrt(real x) { return x > 0 ? ::sqrtf(x) : 0; }
};

int main() {
    robo::zconveyer_t<MyMath> cv;

    // Create a sample with base period 1000 (e.g., 1 kHz)
    auto* sample = new robo::zconveyer_t<MyMath>::sample_s(1000, cv);

    // Create a counter with prescaler 1 (i.e., runs every sample)
    auto* counter = new robo::zconveyer_t<MyMath>::counter_s(1, 0, *sample);

    // Create an input signal (must be attached to some source)
    robo::zconveyer_t<MyMath>::signal_s input(*counter, {64}); // buffer size 64

    // Create a scale actor
    robo::zconveyer_t<MyMath>::scale_s::config_s scale_cfg;
    scale_cfg.actor.signal.size = 64; // output buffer size
    scale_cfg.delay = 2;
    scale_cfg.offset = 0;
    scale_cfg.gain = 2.0;
    auto* scale = new robo::zconveyer_t<MyMath>::scale_s(*counter, input, scale_cfg);

    // Create a sliding mean actor
    robo::zconveyer_t<MyMath>::slide_mean_s::config_s mean_cfg;
    mean_cfg.actor.signal.size = 64;
    mean_cfg.len = 10;
    auto* mean = new robo::zconveyer_t<MyMath>::slide_mean_s(*counter, *scale, mean_cfg);

    // External loop: feed input and run the sample
    for (int t = 0; t < 1000; ++t) {
        input.put(some_external_value(t));
        sample->run();  // triggers counters and actors
        float output = mean->value();
        // use output...
    }
}
\end{lstlisting}

\subsubsection{Потенциальные проблемы и замечания}

\begin{enumerate}
    \item \textbf{Незавершённость кода}: в файле присутствуют опечатки и несоответствия. Например, в методе \texttt{get\_sample} используется \texttt{event\_s} вместо \texttt{sample\_s}. Метод \texttt{get\_counter} пытается обратиться к полю \texttt{id} у \texttt{counter}, которое не определено. Это указывает на то, что код, вероятно, не компилируется в текущем виде.
    \item \textbf{Опасный макрос \texttt{Z\_CONFIG\_S}}: он выполняет \texttt{reinterpret\_cast} для получения конфигурации нужного типа. Если фактический тип конфигурации не совпадает с ожидаемым, возникает неопределённое поведение.
    \item \textbf{Управление памятью}: все объекты (\texttt{sample\_s}, \texttt{counter\_s}, \texttt{actor\_s}) создаются через \texttt{new}, но нет автоматического удаления. Деструктор \texttt{zconveyer\_t} удаляет только \texttt{default\_sample}, но не все созданные объекты. Это приводит к утечкам.
    \item \textbf{Сложность и связанность}: классы тесно переплетены, и для понимания требуется изучить весь код. Отсутствует документация по предполагаемому использованию.
    \item \textbf{Использование списков \texttt{unidir::base\_t}}: они предполагают, что объекты содержат поле \texttt{next\_}, но не все классы его имеют (например, \texttt{sample\_s} имеет \texttt{next\_}, но \texttt{counter\_s} тоже должен его иметь, и он есть). Однако в некоторых местах используются неинициализированные указатели.
    \item \textbf{Неэффективность}: многие операции выполняются в циклах, а буферы сигналов копируют значения, хотя можно было бы использовать ссылки.
    \item \textbf{Отсутствие синхронизации}: конвейер не является потокобезопасным, но может использоваться в однопоточном контексте реального времени.
    \item \textbf{Зависимость от \texttt{matrix\_t}}: для кубической аппроксимации используются матрицы, что может быть избыточно для встраиваемых систем.
\end{enumerate}

\subsubsection{Макросы конфигурации}

Файл не использует специальных макросов, но требует определения типа \texttt{math} (как и \texttt{matrix\_t}) и включает \texttt{robosd\_list.hpp} и \texttt{robosd\_matrix.hpp}.

\subsection{Транзакционный протокол (\texttt{robosd\_tran.h})}

Файл \texttt{robosd\_tran.h} содержит определения структур и перечислений, используемых для организации обмена данными между компонентами системы в рамках транзакционного протокола. Протокол предполагает наличие запросов (GET, PUT, EXCHANGE, REBOOT\_ME) и статусов выполнения (NONE, EXECUTE, COMPLETE, REFUSE). Заголовок транзакции может интерпретироваться либо как пара (идентификатор устройства + команда), либо как единое хеш-значение, что позволяет гибко адаптировать протокол под разные задачи (например, адресация в CAN-сети или идентификация сервисов).

\subsubsection{Настраиваемые типы}

В начале файла определён ряд макросов, которые пользователь может переопределить для изменения разрядности соответствующих полей:

\begin{itemize}
    \item \texttt{ROBO\_TRAN\_COMMAND\_ID} — тип идентификатора команды (по умолчанию \texttt{uint8\_t}).
    \item \texttt{ROBO\_TRAN\_DEV\_ID} — тип идентификатора устройства (по умолчанию \texttt{uint8\_t}).
    \item \texttt{ROBO\_TRAN\_HEADER\_IX} — тип для хеша заголовка (объединение dev\_id и command, по умолчанию \texttt{uint16\_t}).
    \item \texttt{ROBO\_PHY\_CHAN\_IX} — тип индекса физического канала (по умолчанию \texttt{uint8\_t}). Специальное значение \texttt{ROBO\_PHY\_CHAN\_AUTO} (0xFF) означает автоматический выбор канала.
\end{itemize}

Эти макросы позволяют подстроить протокол под конкретную аппаратную платформу и требования по размеру пакетов.

\subsubsection{Перечисления}

\begin{itemize}
    \item \texttt{robo\_tran\_status\_t} — статус выполнения транзакции:
        \begin{itemize}
            \item \texttt{ROBO\_TRAN\_NONE} — нет статуса.
            \item \texttt{ROBO\_TRAN\_EXECUTE} — флаг выполнения (0x10). Может комбинироваться с подфлагами: \texttt{ROBO\_TRAN\_EXECUTE\_START} (начало выполнения) и \texttt{ROBO\_TRAN\_EXECUTE\_PHY} (выполнение на физическом уровне).
            \item \texttt{ROBO\_TRAN\_COMPLETE} — успешное завершение.
            \item \texttt{ROBO\_TRAN\_REFUSE} — отказ.
        \end{itemize}
    \item \texttt{robo\_tran\_request\_t} — тип запроса:
        \begin{itemize}
            \item \texttt{ROBO\_TRAN\_REQUEST\_GET} — чтение данных.
            \item \texttt{ROBO\_TRAN\_REQUEST\_PUT} — запись данных.
            \item \texttt{ROBO\_TRAN\_EXCANGE} — обмен данными (вероятно, чтение+запись).
            \item \texttt{ROBO\_TRAN\_REBOOT\_ME} — запрос на перезагрузку устройства.
        \end{itemize}
\end{itemize}

\subsubsection{Структура заголовка \texttt{robo\_tran\_header\_t}}

Объединение позволяет интерпретировать заголовок двумя способами:

\begin{lstlisting}[caption={Header union}]
typedef union robo_tran_header_s {
    struct {
        robo_tran_dev_id_t dev_id;
        robo_tran_command_id_t command;
    };
    robo_tran_header_ix_t hash;
} robo_tran_header_t;
\end{lstlisting}

Поля \texttt{dev\_id} и \texttt{command} занимают младшие и старшие байты соответственно (порядок байт зависит от платформы). Поле \texttt{hash} позволяет рассматривать заголовок как единое целое, что удобно для быстрого сравнения или использования в хеш-таблицах.

\subsubsection{Основная структура транзакции \texttt{robo\_tran\_t}}

\begin{lstlisting}[caption={Transaction structure}]
typedef struct robo_tran_s {
    size_t size_actual;
    size_t size_max;
    uint8_t * data;
    robo_tran_status_t status;
    robo_tran_request_t request;
    robo_tran_header_t header;
    robo_phy_chan_ix_t phy_chan;
} robo_tran_t;
\end{lstlisting}

Поля:
\begin{itemize}
    \item \texttt{size\_actual} — фактический размер данных в буфере \texttt{data}.
    \item \texttt{size\_max} — максимальная ёмкость буфера.
    \item \texttt{data} — указатель на буфер с данными транзакции.
    \item \texttt{status} — текущий статус выполнения.
    \item \texttt{request} — тип запроса.
    \item \texttt{header} — заголовок (идентификатор устройства и команда).
    \item \texttt{phy\_chan} — индекс физического канала (например, номер CAN-интерфейса или последовательного порта).
\end{itemize}

Данная структура не владеет памятью под \texttt{data} — предполагается, что буфер предоставляется пользователем и его время жизни контролируется внешним кодом.

\subsubsection{Предполагаемое использование}

Транзакционный протокол может применяться для обмена данными между узлами распределённой системы (например, по CAN, Modbus или другой шине). Запросы \texttt{GET} и \texttt{PUT} соответствуют классическим операциям чтения и записи регистров или параметров. \texttt{EXCHANGE} может использоваться для атомарного обмена (например, запрос статуса и получение данных). \texttt{REBOOT\_ME} — служебная команда для перезагрузки удалённого устройства.

Статусы \texttt{EXECUTE\_START} и \texttt{EXECUTE\_PHY} позволяют отслеживать многоэтапные транзакции, например, когда сначала начинается выполнение на логическом уровне, а затем оно передаётся на физический уровень (драйвер).

\subsubsection{Потенциальные проблемы и замечания}

\begin{enumerate}
    \item \textbf{Порядок байт (endianness)}: объединение \texttt{header} предполагает, что поля \texttt{dev\_id} и \texttt{command} располагаются в памяти последовательно. На little-endian системах младший байт будет \texttt{dev\_id}, старший — \texttt{command}. При передаче по сети может потребоваться приведение к единому порядку.
    \item \textbf{Выравнивание}: компилятор может добавить выравнивание между полями структуры внутри объединения, что нарушит предполагаемое соответствие с полем \texttt{hash}. В стандарте C гарантируется, что первый элемент структуры совпадает с началом объединения, но последующие поля могут иметь смещение. Для переносимости лучше использовать явное приведение через указатели или гарантировать упаковку (\texttt{\#pragma pack}).
    \item \textbf{Отсутствие владения буфером}: структура хранит только указатель на данные, но не управляет памятью. Это гибко, но требует осторожности, чтобы не было висячих указателей.
    \item \textbf{Отсутствие синхронизации}: предполагается, что транзакции обрабатываются в одном контексте или защищены внешними средствами.
    \item \textbf{Неполнота протокола}: определения только заголовка и структуры, но нет спецификации формата данных, кодировки, способа фрагментации и т.д. Это оставлено для вышележащих уровней.
\end{enumerate}

\subsubsection{Макросы конфигурации}

Файл определяет несколько макросов, которые могут быть переопределены пользователем до включения заголовка:

\begin{itemize}
    \item \texttt{ROBO\_TRAN\_COMMAND\_ID}
    \item \texttt{ROBO\_TRAN\_DEV\_ID}
    \item \texttt{ROBO\_TRAN\_HEADER\_IX}
    \item \texttt{ROBO\_PHY\_CHAN\_IX}
\end{itemize}

Если они не определены, используются стандартные типы (\texttt{uint8\_t} и \texttt{uint16\_t}).

\subsection{Общие определения и утилиты (\texttt{robosd\_common.hpp})}

Файл \texttt{robosd\_common.hpp} является центральным заголовочным файлом фреймворка, который включается практически во все остальные модули. Он задаёт базовые типы, макросы, математические константы, полезные шаблонные функции и классы, а также настройки, управляющие поведением различных подсистем. Этот файл служит фундаментом для всего \texttt{robosd} и должен быть включён в любую единицу трансляции, использующую фреймворк.

\subsubsection{Подключаемые заголовки и условная компиляция}

\begin{itemize}
    \item \texttt{robosd\_app\_tuning.hpp} — предполагаемый пользовательский файл с основными макросами настройки приложения.
    \item \texttt{robosd\_target.hpp} — платформозависимый заголовок, определяющий экспорт/импорт символов и другие архитектурные особенности.
    \item Стандартные заголовки: \texttt{<limits>}, \texttt{<algorithm>}, \texttt{<type\_traits>}, \texttt{<stdarg.h>}, \texttt{<stdint.h>}.
\end{itemize}

В файле также переопределяются глобальные операторы \texttt{new} и \texttt{delete} (их реализация находится в \texttt{robosd\_alocator.cpp}), что позволяет интегрировать собственный аллокатор памяти.

\subsubsection{Настраиваемые типы и макросы}

Пользователь может переопределить следующие макросы до включения \texttt{robosd\_common.hpp}:

\begin{description}
    \item[\texttt{ROBO\_TYPE\_RANDOM}] — тип, используемый для генерации случайных чисел. По умолчанию \texttt{int}.
    \item[\texttt{ROBO\_TYPE\_TIME\_US}] — тип для представления микросекунд. По умолчанию \texttt{uint32\_t}.
    \item[\texttt{ROBO\_TYPE\_TIME\_MS}] — тип для миллисекунд. По умолчанию \texttt{uint32\_t}.
    \item[\texttt{ROBO\_UNICODE\_ENABLED}] — если равен 1, включает поддержку широких символов (\texttt{wchar\_t}) и макросы \texttt{RT}, \texttt{RT8}. По умолчанию 0.
    \item[\texttt{ROBO\_APP\_MODULE\_ENABLED}] — включает поддержку модулей (динамическую загрузку). По умолчанию 0.
    \item[\texttt{ROBO\_APP\_MATH\_SHIFT\_ENABLE}] — если 1, использует аппаратные сдвиги для отрицательных чисел (в противном случае эмулирует). По умолчанию 0.
    \item[\texttt{ROBO\_STD\_ARGS}] — позволяет заменить стандартный \texttt{<stdarg.h>} на другой заголовок.
    \item[\texttt{ROBO\_EXPORT\_FUNCTION\_PREFIX}] — префикс для экспортируемых функций (например, пустая строка или \texttt{"\_"}).
    \item[\texttt{ROBO\_APP\_DEBUG\_ENABLED}] — включает отладочные ассерты. По умолчанию 1.
    \item[\texttt{ROBO\_APP\_SYSTEM\_ENABLED}] — включает системный слой. По умолчанию 0.
    \item[\texttt{ROBO\_APP\_RTTI\_ENABLED}] — разрешает использование RTTI. По умолчанию 1.
    \item[\texttt{ROBO\_APP\_REALTIME\_TYPE}] — тип реализации реального времени. По умолчанию \texttt{ROBO\_APP\_TYPE\_SPECIFIC}.
    \item[\texttt{ROBO\_APP\_ULTRACOMPACT}] — флаг ультракомпактного режима. По умолчанию 0.
\end{description}

Также определены константы для выбора типа реализации:
\begin{lstlisting}[caption={Platform type constants}]
#define ROBO_APP_TYPE_NONE 0
#define ROBO_APP_TYPE_NATIVE 1
#define ROBO_APP_TYPE_STD 2
#define ROBO_APP_TYPE_WIN 3
#define ROBO_APP_TYPE_LINUX 4
#define ROBO_APP_TYPE_RASPBERRY 5
#define ROBO_APP_TYPE_UBUNTU 6
#define ROBO_APP_TYPE_ASTRA 7
#define ROBO_APP_TYPE_SPECIFIC 8
#define ROBO_APP_TYPE_DUMMY 9
\end{lstlisting}

\subsubsection{Макросы для работы с Unicode}

В зависимости от \texttt{ROBO\_UNICODE\_ENABLED} определяются:

\begin{itemize}
    \item \texttt{ROBO\_CHAR} — тип символа: \texttt{wchar\_t} или \texttt{char}.
    \item \texttt{ROBO\_CONST\_STRING} — тип строкового литерала: \texttt{wchar\_t const*} или \texttt{char const*}.
    \item \texttt{RT(s)} — преобразует обычную строку в широкую, если нужно: \texttt{L##s} или просто \texttt{s}.
    \item \texttt{RT8(s)} — всегда остаётся узкой строкой (для совместимости).
\end{itemize}

\subsubsection{Псевдонимы типов в пространстве имён \texttt{robo}}

\begin{lstlisting}[caption={Basic type aliases}]
typedef ROBO_TYPE_RANDOM random_t;
typedef ROBO_TYPE_TIME_US time_us_t;
typedef ROBO_TYPE_TIME_MS time_ms_t;
typedef ROBO_CHAR char_t;
typedef ROBO_CONST_STRING cstr;
\end{lstlisting}

\subsubsection{Функции хеширования и аварийного останова}

\begin{itemize}
    \item \texttt{int32\_t hash(cstr src, int32\_t begin = 0)} — простая хеш-функция, обновляющая начальное значение. Используется для генерации ключей в интрузивных списках.
    \item \texttt{int32\_t hash(cstr beg, cstr end, int32\_t begin = 0)} — хеширование диапазона.
    \item \texttt{void crash(char const* function, char const* file, int line)} — аварийный останов с сохранением места ошибки. Вызывается макросом \texttt{ROBO\_APP\_CRASH()}.
\end{itemize}

\subsubsection{Перечисление \texttt{result}}

\begin{lstlisting}
enum class result { complete, resume, panic };
\end{lstlisting}
Используется для возврата состояния из функций жизненного цикла узлов и модулей.

\subsubsection{Математические константы и функции}

В пространстве имён \texttt{robo} определены шаблонные константы:
\begin{lstlisting}
template<typename T> constexpr T pi = T(3.1415926535897932385);
template<typename T> constexpr T deg2rad = pi<T> / T(180.0);
template<typename T> constexpr T rad2deg = T(180.0) / pi<T>;
template<typename T> constexpr T g = T(9.8066);
\end{lstlisting}

Функция \texttt{csqrt} — рекурсивное вычисление квадратного корня на этапе компиляции (constexpr), основанное на методе Ньютона.
\begin{lstlisting}
double constexpr csqrt_helper(double x, double curr, double prev);
template<typename T> T constexpr csqrt(T x);
\end{lstlisting}

Функция \texttt{parity<T>(T n)} — вычисляет чётность количества единичных бит (паритет).

\subsubsection{Классы для сериализации: \texttt{ostram} и \texttt{istram}}

Простые потоки для упаковки/распаковки данных в/из байтового буфера. Работают на этапе компиляции (constexpr).

\begin{lstlisting}[caption={ostram usage example}]
uint8_t buffer[256];
robo::ostram out(buffer, sizeof(buffer));
out.put(42);                 // put int
out << 3.14f;                 // put float
size_t written = out.actual_size();

robo::istram in(buffer, written);
int x; float y;
in.get(x) >> y;               // read back
\end{lstlisting}

\subsubsection{Простые контейнеры}

\begin{itemize}
    \item \texttt{array<T, C>} — обёртка над статическим массивом фиксированного размера.
    \item \texttt{stack\_t<T, S, C>} — наследник \texttt{array}, добавляющий счётчик \texttt{len}. Предназначен для реализации стека.
\end{itemize}

\subsubsection{Утилиты для работы с числами}

\begin{itemize}
    \item \texttt{constexpr T abs(T x)} — абсолютное значение.
    \item \texttt{constexpr T min(T x, T y)}, \texttt{max(T x, T y)} — минимум и максимум.
    \item \texttt{constexpr T fma(T a, T b, T c)} — fused multiply-add: $a*b + c$.
    \item \texttt{constexpr T saturate(T x, H lo, H hi)} — ограничение значения диапазоном.
\end{itemize}

\subsubsection{Пространство имён \texttt{digit}: округление и сдвиги}

Содержит функции для безопасного сдвига и округления целых чисел с учётом знака. Важно для реализации фиксированной точки и цифровой обработки.

\begin{lstlisting}[caption={digit namespace functions}]
namespace digit {
    template<typename T> constexpr T rsh(T x, int8_t sh);  // rounded shift right
    template<typename T> constexpr T lsh(T x, int8_t sh);  // shift left
    template<typename T> constexpr T round(T src, uint8_t shift); // round to nearest
    template<typename D, typename S> constexpr D pack(const S& x, uint8_t shift); // pack with saturation
    template<typename T> T select(const T& mask, const T& src0, const T& src1); // bitwise select
}
\end{lstlisting}

Эти функции учитывают макрос \texttt{ROBO\_APP\_MATH\_SHIFT\_ENABLE} для выбора аппаратной или эмулированной реализации.

\subsubsection{Генерация случайных чисел}

Шаблон \texttt{rand\_t<T, R>} использует статический метод \texttt{R::rand()} для получения случайного числа и масштабирует его в заданный диапазон. Используется в \texttt{system::rand}.

\subsubsection{Аппроксимация арктангенса: \texttt{catan2}}

Функция \texttt{catan2} реализует быстрый расчёт арктангенса с использованием полиномиальной аппроксимации (таблица коэффициентов). Подходит для встраиваемых систем, где нет аппаратного \texttt{atan2}.

\subsubsection{Таблицы для \texttt{atan2}}

Шаблон \texttt{atan2\_table\_t<T, N>} генерирует на этапе компиляции таблицу значений \texttt{atan2(y,1)} для равномерно распределённых \texttt{y} от 0 до $\pi/4$. Может использоваться для быстрой интерполяции.

\subsubsection{Синус через ряд Тейлора: \texttt{csin} и \texttt{taylor\_sin}}

\texttt{taylor\_sin} — аппроксимация синуса рядом Тейлора (10 членов). \texttt{csin} — обобщённая версия для произвольного типа.

\subsubsection{Экспонента матрицы: \texttt{expm}}

Вычисляет матричную экспоненту через ряд, пока норма очередного члена не станет меньше заданного порога. Требует, чтобы тип матрицы поддерживал операции сложения, умножения, деления на скаляр и метод \texttt{norma()}.

\subsubsection{Генератор сигналов: \texttt{generator\_t}}

Позволяет создать таблицу значений синусоиды (или другой функции) на этапе компиляции, используя параметры из класса-конфигурации \texttt{C}, который должен предоставлять типы \texttt{type}, константы \texttt{zero}, \texttt{mag}, \texttt{count}, \texttt{offset}. Применяется для генерации ШИМ-таблиц, волновых таблиц и т.п.

\begin{lstlisting}[caption={Generator usage example}]
struct SineConfig {
    using type = int16_t;
    static constexpr type zero = 0;
    static constexpr type mag = 32767;
    static constexpr int count = 256;
    static constexpr int offset = 0;
};
robo::generator_t<SineConfig> sine_table;
// sine_table.table[i] contains precomputed sine values
\end{lstlisting}

\subsubsection{Антидребезг: \texttt{debouncer\_t}}

Шаблонный класс для подавления дребезга контактов (кнопки, переключатели). Работает с массивом каналов, использует временные метки от класса \texttt{S}, который должен предоставлять статический метод \texttt{time\_us()}. Поддерживает индивидуальные пороги для каждого канала.

\begin{lstlisting}[caption={Debouncer usage example}]
struct Timer {
    static robo::time_us_t time_us() { return robo::system::time_us(); }
};

robo::debouncer_t<4, uint8_t, Timer> debouncer;
bool inputs[4] = { ... };
robo::time_us_t thresholds[4] = { 5000, 5000, 5000, 5000 };
debouncer.update(thresholds, inputs);
uint8_t stable = debouncer.stable(); // stable states after debouncing
\end{lstlisting}

\subsubsection{Обёртки математических функций: \texttt{math\_t}}

Шаблон \texttt{math\_t<T>} предоставляет статические методы \texttt{cos}, \texttt{sin}, \texttt{sqrt}, которые вызывают соответствующие функции из стандартной библиотеки (с приведением к \texttt{double}). Специализация для \texttt{float} использует \texttt{cosf}, \texttt{sinf}, \texttt{sqrtf}. Это позволяет писать код, не зависящий от типа с плавающей точкой.

\subsubsection{Макросы для отладки и ассертов}

\begin{lstlisting}[caption={Debug macros}]
#define ROBO_APP_CRASH() robo::crash(ROBO_APP_PROC_NAME, ROBO_APP_PROC_FILE, ROBO_APP_PROC_LINE)
#define ROBO_APP_ASSERT(x) if(!(x)) ROBO_APP_CRASH()
\end{lstlisting}

Макросы \texttt{ROBO\_APP\_PROC\_NAME}, \texttt{\_FILE}, \texttt{\_LINE} определяются в \texttt{robosd\_target.hpp} (обычно как \texttt{\_\_func\_\_}, \texttt{\_\_FILE\_\_}, \texttt{\_\_LINE\_\_}).

\subsubsection{Макросы для работы с variadic arguments}

В файле определяются макросы \texttt{PP\_THIRD\_ARG}, \texttt{VA\_OPT\_SUPPORTED\_I}, \texttt{VA\_OPT\_SUPPORTED}, которые используются для проверки поддержки \texttt{\_\_VA\_OPT\_\_} компилятором. Это нужно для корректной обработки пустых списков аргументов в макросах логирования.

\subsubsection{Потенциальные проблемы и устаревшие места}

\begin{enumerate}
    \item \textbf{Неоднозначность хеш-функции}: используется простая аддитивная хеш-функция, которая может приводить к коллизиям. Для больших наборов имён стоит рассмотреть более качественные алгоритмы.
    \item \textbf{Использование \texttt{double} в constexpr функциях}: в C++11/14 constexpr с плавающей точкой ограничено, но код использует \texttt{double} в \texttt{csqrt\_helper}, \texttt{catan2}, \texttt{taylor\_sin} и т.д. Это может не компилироваться в старых стандартах. Фреймворк, вероятно, требует C++17 или выше.
    \item \textbf{Пространство имён \texttt{digit}::select} — использует побитовые операции, предполагая, что маска корректно настроена. Для нецелочисленных типов поведение не определено.
    \item \textbf{Класс \texttt{generator\_t}} требует, чтобы \texttt{C::type} был арифметическим, и выполняет приведение \texttt{double} к этому типу с округлением. При малых разрядностях возможно переполнение.
    \item \textbf{Класс \texttt{debouncer\_t}} использует \texttt{time\_us\_t} для хранения времени начала; при длительной работе счётчик может переполниться (если \texttt{time\_us\_t} 32-битный, переполнение через 4295 секунд). Для длительных интервалов дребезга это может быть проблемой.
    \item \textbf{Макросы \texttt{ROBO\_APP\_ASSERT} и \texttt{ROBO\_APP\_CRASH}} не используют информацию о маске логирования; в релизных сборках они могут быть пустыми, что скрывает ошибки.
\end{enumerate}

\subsubsection{Пример полного использования}

\begin{lstlisting}[caption={Mixing common utilities}]
#include "core/robosd_common.hpp"

using namespace robo;

// Compute parity of a number
static_assert(parity(0b1010) == true);  // even number of ones? actually parity(0b1010) -> 1 (odd)

// Generate sine table
struct SineConf {
    using type = int16_t;
    static constexpr type zero = 0;
    static constexpr type mag = 1000;
    static constexpr int count = 64;
    static constexpr int offset = 0;
};
generator_t<SineConf> sine;
// sine.table[0] approx 0, sine.table[16] approx 1000, etc.

// Serialization
uint8_t buffer[64];
ostram os(buffer, 64);
os << 123 << 45.67f;

istram is(buffer, os.actual_size());
int i; float f;
is >> i >> f;
// i == 123, f == 45.67f

// Math helpers
double angle = deg2rad<double> * 90.0;   // 1.5708 rad
double s = taylor_sin(angle);            // approx 1.0
\end{lstlisting}

\section{Библиотека \texttt{burst} — управление реальным временем и устройствами}

Библиотека \texttt{burst} представляет собой высокоуровневый каркас для построения систем управления в реальном времени. Она базируется на ядре \texttt{robosd} и добавляет концепции: плата (\texttt{board}) как корневой диспетчер, устройства (\texttt{dev}) с набором режимов работы (\texttt{mode}), слоты выполнения (\texttt{slot}) для организации циклических вызовов, очередь асинхронных запросов (\texttt{request}) для взаимодействия между frontend и backend, а также систему регистрации переменных (\texttt{vartree}) для конфигурации и мониторинга в реальном времени. Математическая подсистема предоставляет эффективные вычисления с фиксированной точкой и преобразования координат.

\subsection{Состав библиотеки}

\begin{itemize}
    \item \texttt{burst\_common.hpp} — общие макросы и структуры (диапазоны, гистерезис, флаги включения).
    \item \texttt{dev.front.hpp} — определения для внешнего интерфейса (статусы, паники, команды).
    \item \texttt{vartree.hpp} / \texttt{vartree.cpp} — система древовидных переменных.
    \item \texttt{math.hpp} / \texttt{math.cpp} — математические классы с фиксированной точкой, тригонометрия, преобразования.
    \item \texttt{burst.hpp} — главный заголовок с классами \texttt{request}, \texttt{dev}, \texttt{board}, \texttt{timer\_t} и макросами.
    \item \texttt{burst.cpp} — реализация ключевых методов, очереди запросов, обработка паник, циклы.
    \item \texttt{burst.h} — C-обёртки для вызова из чистого C.
\end{itemize}

\subsection{Общие определения (\texttt{burst\_common.hpp})}

Файл содержит настраиваемые макросы и простые шаблонные структуры.

\begin{itemize}
    \item \texttt{range\_s<A>} — структура с полями \texttt{lo} и \texttt{hi}.
    \item \texttt{hyst\_t<R>} — структура для гистерезиса: \texttt{overhi}, \texttt{hi}, \texttt{lo}, \texttt{ultralo}.
    \item Макросы \texttt{BURST\_RANGE\_CONFIG(a)} и \texttt{BURST\_HIST\_CONFIG(a)} для удобного заполнения конфигурационных структур.
    \item Множество макросов для включения отдельных функций защиты (температура, напряжение, ток, потеря мастера). По умолчанию многие выключены.
    \item \texttt{BURST\_PROTECTION\_ENABLED} — общий флаг защиты.
    \item \texttt{BURST\_PANICS\_MASTER\_LOST\_ENABLED} — контроль потери связи с мастером.
    \item \texttt{ROBO\_APP\_BURST\_REALTIME\_SLOT\_ENABLE} — включает отдельный слот реального времени.
    \item \texttt{ROBO\_APP\_BURST\_VARTREE\_ENABLED} — включает систему переменных.
\end{itemize}

Макросы используют префиксы \texttt{ROBO\_APP\_} (глобальные) и \texttt{BURST\_} (локальные). Многие из них должны быть определены пользователем до компиляции.

\subsection{Интерфейс frontend (\texttt{dev.front.hpp})}

Содержит перечисления и структуры, определяющие взаимодействие между \texttt{burst} и внешним миром (например, по протоколу).

\begin{itemize}
    \item \texttt{burst::front::board::commands} — \texttt{none}, \texttt{configure}.
    \item \texttt{burst::front::board::statuses} — \texttt{unknown}, \texttt{startuped}, \texttt{configure}, \texttt{restart}, \texttt{normal}.
    \item \texttt{burst::front::board::panics::bits} и \texttt{masks} — биты паник платы (перенапряжение, недонапряжение, перегрузка по току, температура и т.д.).
    \item \texttt{burst::front::tp\_verb} — уровни трассировки (realtime, backend, frontend, loop).
    \item \texttt{burst::front::dev::panics::bits} — паники устройства: \texttt{master\_lost}, \texttt{board}.
    \item \texttt{burst::front::dev::modes} — пока только \texttt{idle}.
    \item \texttt{burst::front::dev::action\_s} — структура, содержащая поле \texttt{mode} и флаг \texttt{action\_actual}. При условной компиляции может включать поля из \texttt{robo::common::devagent}.
    \item \texttt{burst::front::dev::feedback\_s} — аналогично, содержит \texttt{mode} и \texttt{fault}.
\end{itemize}

Этот файл связывает внутренние механизмы \texttt{burst} с протоколами верхнего уровня (например, servo). Использование \texttt{\#ifdef ROBO\_APP\_BURST\_SERVO\_SIDE} и т.п. указывает на две возможные роли: сторона сервопривода или сторона управляющего устройства.

\subsection{Система древовидных переменных (\texttt{vartree.hpp}, \texttt{vartree.cpp})}

Предназначена для регистрации и доступа к переменным (конфигурационным, статусным) в иерархическом пространстве имён. Ключевые элементы:

\begin{itemize}
    \item \texttt{tags} — \texttt{push}, \texttt{var}, \texttt{pop} для построения дерева.
    \item \texttt{ref\_s} — двухбитовый тег (общий предок всех записей).
    \item \texttt{mode} — перечисление уровней доступа: \texttt{none}, \texttt{tuning}, \texttt{action}, \texttt{config}, \texttt{full}.
    \item \texttt{descriptor} — упакованный дескриптор типа переменной (16 бит). Содержит длину, знак, константность, признак вещественного числа, флаг необходимости реконфигурации. Битовые поля используют \texttt{\#pragma pack(1)}.
    \item \texttt{descriptor\_enco} — constexpr функция для формирования дескриптора.
    \item \texttt{types} — перечисление предопределённых типов (uint8, int8, uint16, int16, uint32, int32, uint64, int64, real, ext, и их константные варианты). Значения вычисляются через \texttt{descriptor\_enco}.
    \item Шаблонные функции \texttt{const\_unsigned\_type}, \texttt{const\_signed\_type}, \texttt{unsigned\_type}, \texttt{signed\_type} — определяют тип по размеру.
    \item \texttt{type\_name} — возвращает строковое имя типа.
    \item \texttt{request} — константы для протокола: \texttt{index}, \texttt{get}, \texttt{put}.
    \item \texttt{error} — коды ошибок.
    \item \texttt{header\_s} — заголовок запроса (2 байта: 2 бита query, 4 бита error, 10 бит ix).
    \item \texttt{record\_s} — запись переменной: дескриптор, имя, ключ (хеш), адрес.
    \item \texttt{reg(ref\_s*)} — добавляет ссылку в глобальный массив \texttt{index}.
    \item \texttt{push(name)} и \texttt{pop()} — создают узлы для построения пути.
    \item \texttt{reg<T>(type, addr, name, reconfig\_need)} — создаёт запись переменной, вычисляет дескриптор, вызывает \texttt{reg}.
    \item \texttt{varreg} — перегрузки для регистрации диапазона или гистерезиса (создаёт несколько переменных в одной группе).
    \item \texttt{reindex} — пересчитывает ключи (хеши) на основе текущего пути. Используется после регистрации всех переменных.
    \item \texttt{if\_configure}, \texttt{reconfig\_query} — для ультракомпактного режима.
    \item \texttt{showpro} и \texttt{show} — вспомогательные классы для печати дерева переменных (например, для FreeMASTER).
\end{itemize}

В \texttt{vartree.cpp} реализованы:

\begin{itemize}
    \item Глобальный массив \texttt{index} (размер \texttt{pool\_size}), счётчик \texttt{count}.
    \item \texttt{reg} — простая вставка.
    \item \texttt{push}/\texttt{pop} — создают объекты \texttt{node\_s} (тег push) или просто \texttt{ref\_s} (pop) и регистрируют.
    \item \texttt{reindex} — обходит все записи, строит путь (имитируя стек), вычисляет хеш для каждой переменной. Использует временные буферы.
    \item \texttt{find} — линейный поиск по ключу.
    \item \texttt{get} — доступ по индексу.
    \item \texttt{proto} — обработчик запросов по протоколу (index, get, put). Реализует чтение/запись переменных с проверкой прав.
    \item \texttt{sprintf} — форматирует значение переменной в строку.
    \item \texttt{first}/\texttt{last} — для итерации.
\end{itemize}

\textbf{Недоработки:}
\begin{itemize}
    \item Линейный поиск — при большом количестве переменных может быть медленным.
    \item Нет защиты от повторной регистрации одинаковых имён.
    \item Использование глобального массива фиксированного размера (\texttt{pool\_size}) может привести к переполнению.
    \item В \texttt{proto} при обработке \texttt{put} для переменных с \texttt{reconfig\_need} проверяется режим конфигурации через глобальную функцию \texttt{if\_configure()}, что создаёт сильную связность с \texttt{board}.
    \item Массив \texttt{index} хранит сырые указатели, нет автоматического управления памятью — требуется явный вызов \texttt{free}.
\end{itemize}

\subsection{Математическая библиотека (\texttt{math.hpp}, \texttt{math.cpp})}

Предоставляет два основных числовых пространства: \texttt{int15} (фиксированная точка 15 бит) и \texttt{real} (float). Вокруг них построены шаблонные классы \texttt{fixed\_point<digit>} и \texttt{float\_point<q>} с множеством полезных функций.

\subsubsection{Класс \texttt{int15}}

\begin{itemize}
    \item Типы: \texttt{discret\_t} = int16\_t, \texttt{signal\_t} = int16\_t, \texttt{long\_signal\_t} = int32\_t, \texttt{usignal\_t} = uint16\_t и т.д.
    \item Константы: \texttt{max = 32767}, \texttt{min = -max}, \texttt{ones = max}.
    \item Статические функции: \texttt{s\_round}, \texttt{s\_frac}, \texttt{l\_round}, \texttt{l\_frac} для преобразования double в фиксированную точку.
    \item \texttt{sin}, \texttt{cos}, \texttt{sqrt}, \texttt{atan2} — реализации с использованием таблиц.
    \item В \texttt{math.cpp} приведена таблица синуса \texttt{mcgenSineTable256} (256 значений) и функции \texttt{mcsinPIxLUT}, \texttt{sin}, \texttt{cos}.
    \item \texttt{sqrt} реализован через аппроксимацию с таблицами \texttt{dsqrtGainArray}, \texttt{dsqrtOffsetArray}, \texttt{dsqrtShift}.
    \item \texttt{atan2} использует таблицу \texttt{tableAtan64} и вспомогательную функцию \texttt{\_atan}.
\end{itemize}

\subsubsection{Класс \texttt{real}}

\begin{itemize}
    \item Типы: \texttt{signal\_t = float}, \texttt{long\_signal\_t = float} и т.д.
    \item Константы: \texttt{pi}, \texttt{discret\_max = 32767}.
    \item Аналогичные функции \texttt{sin}, \texttt{cos}, \texttt{atan2}, но с плавающей точкой.
    \item Таблица синуса \texttt{real\_mcgenSineTable256} (аналогична int15, но значения float).
    \item \texttt{atan2} использует \texttt{robo::atan2\_table\_t} (таблица, генерируемая на этапе компиляции).
\end{itemize}

\subsubsection{Шаблон \texttt{fixed\_point<digit>} и \texttt{float\_point<q>}}

Эти классы объединяют функции и константы, добавляя:

\begin{itemize}
    \item \texttt{sqrt3\_div\_2}, \texttt{sqrt2\_div\_2}, \texttt{one\_div\_sqrt3} — предвычисленные константы.
    \item \texttt{s\_mult}, \texttt{s\_add}, \texttt{s\_dot} — операции с аккумуляцией в расширенном типе.
    \item \texttt{long\_signal\_u} — union для разложения 32-битного числа на два 16-битных (используется в умножении).
    \item \texttt{qa} — структура для квадратичной аппроксимации.
    \item \texttt{lsf} — структура с линейными операциями.
    \item \texttt{pi} — структура для ПИ-регулятора (дифференциальная и квадратичная составляющие).
    \item \texttt{signal2discret\_s} и \texttt{discret2signal} — конвертеры между физическим сигналом и целым кодом.
    \item \texttt{ab\_t}, \texttt{abc\_t}, \texttt{dq\_t} — для преобразований координат (Кларк, Парк).
\end{itemize}

В \texttt{math.cpp} также есть тестовая функция \texttt{sin\_test}.

\textbf{Недоработки:}
\begin{itemize}
    \item Множество закомментированных блоков и устаревших определений (например, большая часть в конце \texttt{math.hpp} закомментирована).
    \item Функция \texttt{atan2} в \texttt{int15} использует 64-битные вычисления, что может быть медленно на 32-битных процессорах.
    \item Таблицы для \texttt{sqrt} и \texttt{atan2} не сопровождаются комментариями о происхождении и точности.
    \item Некоторые макросы (например, \texttt{S(x)}) используются без явного контекста, что затрудняет чтение.
    \item В \texttt{float\_point::lsf} функция \texttt{sat} просто возвращает аргумент (нет ограничений) — вероятно, недоделано.
    \item Отсутствует документация по использованию шаблонных преобразователей.
\end{itemize}

\subsection{Основные классы \texttt{burst} (\texttt{burst.hpp}, \texttt{burst.cpp})}

\subsubsection{Класс \texttt{request}}

Предназначен для асинхронного выполнения действий в другом контексте (frontend/backend). Содержит указатели на делегаты \texttt{query} и \texttt{confirm}. При вызове \texttt{post()} объект помещается в одну из двух глобальных очередей (\texttt{frontend\_queue} или \texttt{backend\_queue}) в зависимости от текущего контекста. Статические методы \texttt{front\_perfom} и \texttt{backend\_perfom} извлекают запросы из соответствующей очереди и выполняют их, переключая статус между \texttt{query} и \texttt{confirm}. Это классический механизм для безопасного переключения контекста.

\subsubsection{Класс \texttt{dev} — устройство}

Базовый класс для всех устройств. Содержит:

\begin{itemize}
    \item Ссылки на конфигурацию (\texttt{config\_s}), состояние (\texttt{present\_s}), действие (\texttt{action\_s}) и обратную связь (\texttt{feedback\_s}). Эти структуры должны быть определены пользователем.
    \item Карту режимов (\texttt{modes\_}) — список уникальных режимов (наследников \texttt{mode}).
    \item Текущий режим (\texttt{actual\_mode\_}) и методы для переключения (\texttt{switch\_to\_mode}).
    \item Методы \texttt{raise\_panic}, \texttt{reset\_panic}, \texttt{reset\_panics} для управления битами паники.
    \item Циклические функции \texttt{loopA}, \texttt{loopB}, \texttt{loopC}, \texttt{frontend\_loop}, которые делегируют вызов текущему режиму.
    \item Поддержка дерева переменных: \texttt{regvar\_present}, \texttt{regvar\_action}, \texttt{regvar\_conf} (определены в \texttt{burst.cpp}).
\end{itemize}

Класс \texttt{mode} — вложенный в \texttt{dev}, должен быть переопределён пользователем. Содержит чисто виртуальные методы \texttt{applay\_action}, \texttt{start}, \texttt{stop}, \texttt{loopA/B/C}, \texttt{frontend\_loop}. Регистрируется в устройстве по уникальному идентификатору.

\subsubsection{Класс \texttt{board} — плата}

Синглтон, управляющий всем. Содержит:

\begin{itemize}
    \item Список устройств (\texttt{dev::map}).
    \item Слоты выполнения (\texttt{slots}) — набор объектов \texttt{slot}, каждый из которых содержит список делегатов (наследников \texttt{slot::delegat}), вызываемых в определённый момент цикла.
    \item Предопределённые слоты: \texttt{begin}, \texttt{start}, \texttt{startup}, \texttt{realtime}, \texttt{backend}, \texttt{frontend}, \texttt{halt}, \texttt{ready}, \texttt{reconfig} и массив \texttt{periodic} (число задаётся макросом \texttt{ROBO\_APP\_BURST\_SLOT\_COUNT}).
    \item Состояние платы (\texttt{present\_s}) — паники, время последней паники, команда, статус, уровень трассировки.
    \item Конфигурация платы (\texttt{config\_s}) — версия, флаг ожидания конфигурации при старте, параметры защиты.
    \item Методы жизненного цикла: \texttt{begin}, \texttt{realtime\_loop}, \texttt{backend\_loop}, \texttt{frontend\_loop}, \texttt{handle\_panic}, \texttt{raise\_panic}, \texttt{reset\_panic}.
    \item Статические методы для доступа извне: \texttt{begin}, \texttt{backend\_loop}, \texttt{frontend\_loop} и т.д.
\end{itemize}

Класс \texttt{slot::delegat} — базовый для всех функций, которые можно привязать к слоту. Имеет методы \texttt{attach} и \texttt{dettach} для добавления в конкретный слот. Предопределены удобные наследники: \texttt{simple} (обычная функция), \texttt{lambda} (объект \texttt{robo::lambda}), \texttt{member} (метод класса).

Класс \texttt{timer\_t< K, D>} — шаблон для создания периодических таймеров. Параметр \texttt{K} — тип ключа (например, \texttt{slot::kind}), \texttt{D} — тип делегата (например, \texttt{slot::simple}). Таймер хранит период, флаг однократности, состояние и в операторе \texttt{()} проверяет, не пора ли вызвать делегат.

В \texttt{burst.cpp} реализованы:

\begin{itemize}
    \item Глобальные очереди \texttt{backend\_queue} и \texttt{frontend\_queue} (тип \texttt{robo::safe\_ring\_t}).
    \item Методы \texttt{request::post}, \texttt{front\_perfom}, \texttt{backend\_perfom}.
    \item Конструкторы \texttt{dev} и \texttt{mode}.
    \item \texttt{dev::raise\_panic}, \texttt{reset\_panic}, \texttt{perform\_panic}, \texttt{switch\_to\_mode}, \texttt{loopA/B/C}, \texttt{frontend\_loop}.
    \item \texttt{board::begin\_} — инициализация, запуск слотов \texttt{begin} и \texttt{start}, инициализация сети и переменных, запуск системного таймера.
    \item \texttt{board::realtime\_loop\_}, \texttt{backend\_loop\_}, \texttt{frontend\_loop\_} — реализация циклов с вызовом соответствующих слотов.
    \item Обработка паник (\texttt{raise\_panic\_}, \texttt{reset\_panic\_}, \texttt{handle\_panic\_}).
    \item Защита (\texttt{realtime\_protection\_}, \texttt{frontend\_protection\_}) — проверка напряжения, тока, температуры, состояния дружественных плат.
    \item Регистрация переменных платы (\texttt{regvar\_conf}).
\end{itemize}

\subsection{C-обёртки (\texttt{burst.h})}

Предоставляет функции \texttt{burst\_begin}, \texttt{burst\_begin\_ps}, \texttt{burst\_realtime\_loop}, \texttt{burst\_backend\_loop}, \texttt{burst\_frontend\_loop}, \texttt{burst\_core\_crash\_} для вызова из чистого C. Макрос \texttt{burst\_core\_crash} упрощает вызов.

\subsection{Ключевые идеи и достоинства}

\begin{itemize}
    \item \textbf{Разделение контекстов} через очереди запросов позволяет безопасно взаимодействовать frontend и backend.
    \item \textbf{Слоты выполнения} дают гибкость в организации периодических задач: можно назначать функции на разные фазы цикла (начало, запуск, реальное время, backend, frontend и т.д.) с возможностью задания нескольких периодических слотов.
    \item \textbf{Устройства с режимами} упрощают реализацию конечных автоматов для управления оборудованием.
    \item \textbf{Дерево переменных} предоставляет единый способ конфигурации и мониторинга, легко интегрируется с протоколами (FreeMASTER, CANopen).
    \item \textbf{Математика с фиксированной точкой} обеспечивает детерминизм и высокую производительность на встраиваемых системах без FPU.
    \item \textbf{Широкие возможности настройки} через макросы позволяют адаптировать библиотеку под конкретные требования (включение/отключение защиты, выбор количества слотов и т.д.).
\end{itemize}

\subsection{Недоработки и несуразности}

\begin{enumerate}
    \item \textbf{Сложность и избыточность}. Многие классы тесно переплетены, что затрудняет понимание и модификацию. Например, \texttt{board} знает о \texttt{vartree}, а \texttt{vartree} вызывает \texttt{board::if\_configure()}.
    \item \textbf{Отсутствие документации по использованию}. Нет примеров создания собственных устройств и режимов.
    \item \textbf{Неполнота реализации защиты}. В \texttt{board::frontend\_protection\_} используются функции \texttt{temper\_get\_lo\_pp}, \texttt{voltage\_get\_pp} и т.п., которые не определены в библиотеке — предполагается, что пользователь предоставит их реализацию. Это не очевидно.
    \item \textbf{Проблемы с потоками}. Очереди запросов (\texttt{safe\_ring\_t}) используют \texttt{system::guard}, что корректно, но в остальном коде синхронизация отсутствует — предполагается, что вызовы циклов происходят из одного потока.
    \item \textbf{Макросы с неочевидными именами}. Например, \texttt{DEV\_PRESENT\_S(p)} раскрывается в объявление ссылки на структуру \texttt{present\_s}. Это затрудняет отладку.
    \item \textbf{Потенциальные ошибки в алгоритмах}. В \texttt{dev::switch\_to\_mode} используется \texttt{mode\_id == front::dev::modes::idle} как особый случай, но в \texttt{mode} нет виртуального деструктора — возможно, проблема с удалением.
    \item \textbf{Незавершённые части}. В \texttt{math.hpp} много закомментированного кода, некоторые функции не реализованы (например, \texttt{transform::forward}).
    \item \textbf{Жёсткая связь с \texttt{robosd\_system}}. Библиотека полагается на системные функции (таймер, очереди, guard), что усложняет перенос на другие платформы.
    \item \textbf{Управление памятью}. В \texttt{vartree} используются \texttt{new} без соответствующих \texttt{delete} (кроме \texttt{free}), но \texttt{free} не вызывается автоматически. Пользователь должен явно вызывать \texttt{burst::var::free()} при завершении.
    \item \textbf{Ошибки в коде}. Например, в \texttt{board::begin\_} после \texttt{system::start} вызывается \texttt{raise\_panic\_(front::board::panics::bits::config)} — это поднимает панику конфигурации, хотя плата только запускается. Возможно, это intentional, но выглядит странно.
\end{enumerate}

\subsection{Заключение}

Библиотека \texttt{burst} является мощным инструментом для построения встраиваемых систем управления, но требует от разработчика глубокого понимания её внутреннего устройства. Многие механизмы реализованы качественно, однако код содержит следы эволюции и нуждается в рефакторинге и документировании. При использовании следует быть готовым к доработке недостающих частей и тщательному тестированию.

\end{document}